\documentclass[10pt,a4paper]{article}
\AtBeginDocument{\fontsize{14pt}{17pt}\selectfont}
\usepackage[a4paper,margin=1.5cm]{geometry}  % comfortable page margins
\usepackage{microtype}                       % improved character spacing/justification
\usepackage{enumitem}                        % fine-grained list control

%% fix mathbb font for digits
\DeclareMathAlphabet{\mathbb}{U}{bbold}{m}{n}
\usepackage{amsmath,amssymb,amsthm}          % math environments, symbols, theorem-like environments
\usepackage{booktabs}                        % rules in tables
\usepackage{tikz}                            % schematic figures
\usetikzlibrary{arrows.meta,decorations.markings,calc,patterns,positioning,shapes.geometric}
\usepackage{caption}                         % caption formatting
\usepackage[T1]{fontenc}
\usepackage{lmodern}
\usepackage[scr=boondoxo]{mathalfa}       % lowercase script letters: the symbol \lds of the logarithmic density of states
\usepackage[colorlinks=true,linkcolor=blue,citecolor=blue]{hyperref}

\theoremstyle{plain}
\newtheorem{proposition}{Proposition}
\newtheorem{lemma}{Lemma}
%% assumptions on the spectrum are numbered A1, A2, ...
\newtheorem{assumption}{Assumption}
\renewcommand{\theassumption}{A\arabic{assumption}}

\newcommand{\Var}{\operatorname{Var}}
\newcommand{\Cov}{\operatorname{Cov}}
\newcommand{\Tr}{\operatorname{Tr}}
\newcommand{\cN}{\mathcal{N}}   % level-counting function
\newcommand{\cA}{{\mathcal{A}_0}}   % amplitude of the logarithmic density of states
\newcommand{\evalat}{\bigg\vert}  % evaluation bar, as in (d/d beta)|_mu
\newcommand{\lndos}{{logdos}}  % abbreviation for the logarithmic density of states
\newcommand{\hfield}{\mathscr{h}}  % parameter (field or coupling) of the Hamiltonian
\newcommand{\enrg}{\mathcal{E}}  % energy eigenvalue of H - mu N
\newcommand{\fen}[1][T]{\mathscr{F}_{#1}}  % F_T(eps) = eps - T s(eps): free energy per unit size at a prescribed energy
\newcommand{\lds}{\mathscr{s}}  % the logarithmic density of states (microcanonical entropy density) as a function of the energy
\newcommand{\Tmax}{T_{\max}}  % limiting temperature of the linear bound on the logdos
\newcommand{\etaD}{\eta_{\rm D}}  % Dirichlet eta function (eta is the anomalous exponent)


\title{Phase transitions and power laws\\ from the many-body density of states\\
\large A spectral formulation of thermodynamics beyond the order parameter}
\author{A. Nejati}
\date{}

\begin{document}
\maketitle

\begin{abstract}
The grand partition function of a many-body Hamiltonian is the Laplace
transform of its density of states, and for a system of large size the
thermodynamic potential is the Legendre transform of the logarithmic density
of states $\lds(\varepsilon)$, the microcanonical entropy per unit size. This
note uses these two facts to formulate phase transitions, crossovers and
power laws directly in terms of the microscopic spectrum, instead of the
free energy. Every thermodynamic
quantity at temperature $T$ is expressed through $\lds$ and its derivatives at
the single energy $\varepsilon_*(T)$ at which the slope of $\lds$ is $1/T$, and
an assumption on the form of $\lds$ is converted into a statement on the
observables. A power law $\lds\propto(\varepsilon-\varepsilon_0)^{\sigma}$
above the ground-state energy $\varepsilon_0$ is equivalent to a heat capacity
$\propto T^{\sigma/(1-\sigma)}$. Homogeneity of $\lds$ in energy and field gives
the scaling functions of a quantum critical point and, about an energy in the
interior of the spectrum, those of a thermal critical point, whose static
exponents follow from two exponents of the spectrum. Homogeneity in the wave vector of a modulated field and in the size
gives the correlation length, the Fisher relation and finite-size scaling;
hyperscaling is the invariance of the total number of levels under the
rescaling. Discontinuous, continuous and quantum transitions and crossovers
are classified by the motion of $\varepsilon_*$ through the spectrum. Since
the input is a property of the spectrum and not a symmetry, Goldstone modes,
degenerate manifolds of ordering wave vectors, competing orders, residual
entropy, glasses, random-singlet phases and non-Fermi liquids are described
in the same terms. Where both apply, the formulation is equivalent to the
scaling hypothesis for the free energy; it relates the exponents to the
spectrum and does not compute the spectrum.
\end{abstract}

\tableofcontents

\section{Introduction}
\label{sec:intro}

The equilibrium thermodynamics of a many-body Hamiltonian $\hat H$ follows
from its spectrum in two well-known steps.
First, the grand partition function is the
Laplace transform of the many-body density of states $\rho(\enrg)$,
Eq.~\eqref{eq:filt_laplace}; this is exact for a system of any size. Second,
for a system of large size $V$ the number of levels is exponential in $V$,
$\rho_V(V\varepsilon)\asymp e^{V\lds(\varepsilon)/k_B}$ (equality of the
logarithms to leading order in $V$), and the Laplace transform reduces to a
Legendre transform: the thermodynamic potential per unit size is
$\omega(T)=\min_\varepsilon[\varepsilon-T\lds(\varepsilon)]$.
The function $\lds(\varepsilon)$, the
\emph{logarithmic density of states} (hereafter \emph{\lndos}), corresponds to the
microcanonical entropy per unit size, Boltzmann's $k_B\ln W$ with $W$ written
in terms of the level count, and the second step is basically the equivalence of the
ensembles.

In the usual presentation the assumptions that produce power laws and
exponents are made on the free energy: the scaling
hypothesis~\cite{Widom1965} postulates that its singular\footnote{A function
is \emph{analytic} at a point if it equals its Taylor series in a
neighbourhood of the point. A point at which it is not analytic is a
\emph{singularity} of the function: a jump, a kink (a jump of the first
derivative), a divergence, or a point at which all derivatives are finite
but the Taylor series does not sum to the function (an \emph{essential
singularity}, such as $e^{-1/t}$ at $t=0$). The \emph{singular part} of a
function is what remains after an analytic function has been subtracted. A
phase transition is a singularity of a thermodynamic potential as a function
of $T$ or of a parameter of the Hamiltonian.} part is a generalised
homogeneous function, and Landau's theory expands it in an order parameter.
The aim of this note is to make the assumptions on the spectrum instead, and
to derive the thermodynamic statements from them. The thermodynamic
potential is then a secondary quantity: at the temperature $T$ it depends on
the spectrum only through the form of $\lds$ at the single energy
$\varepsilon_*(T)$ at which the slope of $\lds$ is $1/T$, and every
singularity, power law and exponent of the thermodynamics is a property of
the form of $\lds$ near that energy. The article is organised as follows.
\begin{enumerate}[label=(\roman*),itemsep=2pt]
\item \emph{Exact relations} (Section~\ref{sec:filters}). The density of
  states $\rho$ is a distribution, and the thermodynamic quantities of a
  finite system are its pairings with test functions built from the
  Boltzmann factor. A field moves the levels without creating or removing
  any, and the static responses are derivatives of $\rho$ with respect to
  the field.
\item \emph{Large size} (Section~\ref{sec:large}). The Legendre
  transform~\eqref{eq:filt_legendre} gives the entropy, the heat capacity
  and the responses in terms of $\lds$ and its derivatives at
  $\varepsilon_*$: the heat capacity is the inverse curvature of the \lndos,
  Eq.~\eqref{eq:filt_curvmain}.
\item \emph{Transitions and crossovers} (Section~\ref{sec:transitions}) are
  classified by the motion of $\varepsilon_*(T)$ through the spectrum.
\item \emph{Power laws and exponents} (Section~\ref{sec:filt_lowenergy}). A
  power law of $\lds$ above the ground state (Assumption~\ref{ass:power})
  gives power laws in $T$; homogeneity in energy and field
  (Assumption~\ref{ass:hom0}) gives the scaling functions of a quantum
  critical point, worked out for the Lifshitz transition, which has no order
  parameter; homogeneity about an energy in the interior of the
  spectrum (Assumption~\ref{ass:crit}) gives the critical exponents and the
  scaling relations; homogeneity in the wave vector and in the size
  (Assumptions~\ref{ass:q} and~\ref{ass:V}) gives the correlation length,
  finite-size scaling and hyperscaling.
\item \emph{Applications} (Section~\ref{sec:beyond}) to spectra whose form is
  not due to a broken symmetry: degenerate manifolds of wave vectors,
  competing orders, residual entropy, disorder and non-Fermi liquids.
\end{enumerate}

\paragraph{Contribution and relation to earlier work.} The contribution of
this note is a reformulation of phase transitions, crossovers and power laws
directly in terms of the microscopic spectrum, instead of the free energy.
The objects of the description are the features of the graph of $\lds$ over
the energy axis: its slope, its curvature, its lower and upper edges, the
intervals on which it is linear and the points at which it is non-analytic.
They are referred to below as the \emph{form} of the \lndos, and a
thermodynamic phenomenon is a property of this form at the energy selected
by the temperature. Since the form is defined for every Hamiltonian, the
description does not presuppose a broken symmetry or an order parameter, and
it covers the cases outside Landau's paradigm in the same terms as those
inside it (Section~\ref{sec:beyond}).

The individual ingredients are known. The reduction of the Laplace transform
to a Legendre transform (Section~\ref{sec:large}) is the equivalence of the
ensembles~\cite{Schwabl,Ruelle} and, in the language of probability, a
large-deviation principle~\cite{Touchette2009}. The stretched-exponential
growth of the level count of a Fermi gas is Bethe's
formula~\cite{Bethe1936}, and its counterpart for conformal field theories in
one dimension is Cardy's formula~\cite{Cardy1986}. The closest earlier work
is microcanonical thermodynamics, which locates and classifies the
transitions of finite systems by their microcanonical
entropy~\cite{Gross2001,Behringer2006,QiBachmann2018}, obtained
numerically~\cite{WangLandau2001}, and analyses its finite-size scaling at a
critical point, where the exponents with respect to the energy are the usual
ones divided by $1-\alpha$ if the heat capacity
diverges~\cite{Kastner2000}. The present formulation differs in three
respects. It is set in the limit of large size, in which $\lds$ is a property
of the Hamiltonian alone and non-analyticity is sharply defined. Its
assumptions are homogeneity properties of the level count in the energy, the
field, and the size, from which the exponents, the
correlation length, finite-size scaling and hyperscaling follow by one
method. And it applies the same description at the lower edge of the
spectrum, to low-temperature power laws, quantum criticality, frustration,
disorder and non-Fermi liquids. Where the scaling hypothesis applies, the
two formulations are equivalent (Section~\ref{sec:filt_rg}).

\paragraph{Conventions.} The representation is the eigenbasis of the full
many-body Hamiltonian, which is assumed to be known; the formulation relates
the exponents to the spectrum and does not compute the spectrum. The
chemical potential $\mu$ is fixed throughout. A parameter of the Hamiltonian
(a field or a coupling constant) is denoted by $\hfield$.

\section{Exact relations for a finite system}
\label{sec:filters}
This section assumes nothing about $\hat H$ except that its spectrum is
bounded from below. All statements are identities for a system of finite
size.

\subsection{Spectrum, equilibrium state and partition function}
\label{sec:filt_exact}\label{sec:filt_laplace}

Let $ \hat H $ be the many-body Hamiltonian, interacting or not, of a system
whose total particle number is conserved\footnote{The
conservation law is required only for the species to which $ \mu $ is
attached. Phonons, photons and magnons are not conserved and are treated
with $ \mu = 0 $. For a Hamiltonian with no conserved number, such as the
Bardeen--Cooper--Schrieffer mean-field Hamiltonian, the construction applies
directly to $ \hat{\mathcal{H}} = \hat H_{\rm MF} - \mu\hat N $, and
$ \mu $ is fixed afterwards by $ \langle\hat N\rangle = N $.},
$ [\hat H, \hat N] = 0 $. The two operators have a common orthonormal
eigenbasis $ \{|a\rangle\} $ of the Fock space, and we abbreviate
\[
  \hat{\mathcal{H}} := \hat H - \mu\hat N , \qquad
  \hat{\mathcal{H}}|a\rangle = \enrg_a|a\rangle , \qquad \enrg_a = E_a - \mu N_a .
\]
The lowest eigenvalue is $\enrg_0$. In contact with a reservoir at
temperature $T$ and chemical potential $\mu$ the equilibrium state is
\[
  \hat\rho_{\rm eq} = \frac{e^{-\beta\hat{\mathcal{H}}}}{\Xi}
  = \sum_ap_a\,|a\rangle\langle a| , \qquad
  p_a = \frac{e^{-\beta \enrg_a}}{\Xi} , \qquad
  \Xi = \Tr e^{-\beta\hat{\mathcal{H}}} = \sum_ae^{-\beta\enrg_a} , \qquad
  \beta=\frac1{k_BT} ,
\]
with the grand partition function $\Xi$ and the grand potential
$\Omega=-k_BT\ln\Xi$.

The spectrum of $\hat{\mathcal H}$ is described by the \emph{many-body
density of states} $\rho$ or, equivalently, by the \emph{level count} $\cN$,
\begin{equation}\label{eq:filt_gN}
  \rho(\enrg) := \sum_a\delta(\enrg - \enrg_a) , \qquad
  \cN(\enrg) := \sum_a\Theta(\enrg-\enrg_a) = \int_{-\infty}^{\enrg}d\enrg'\,\rho(\enrg') ,
\end{equation}
with $\Theta$ the Heaviside step function: $\cN(\enrg)$ is the number of
eigenstates with $\enrg_a\le\enrg$, a non-decreasing step function, and
$\rho=d\cN/d\enrg$. Then
\begin{equation}\label{eq:filt_laplace}
  \Xi(\beta) = \int_{-\infty}^{+\infty} d\enrg\; \rho(\enrg)\,e^{-\beta \enrg}
  = \beta\int_{-\infty}^{+\infty} d\enrg\;\cN(\enrg)\,e^{-\beta\enrg} :
\end{equation}
the grand partition function is the Laplace transform of the many-body
density of states, with $\beta$ the variable conjugate to $\enrg$. The second
form follows from the first by partial integration\footnote{The boundary
term $[\cN(\enrg)e^{-\beta\enrg}]_{-\infty}^{\infty}$ vanishes: $\cN=0$ below
the ground state, and if $\Xi$ converges at some $0<\beta'<\beta$ (always,
for a finite Fock space), the $\cN(\enrg)$ levels below $\enrg$ contribute at
least $\cN(\enrg)e^{-\beta'\enrg}$ to $\Xi(\beta')$, so that
$\cN(\enrg)e^{-\beta\enrg}\le\Xi(\beta')e^{-(\beta-\beta')\enrg}\to0$ for
$\enrg\to\infty$.}. The probability density of the energy in equilibrium is
the density of states weighted by the Boltzmann factor,
\[
  P_T(\enrg) := \frac{\rho(\enrg)\,e^{-\beta \enrg}}{\Xi} .
\]
A parameter $\hfield$ of the Hamiltonian changes the levels,
$\enrg_a(\hfield)$, and hence $\rho(\enrg(\hfield))$ and $\cN(\enrg(\hfield))$;
the Boltzmann factor $e^{-\beta\enrg}$ does not depend on it.

\subsection{The density of states as a distribution}
\label{sec:filt_dist}

A sum of delta functions has no values; $\rho$ is a \emph{distribution}
(generalised function)\footnote{\label{footnote:distribution}
A \emph{distribution} is a linear functional
$\varphi\mapsto\langle\rho,\varphi\rangle$ on a space of smooth \emph{test
functions}. An ordinary function $F$ with finite integrals over bounded
intervals defines one by $\langle F,\varphi\rangle=\int d\enrg\,F\varphi$;
the delta function $\delta(\enrg-\enrg_a)$ is the functional
$\varphi\mapsto\varphi(\enrg_a)$, which is not of this form. The derivative
of a distribution is defined by
$\langle\rho',\varphi\rangle:=-\langle\rho,\varphi'\rangle$~\cite{GelfandShilov,Zemanian}.}:
the linear functional that assigns to a test function $\varphi$ the number
\begin{equation}\label{eq:filt_pairing}
  \langle \rho,\varphi\rangle:=\int d\enrg\;\rho(\enrg)\,\varphi(\enrg)
  =\sum_a\varphi(\enrg_a)=\Tr\varphi(\hat{\mathcal H}) .
\end{equation}
The last form is independent of the basis. Since the support of $\rho$ is
bounded below, the pairing is defined for every continuous $\varphi$ that
decays for $\enrg\to+\infty$ faster than the levels accumulate; no condition
is needed if the Fock space has a finite dimension $D$. Two properties are
used below.
\begin{itemize}\setlength{\itemsep}{2pt}
\item \emph{Level count.} $\rho=d\cN/d\enrg$ in the sense of distributions.
  Unlike $\rho$, the level count $\cN$ is an ordinary function, with values
  that can be bounded and compared. The assumptions on the spectrum
  (Sections~\ref{sec:large} and~\ref{sec:filt_lowenergy}) are therefore
  stated for $\cN$. A value of $\rho$, the number of levels per unit energy,
  is used only for large size, where it is the derivative of the smooth
  function that replaces $\cN$ to logarithmic accuracy
  (Eq.~\eqref{eq:filt_rhosmooth}).
\item \emph{Dependence on a parameter.} Let
  $\hat{\mathcal H}(\hfield)=\hat{\mathcal H}-\hfield\hat A$ with a Hermitian
  operator $\hat A$. A level moves with the field at the rate
  $\partial\enrg_a/\partial \hfield=-A_{aa}$, with
  $A_{aa}=\langle a|\hat A|a\rangle$ in the eigenbasis at $\hfield$ (the
  Hellmann--Feynman theorem, Appendix~\ref{app:lehmann}). Hence
  $\partial_\hfield\sum_a\varphi(\enrg_a)=-\sum_aA_{aa}\varphi'(\enrg_a)$ for
  every test function, i.e.
  \begin{equation}\label{eq:filt_gA}
    \frac{\partial \rho}{\partial \hfield}=\frac{\partial \rho_A}{\partial\enrg} ,
    \qquad
    \rho_A(\enrg):=\sum_aA_{aa}\,\delta(\enrg-\enrg_a) ,
    \qquad
    \langle \rho_A,\varphi\rangle=\Tr\big[\hat A\,\varphi(\hat{\mathcal H})\big] ,
  \end{equation}
  with $\rho_A$ the density of states weighted by $\hat A$.
  Eq.~\eqref{eq:filt_gA} holds whether or not $\hat A$ commutes with
  $\hat{\mathcal H}$ (Appendix~\ref{app:lehmann}). It is a continuity
  equation on the energy axis, $\partial_\hfield\rho+\partial_\enrg j=0$, with
  the current of levels
  $j=-\rho_A=\sum_a(\partial\enrg_a/\partial\hfield)\,\delta(\enrg-\enrg_a)$
  (density times velocity, with $\hfield$ in the role of time): a field does
  not create or remove levels, it moves them. In integrated form,
  $\partial_\hfield\cN=\rho_A$: the number of levels below $\enrg$ changes
  only through the levels that cross $\enrg$.
\end{itemize}

\paragraph{Thermodynamics as a Laplace transform.} Thermodynamics pairs
$\rho$ with the one-parameter family of test functions generated by the
Boltzmann factor $e_\beta(\enrg):=e^{-\beta\enrg}$,
\begin{equation}\label{eq:filt_map}
  \Xi(\beta)=\langle \rho,e_\beta\rangle ,
\end{equation}
and every equilibrium quantity is such a pairing: the moments of the energy
are $\langle\rho,\enrg^{n}e_\beta\rangle/\Xi$, and the mean of $\hat A$ is
$\langle\rho_A,e_\beta\rangle/\Xi$
(Sections~\ref{sec:filt_mb} and~\ref{sec:filt_resp}). The map
$\rho\mapsto\Xi(\beta)$ has four properties.
\begin{itemize}\setlength{\itemsep}{2pt}
\item \emph{Linearity and additivity.} The map is linear in $\rho$. For two
  systems without mutual interaction the levels add,
  $\enrg=\enrg^{(1)}_{a_1}+\enrg^{(2)}_{a_2}$, so that the density of states is
  the convolution\footnote{For sums of delta functions, $\rho_1*\rho_2$ is the
  sum of delta functions located at the sums of the levels:
  $\langle\rho_1*\rho_2,\varphi\rangle=\sum_{a_1,a_2}\varphi(\enrg^{(1)}_{a_1}+\enrg^{(2)}_{a_2})$.
  It counts the ways of sharing the energy between the two
  systems~\cite{Zemanian}.} $\rho=\rho_1*\rho_2$. Since
  $e_\beta(\enrg'+\enrg'')=e_\beta(\enrg')\,e_\beta(\enrg'')$, the Laplace
  transform of the convolution is the product $\Xi=\Xi_1\Xi_2$, and
  $\Omega=\Omega_1+\Omega_2$. The exponentials are the only continuous weight
  functions with this property.
\item \emph{Regularisation.} The sum $\Xi(\beta)=\sum_ae^{-\beta\enrg_a}$ is an
  analytic function of $\beta$ wherever it converges. The thermal quantities
  of a finite system are therefore smooth functions of $T$; a phase
  transition can arise only in the limit $V\to\infty$.
\item \emph{Injectivity (1-to-1 property).} A distribution with support bounded below is
  determined by its Laplace transform (Lerch's theorem~\cite{Widder}): two
  spectra with the same $\Xi(\beta)$ for all $\beta$ are the same spectrum.
\item \emph{Finite resolution.} The test functions $\enrg^{n}e_\beta$ vary on
  the scale $k_BT$. Thermodynamics at temperature $T$ does not resolve
  structure of $\rho$ finer than $k_BT$; this requires test functions that
  are narrow in energy, i.e.\ a spectroscopy. The reconstruction of $\rho$
  from $\Xi(\beta)$ at real $\beta$ is unique but unstable
  (Appendix~\ref{app:generating}).
\end{itemize}

\subsection{Energy, entropy and heat capacity}
\label{sec:filt_mb}

Since $\hat{\mathcal H}$ does not depend on $T$ at fixed $\mu$, each
derivative of Eq.~\eqref{eq:filt_laplace} with respect to $\beta$ produces a
factor $-\enrg$:
\begin{equation}\label{eq:filt_moments}
  \langle\hat{\mathcal H}^{n}\rangle
  = \frac{(-1)^{n}}{\Xi}\frac{\partial^{n}\Xi}{\partial\beta^{n}}\evalat_{\mu}
  = \frac{1}{\Xi}\int d\enrg\;\enrg^{n}\rho(\enrg)\,e^{-\beta\enrg} ,
  \quad
  \langle\hat{\mathcal H}\rangle = -\frac{\partial\ln\Xi}{\partial\beta}\evalat_{\mu} ,
  \quad
  \Var(\hat{\mathcal H}) = \frac{\partial^{2}\ln\Xi}{\partial\beta^{2}}\evalat_{\mu} .
\end{equation}
The entropy of $\hat\rho_{\rm eq}$ is $S=-k_B\sum_ap_a\ln p_a$; with
$\ln p_a=-\beta\enrg_a-\ln\Xi$,
\begin{equation}\label{eq:filt_mb}
  S = k_B\Big(\ln\Xi - \beta\,\frac{\partial\ln\Xi}{\partial\beta}\Big)
  = \frac{\langle\hat{\mathcal{H}}\rangle - \Omega}{T} , \qquad
  C_\mu := T\,\frac{\partial S}{\partial T}\bigg|_{\mu}
  = k_B\beta^{2}\,\frac{\partial^{2}\ln\Xi}{\partial\beta^{2}}
  = \frac{\Var(\hat{\mathcal{H}})}{k_BT^{2}} .
\end{equation}
Equivalently, since $T\,dS=d\langle\hat{\mathcal H}\rangle$ at fixed $\mu$,
\begin{equation}\label{eq:filt_CU}
  C_\mu = \frac{\partial\langle\hat{\mathcal H}\rangle}{\partial T}\bigg|_{\mu}
  = \frac{1}{k_BT^{2}}\int d\enrg\; P_T(\enrg)\,\big(\enrg - \langle\hat{\mathcal{H}}\rangle\big)^{2}\;\ge0 .
\end{equation}
The heat capacity at fixed $\mu$ is the variance of the
energy.\footnote{$C_\mu$ differs from the heat capacity at fixed particle
number,
$C_N=C_\mu-T(\partial N/\partial T)_\mu^{2}/(\partial N/\partial\mu)_T$. The
two have the same low-temperature exponents in the examples of
Section~\ref{sec:filt_power_mb}; at a critical point the constraint of a
fixed density can change the exponents (Fisher
renormalisation~\cite{Fisher1968}).}

\subsection{Response and correlation functions}
\label{sec:filt_corr}\label{sec:filt_resp}

Let fields $\hfield_i$ couple to Hermitian operators $\hat A_i$,
$ \hat{\mathcal{H}}(\{\hfield_i\}) = \hat{\mathcal{H}} - \sum_i\hfield_i\hat A_i $.
Since the Boltzmann factor in Eq.~\eqref{eq:filt_laplace} is independent of
the fields, every derivative with respect to a field acts on $ \rho $ alone:
\begin{equation}\label{eq:Xi-derivatives}
  \Xi^{(i)} := \frac{\partial\Xi}{\partial \hfield_i}
  = \int d\enrg\;\frac{\partial\rho}{\partial \hfield_i}\,e^{-\beta\enrg} ,
  \qquad
  \Xi^{(ij)} := \frac{\partial^{2}\Xi}{\partial \hfield_i\,\partial \hfield_j}
  = \int d\enrg\;\frac{\partial^{2}\rho}{\partial \hfield_i\,\partial \hfield_j}\,e^{-\beta\enrg} .
\end{equation}
By Eq.~\eqref{eq:filt_gA} and a partial integration,
$\Xi^{(i)}=\langle\partial_\enrg\rho_{A_i},e_\beta\rangle
=\beta\langle\rho_{A_i},e_\beta\rangle=\beta\,\Xi\,\langle\hat A_i\rangle$,
whether or not the $\hat A_i$ commute with $\hat{\mathcal H}$ and with each
other. Hence $\langle\hat A_i\rangle=-\partial\Omega/\partial \hfield_i$ and
\begin{equation}\label{eq:filt_param}
  \langle\hat A_i\rangle = k_BT\,\frac{\Xi^{(i)}}{\Xi} , \qquad
  \chi_{ij} := \frac{\partial\langle\hat A_i\rangle}{\partial \hfield_j}\bigg|_{T,\mu}
  = k_BT\left[\frac{\Xi^{(ij)}}{\Xi} - \frac{\Xi^{(i)}\,\Xi^{(j)}}{\Xi^{2}}\right] :
\end{equation}
the conjugate quantity is the Laplace transform of
$\partial \rho/\partial \hfield_i$, and the static response function
$\chi_{ij}=\chi_{ji}$ that of the mixed second derivative of $\rho$ less the
product of the first ones, each normalised by $ \Xi $. In terms of the
levels (Appendix~\ref{app:lehmann}),
\begin{equation}\label{eq:filt_corr}
  \chi_{ij}=\frac{1}{k_BT}\Big[\sum_ap_a\,A^{i}_{aa}A^{j}_{aa}
  -\langle\hat A_i\rangle\langle\hat A_j\rangle\Big]
  -\sum_ap_a\,\frac{\partial^{2}\enrg_a}{\partial \hfield_i\,\partial \hfield_j} ,
  \qquad A^{i}_{aa}:=\langle a|\hat A_i|a\rangle :
\end{equation}
the covariance, over the Boltzmann-weighted levels, of the slopes
$\partial\enrg_a/\partial\hfield_i=-A^{i}_{aa}$ of the levels with respect to
the two fields, plus the mean curvature of the levels.

\paragraph{Correlation function.} Let $\hat A_i=\hat m_i$ be the
order-parameter operator of the component (site) $i$ at position
$\mathbf r_i$, and $\hfield_i$ the local field at this site. If all
$\hat m_i$ commute with $\hat{\mathcal H}$ (a classical model such as the
Ising model), the fields shift the levels linearly, the last term of
Eq.~\eqref{eq:filt_corr} vanishes, and
$k_BT\chi_{ij}=G_{ij}:=\langle\hat m_i\hat m_j\rangle
-\langle\hat m_i\rangle\langle\hat m_j\rangle$ exactly\footnote{In a
quantum model $k_BT\chi_{ij}$ is the zero-frequency part of the correlation
function and differs from the equal-time $G_{ij}$
(Appendix~\ref{app:lehmann}).}: two sites are positively correlated if the
levels that are lowered by a field at $i$ are also lowered by a field at
$j$.

\paragraph{Modulated field.} For a translation-invariant system (a lattice
of $V$ sites with periodic boundary conditions)
$\chi_{ij}=\chi(\mathbf r_i-\mathbf r_j)$, with the Fourier transform
$\chi(\mathbf q)=\sum_{\mathbf r}e^{-i\mathbf q\cdot\mathbf r}\chi(\mathbf r)$
at the wave vectors allowed by the boundary conditions. A modulated field
$\hfield_i=\hfield_{\mathbf q}\cos(\mathbf q\cdot\mathbf r_i)$ of amplitude
$\hfield_{\mathbf q}$ is a single parameter of the Hamiltonian, and
\begin{equation}\label{eq:filt_chiq_exact}
  k_BT\,\frac{\partial^{2}\ln\Xi}{\partial \hfield_{\mathbf q}^{2}}\bigg|_{\hfield_{\mathbf q}=0}
  =\sum_{ij}\cos(\mathbf q\cdot\mathbf r_i)\cos(\mathbf q\cdot\mathbf r_j)\,\chi_{ij}
  =\tfrac12V\chi(\mathbf q)
\end{equation}
for every allowed $\mathbf q$ such that $2\mathbf q$ is not a
reciprocal-lattice vector\footnote{By
$\cos a\cos b=\tfrac12[\cos(a-b)+\cos(a+b)]$, the first term gives
$\tfrac12V\chi(\mathbf q)$, and the second vanishes because the sum of
$e^{2i\mathbf q\cdot\mathbf r_j}$ over the sites of a periodic lattice
vanishes unless $2\mathbf q$ is a reciprocal-lattice vector. In the latter
case (e.g.\ $\mathbf q=0$) the factor $\tfrac12$, the mean of $\cos^{2}$, is
replaced by $1$.}. The susceptibility at wave vector $\mathbf q$ is the
second derivative of $\ln\Xi$ with respect to the amplitude of a modulated
field.

\section{The limit of large size}
\label{sec:large}

For a macroscopic system the Laplace transform~\eqref{eq:filt_laplace}
simplifies, because the number of levels is exponential in the size: the
thermodynamic potential becomes the Legendre transform of the \lndos. This
section restates the equivalence of the microcanonical and the grand
canonical ensembles~\cite[Ch.~2]{Schwabl}, \cite{Ruelle} in the terms of
Section~\ref{sec:filters}, so that the observables can be expressed through
the spectrum.

\subsection{The logarithmic density of states}
\label{sec:large_ext}

The \emph{size} $V$ of a system is the number of its components (lattice
sites, spins, orbitals) at fixed boundary conditions\footnote{A volume is
not a property of a Hamiltonian. A system of size $V$ is one member of a
sequence of Hamiltonians $\hat H_V$ built from the same local couplings with
the same boundary conditions, e.g.\ a hypercubic lattice of $V=L^{d}$ sites
with periodic boundary conditions, and ``large size'' is the limit of this
sequence. For particles in a continuum, $V$ is the volume of the box in
units of a microscopic volume. A length enters only through the geometry of
the couplings, as $L=V^{1/d}$ lattice spacings.}. The levels of a lattice
system with finitely many states per component and couplings of bounded
strength are $D=e^{O(V)}$ in number\footnote{$f=O(g)$ means that $|f/g|$
stays bounded in the limit considered, and $f=o(g)$ that $f/g\to0$.} and lie
in an energy range of width $O(V)$. The density of states $\rho_V$ therefore
has no limit for $V\to\infty$: the number of levels in a fixed range of
$\enrg/V$ grows as $e^{O(V)}$. What has a limit is the logarithm of the level
count per unit size, at a fixed energy per unit size
$\varepsilon:=\enrg/V$. For every finite system
\begin{equation}\label{eq:filt_sV}
  \lds_V(\varepsilon):=\frac{k_B}{V}\ln\cN_V(V\varepsilon) ,
  \qquad\text{i.e.}\qquad
  \cN_V(\enrg)=\exp\Big[\frac{V}{k_B}\,\lds_V\Big(\frac{\enrg}{V}\Big)\Big] ,
\end{equation}
is a non-negative, non-decreasing step function above the ground state, and
the second form of the Laplace transform~\eqref{eq:filt_laplace} reads
exactly
\begin{equation}\label{eq:filt_laplace_V}
  \Xi=\beta V\int_{\enrg_0/V}^{\infty}d\varepsilon\;
  \exp\left\{-\frac{V}{k_BT}\big[\varepsilon-T\,\lds_V(\varepsilon)\big]\right\} .
\end{equation}
The exponent is proportional to $V$. It is the sum of the logarithm of the
Boltzmann factor, which decreases with the energy, and the logarithm of the
level count, which increases.

For a macroscopic system the ground-state energy per unit size
$\varepsilon_0:=\lim_{V\to\infty}\enrg_0/V$ exists, and so does the
\emph{logarithmic density of states} (\lndos)
\begin{equation}\label{eq:filt_sdef}
  \lds(\varepsilon) := \lim_{V\to\infty}\lds_V(\varepsilon)
  =\lim_{V\to\infty}\frac{k_B}{V}\,\ln\cN_V(V\varepsilon)
  \qquad(\varepsilon>\varepsilon_0) ,
\end{equation}
uniformly\footnote{\emph{Uniformly}: the largest deviation
$\max_\varepsilon|\lds_V-\lds|$ on the interval tends to zero. This allows
the minimum over $\varepsilon$ in Eq.~\eqref{eq:filt_legendre} to be taken
after the limit $V\to\infty$.} on bounded intervals of $\varepsilon$, with
the linear bound
\begin{equation}\label{eq:filt_linbound}
  \lds(\varepsilon)\le c_0+\frac{\varepsilon}{\Tmax}
\end{equation}
for some constants $c_0$ and $0<\Tmax\le\infty$ (for a Fock space of finite
dimension $D=q^{V}$, $\lds\le k_B\ln q$ and $\Tmax=\infty$). For short-range
interactions the existence of the limit is a theorem~\cite{Ruelle}, proven
by dividing the system into parts whose mutual (surface) interaction is
negligible, and $\lds$ is \emph{concave}\footnote{A function is
\emph{concave} if its graph lies on or above each of its chords; if it is
twice differentiable, its second derivative is nowhere positive.}. For
interactions that decay more slowly than $r^{-d}$ in $d$ dimensions the
limit exists after the couplings are rescaled with the size (Kac
prescription~\cite{Campa2009}), and $\lds$ need not be concave. The results
below do not require concavity unless stated.

\paragraph{Properties.} The number of levels below the energy $V\varepsilon$
is $e^{V\lds(\varepsilon)/k_B}$ up to a factor that is not exponential in
$V$. Since $\cN_V\ge1$ is non-decreasing, $\lds$ is non-negative and
non-decreasing. Its limit
$\lds_0:=\lim_{\varepsilon\to\varepsilon_0^{+}}\lds(\varepsilon)$ is the
residual entropy per unit size; it is non-zero only if the number of levels
within an energy $o(V)$ of the ground state is exponential in $V$. Where
$\lds$ is differentiable with a positive slope\footnote{For a finite Fock
space $\lds$ is constant above the energy at which the levels are densest.
Defined through the level count, it describes the part of the spectrum
below that energy, which is the part selected by positive temperatures.},
the level count may be replaced by the smooth function
$e^{V\lds(\enrg/V)/k_B}$, whose derivative is an ordinary function,
\begin{equation}\label{eq:filt_rhosmooth}
  \tilde\rho_V(\enrg):=\frac{d}{d\enrg}\,e^{V\lds(\enrg/V)/k_B}
  =\frac{1}{k_BT_\varepsilon}\,e^{V\lds(\varepsilon)/k_B} ,
  \qquad
  \frac{1}{T_\varepsilon}:=\frac{\partial\lds}{\partial\varepsilon} ,
\end{equation}
so that the number of levels per unit energy obeys
\begin{equation}\label{eq:filt_rhos}
  \rho_V(V\varepsilon)\asymp e^{V\lds(\varepsilon)/k_B} ,
\end{equation}
where $\asymp$ denotes equality \emph{to logarithmic
accuracy}\footnote{\label{fn:logacc}Two positive quantities agree to
logarithmic accuracy if their logarithms agree to leading order in $V$. A
factor that is not exponential in $V$ is dropped: since $\Xi$ is linear in
$\cN$, a factor between $V^{-p}$ and $V^{p}$ changes $\ln\Xi$ by at most
$p\ln V$ and $\Omega/V$ by $O(\ln V/V)$.}. Likewise the number of levels in a
shell of width $\Delta\ll V$ below $\enrg$ is
$\cN_V(\enrg)-\cN_V(\enrg-\Delta)=\cN_V(\enrg)\,(1-e^{-\Delta/k_BT_\varepsilon})$,
which differs from $\cN_V$ by a factor that is not exponential in $V$.
Boltzmann's entropy $k_B\ln W$, with $W$ the number of levels in a shell, is
therefore $V\lds(\varepsilon)+O(1)$ for every width of the
shell~\cite[Sec.~2.3]{Schwabl}: \emph{the logarithmic density of states is the microcanonical
entropy per unit size}.

If the Hamiltonian depends on a parameter $\hfield$, the level count and the
\lndos{} depend on it only through the levels $\enrg_a(\hfield)$. As for
$\rho(\enrg(\hfield))$ in Section~\ref{sec:filt_exact}, this is denoted by
\[
  \cN_V(\enrg(\hfield)):=\sum_a\Theta\big(\enrg-\enrg_a(\hfield)\big) ,
  \qquad
  \lds(\varepsilon(\hfield)):=\lim_{V\to\infty}\frac{k_B}{V}\ln\cN_V(V\varepsilon(\hfield)) .
\]
A derivative with respect to $\hfield$ acts on the levels, at a fixed value
of the energy at which the level count is evaluated:
$\partial\cN_V/\partial\hfield=\rho_A$, Eq.~\eqref{eq:filt_gA}. In the
homogeneity assumptions of Section~\ref{sec:filt_lowenergy}, which concern
the joint dependence on the energy and on the parameter, the same function
is written $\lds(\delta\varepsilon,\hfield)$.

\subsection{The thermodynamic potential as a Legendre transform}
\label{sec:large_laplace}

Extensive quantities are referred to unit size:
\begin{equation}\label{eq:filt_intensive}
  \omega := \frac{\Omega}{V}=-\frac{k_BT}{V}\ln\Xi , \qquad
  u-\mu n:=\frac{\langle\hat{\mathcal H}\rangle}{V} , \qquad
  c_\mu := \frac{C_\mu}{V} , \qquad
  a := \frac{\langle\hat A\rangle}{V} ,
\end{equation}
with $u=\langle\hat H\rangle/V$ the energy density and
$n=\langle\hat N\rangle/V$ the particle density. Dividing
Eqs.~\eqref{eq:filt_moments}--\eqref{eq:filt_CU} by $V$, with
$\ln\Xi=-\beta V\omega$, gives at every size
\begin{equation}\label{eq:filt_exactV}
  u-\mu n=\frac{\partial(\beta\omega)}{\partial\beta}\bigg|_{\mu} ,\qquad
  \frac{S}{V}=-\frac{\partial\omega}{\partial T}\bigg|_{\mu} ,
  \qquad
  c_\mu=-T\,\frac{\partial^{2}\omega}{\partial T^{2}}\bigg|_{\mu} .
\end{equation}
The spectrum enters these relations through $\omega$ alone.

\paragraph{Thermodynamic potential.} For $V\to\infty$, $\lds_V$ in the exact
Eq.~\eqref{eq:filt_laplace_V} is replaced by its limit:
\begin{equation}\label{eq:filt_laplace_log}
  \Xi\asymp V\int_{\varepsilon_0}^{\infty}d\varepsilon\;
  \exp\left\{-\frac{V}{k_BT}\,\fen(\varepsilon)\right\} ,
  \qquad \fen(\varepsilon):=\varepsilon-T\,\lds(\varepsilon) .
\end{equation}
The notation recalls the Helmholtz free energy $F=E-TS$: $V\fen(\varepsilon)$
is the free energy of the levels at the energy $V\varepsilon$, and the
integrand, divided by $\Xi$, is the probability density $P_T$ of the energy
per unit size. By the bound~\eqref{eq:filt_linbound},
$\fen(\varepsilon)\ge(1-T/\Tmax)\,\varepsilon-c_0T$, so that for $T<\Tmax$
the function $\fen$ grows without limit at large $\varepsilon$ and has a
minimum\footnote{If the level count grew faster than exponentially in the
energy, $\fen$ would be unbounded below and $\Xi$ would not exist at any
temperature; if it grows exponentially, $\Xi$ exists only for
$T<\Tmax$.}. Since the exponent is proportional to $V$,
the integral is dominated by the neighbourhood of this minimum (the
steepest-descent, or Laplace, method): it equals $e^{-V(\min\fen)/k_BT}$ up
to a factor that is not exponential in $V$. Let $\varepsilon_*(T)$ be the
minimiser, the \emph{thermal energy}. Then
\begin{equation}\label{eq:filt_legendre}
  \omega(T) = \min_{\varepsilon\ge\varepsilon_0}\big[\varepsilon - T\,\lds(\varepsilon)\big]
  = \varepsilon_* - T\,\lds(\varepsilon_*) ,
\end{equation}
whether or not $\lds$ is concave. If $\varepsilon_*>\varepsilon_0$ and $\lds$ is
differentiable at $\varepsilon_*$, the condition for the minimum is
\begin{equation}\label{eq:filt_stat}
  \frac{\partial \lds}{\partial\varepsilon}\bigg|_{\varepsilon_*}=\frac1T :
\end{equation}
the temperature is the inverse slope of the \lndos{} at the thermal energy;
equivalently, $\partial\ln\cN/\partial\enrg=1/k_BT$. Geometrically
(Fig.~\ref{fig:transitions}), the straight line
\begin{equation}\label{eq:filt_tangent}
  \ell_T(\varepsilon):=\frac{\varepsilon-\omega(T)}{T}
\end{equation}
of slope $1/T$ lies on or above the graph of $\lds$, touches it at
$\varepsilon_*$, and crosses the energy axis at $\omega$. With increasing $T$
the point of contact moves from the ground state ($T\to0$) to the energy at
which $\lds$ reaches its maximum ($T\to\infty$).

\paragraph{Energy, entropy and heat capacity.} Insert
Eq.~\eqref{eq:filt_legendre} into the exact
relations~\eqref{eq:filt_exactV}.\footnote{The derivatives may be taken after
the limit $V\to\infty$: $\ln\Xi$ is a convex function of $\beta$ at every
size, since $\partial^{2}\ln\Xi/\partial\beta^{2}=\Var(\hat{\mathcal H})\ge0$,
and the derivatives of a convergent sequence of convex functions converge to
the derivative of the limit wherever the latter exists.}
Since $\fen(\varepsilon)$ is stationary at $\varepsilon_*$, only its explicit
dependence on $T$ contributes to the derivative of the minimum
(Lemma~\ref{lem:envelope}), $\partial\omega/\partial T=-\lds(\varepsilon_*)$,
and
\begin{equation}\label{eq:filt_SI}
  \frac{S}{V}=\lds\big(\varepsilon_*(T)\big) , \qquad
  u-\mu n=\omega+T\,\frac{S}{V}=\varepsilon_*(T) , \qquad
  c_\mu=\frac{d\varepsilon_*}{dT} .
\end{equation}
The minimiser is the thermodynamic energy density, and the \lndos{} at the
minimiser is the thermodynamic entropy density:
Eq.~\eqref{eq:filt_legendre} is $\Omega=U-TS-\mu N$ with $S$ the
microcanonical entropy at the mean energy, the equivalence of the ensembles.
Differentiating Eq.~\eqref{eq:filt_stat} with respect to $T$ gives
$\lds_{\varepsilon\varepsilon}\,d\varepsilon_*/dT=-1/T^{2}$, with
$\lds_{\varepsilon\varepsilon}:=\partial^{2}\lds/\partial\varepsilon^{2}$ the
\emph{curvature} of the \lndos. Hence, wherever $\lds$ is twice
differentiable at $\varepsilon_*$,
\begin{equation}\label{eq:filt_curvmain}
  c_\mu=\frac{d\varepsilon_*}{dT}
  =-\frac{1}{T^{2}\,\lds_{\varepsilon\varepsilon}(\varepsilon_*)} :
\end{equation}
the heat capacity is the inverse curvature of the \lndos. A large heat
capacity corresponds to a range of energies over which the slope of $\lds$
is nearly constant. The same curvature fixes the distribution of the
energy: the expansion of $\fen$ about its minimum gives
$P_T\propto\exp[-V|\lds_{\varepsilon\varepsilon}|(\varepsilon-\varepsilon_*)^{2}/2k_B]$,
a Gaussian of variance $k_B/V|\lds_{\varepsilon\varepsilon}|=k_BT^{2}c_\mu/V$,
in agreement with the exact $\Var(\hat{\mathcal H})=k_BT^{2}C_\mu$ of
Eq.~\eqref{eq:filt_mb}.

\paragraph{Response to a field.} For $\lds=\lds(\varepsilon(\hfield))$ the
minimiser is $\varepsilon_*(T,\hfield)$, and by the same argument only the
explicit dependence of $\fen$ on $\hfield$ contributes to
$\partial\omega/\partial\hfield$. The density conjugate to the field is
\begin{equation}\label{eq:filt_af}
  a=-\frac{\partial\omega}{\partial \hfield}\bigg|_{T}
  =T\,\frac{\partial \lds}{\partial \hfield}\bigg|_{\varepsilon_*} .
\end{equation}
This is $\langle\hat A\rangle=k_BT\,\Xi^{(1)}/\Xi$,
Eq.~\eqref{eq:filt_param}, on the logarithmic scale. Explicitly, with
$\partial_\hfield\cN=\rho_A$ and $\tilde\rho=\cN/k_BT$ at $\varepsilon_*$
(Eq.~\eqref{eq:filt_rhosmooth}),
$T\,\partial\lds/\partial\hfield=(k_BT/V)\,\rho_A/\cN=\rho_A/V\tilde\rho$: the
conjugate density is the mean of $A_{aa}/V$ over the levels at the thermal
energy. Differentiating Eq.~\eqref{eq:filt_af} once more, with
$\partial\varepsilon_*/\partial\hfield=-\lds_{\varepsilon\hfield}/\lds_{\varepsilon\varepsilon}$
from the derivative of Eq.~\eqref{eq:filt_stat}, gives the susceptibility per
unit size,
\begin{equation}\label{eq:filt_chis}
  \chi:=\frac{\partial a}{\partial\hfield}\bigg|_{T}
  =T\left[\lds_{\hfield\hfield}
  -\frac{\lds_{\varepsilon\hfield}^{2}}{\lds_{\varepsilon\varepsilon}}\right]_{\varepsilon_*} ,
\end{equation}
with subscripts denoting partial derivatives.

\paragraph{Modulated field.} The amplitude $\hfield_{\mathbf q}$ of a
modulated field (Section~\ref{sec:filt_corr}) is a parameter like any other.
Let $\lds(\varepsilon(\hfield_{\mathbf q}))$ be the \lndos{} in its
presence. With $\ln\Xi=-V\omega/k_BT$, Eq.~\eqref{eq:filt_chiq_exact} reads
$\chi(\mathbf q)=-2\,\partial^{2}\omega/\partial \hfield_{\mathbf q}^{2}$, and
Eq.~\eqref{eq:filt_chis} gives
\begin{equation}\label{eq:filt_chiq}
  \chi(\mathbf q)=-2\,\frac{\partial^{2}\omega}{\partial \hfield_{\mathbf q}^{2}}
  \bigg|_{\hfield_{\mathbf q}=0}
  =2T\,\frac{\partial^{2}\lds}{\partial \hfield_{\mathbf q}^{2}}
  \bigg|_{\hfield_{\mathbf q}=0,\;\varepsilon=\varepsilon_*} .
\end{equation}
The mixed derivative $\lds_{\varepsilon\hfield_{\mathbf q}}$ of
Eq.~\eqref{eq:filt_chis} vanishes at $\hfield_{\mathbf q}=0$: by
Eq.~\eqref{eq:filt_af}, $T\,\partial\lds/\partial \hfield_{\mathbf q}$ at the
energy $\varepsilon$ is the mean of
$V^{-1}\sum_i\cos(\mathbf q\cdot\mathbf r_i)\,m^{i}_{aa}$, with
$m^{i}_{aa}:=\langle a|\hat m_i|a\rangle$, over the levels at that energy;
in the translation-invariant system the mean of $m^{i}_{aa}$
does not depend on the site, and $\sum_i\cos(\mathbf q\cdot\mathbf r_i)=0$.
The uniform susceptibility is the limit
$\chi=\lim_{\mathbf q\to0}\chi(\mathbf q)$, taken after the limit of large
size; for short-range interactions it coincides with
Eq.~\eqref{eq:filt_chis} for a uniform field.

\paragraph{Corrections at finite size.} Eq.~\eqref{eq:filt_legendre} is exact
for $V\to\infty$. Three contributions to the exponent of the exact
Eq.~\eqref{eq:filt_laplace_V} have been discarded: the prefactor $\beta V$
and the width $O(V^{-1/2})$ of the maximum of the integrand, which add
$O(\ln V)$ to $\ln\Xi$; the difference $\lds_V-\lds$, which adds $o(V)$ (for
short-range interactions and open boundaries a surface term); and the
levels away from $\varepsilon_*$, whose weight relative to the dominant ones
is $e^{-O(V)}$. In $\omega=-(k_BT/V)\ln\Xi$ all three vanish for
$V\to\infty$.

In summary, every thermodynamic quantity at $(T,\hfield)$ is a function of
$\lds$ and of its derivatives at the single point $(\varepsilon_*,\hfield)$
(Table~\ref{tab:relations} in Section~\ref{sec:summary}); the temperature
enters through $\varepsilon_*(T,\hfield)$ and through explicit factors of $T$
alone. A property of $\omega$ is the image of a property of the spectrum: a
power law of $\omega$ in $T$ is a power law of $\ln\cN$ in the excitation
energy, a scaling form of $\omega$ is a homogeneity of $\ln\cN$, and a
singularity of $\omega$ is a singularity of $\lds$ at the thermal energy or
a discontinuity of the thermal energy. The following two sections make
these statements precise.

\section{Phase transitions and crossovers}
\label{sec:transitions}

For a finite system $\Xi(\beta)$ is a sum of analytic functions, and nothing
is singular (Section~\ref{sec:filt_dist}). The limit $V\to\infty$ replaces
the sum by the minimisation~\eqref{eq:filt_legendre}, and the minimum of a
family of functions can be non-analytic in its parameters. If $\fen$ is
analytic at its minimiser, with a non-zero second derivative, and the
minimiser is unique, the stationarity condition~\eqref{eq:filt_stat} can be
solved for an analytic $\varepsilon_*(T,\hfield)$, and $\omega$ is analytic
(Appendix~\ref{sec:tools_power}). The thermodynamic potential is therefore
non-analytic only where
\begin{enumerate}[label=(\alph*),itemsep=0pt]
\item the minimiser $\varepsilon_*$ jumps between two separate minima of
  $\fen$;
\item $\lds$ is non-analytic, or its curvature vanishes, at $\varepsilon_*$;
\item the lower edge $\varepsilon_0(\hfield)$ of the spectrum, at which the
  minimiser lies for $T\to0$, is non-analytic in $\hfield$.
\end{enumerate}
These are the discontinuous, the continuous and the quantum transition
(Fig.~\ref{fig:transitions}). The classification is by the motion of the
thermal energy through the spectrum, not by a symmetry.\footnote{If $\lds$
is analytic, the non-analyticities of the minimum are generically two minima of
equal value and, when a second parameter is tuned, a minimum at which the second and third derivatives of
$\fen$ vanish (singularity theory~\cite{Arnold}). The latter has
$\lds_{\rm sing}\propto-(\varepsilon-\varepsilon_c)^{4}$, i.e.\ $\sigma=4$ in
Assumption~\ref{ass:crit} and $c_\mu\propto|t|^{-2/3}$: the mean-field
critical end point of a line of discontinuous transitions, at which the
energy responds to the temperature as a mean-field order parameter responds
to its field. Any other exponent requires $\lds$ itself to be non-analytic at
$\varepsilon_c$. In convex analysis, $\omega(T)$ is, up to signs, the
Legendre--Fenchel transform of $\lds$; it is concave in $T$ and has a kink
exactly where the minimiser is not unique~\cite{Rockafellar,Touchette2009}.}

\begin{figure}[htbp]
\centering
% common styles: energy to the right; upper panels show the logdos, lower panels psi_T
\tikzset{
  ldsgraph/.style={>=Stealth,font=\small,x=1.02cm,y=1.5cm},
  psigraph/.style={>=Stealth,font=\small,x=1.02cm,y=0.8cm},
  lds/.style={very thick,blue!60!black},
  edge/.style={thick,orange!85!black},
  guide/.style={thin,dotted,gray!80!black},
  alt/.style={very thick,red!75!black,dashed},
  psiL/.style={thick,blue!50},
  psiC/.style={very thick,orange!85!black},
  psiH/.style={thick,red!70!black},
}
% ---------------- upper row: logdos
\begin{tikzpicture}[ldsgraph]
  % panel (a): discontinuous transition, linear stretch of the concave logdos
  \node[font=\normalsize] at (-0.4,2.4) {(a)};
  \draw[->] (-0.3,0) -- (4.4,0) node[below left]{$\varepsilon$};
  \draw[->] (0,-0.2) -- (0,2.4) node[above]{$\lds$};
  \draw[lds,domain=0:1,samples=40,smooth] plot (\x,{sqrt(\x)});
  \draw[lds] (1,1) -- (2.2,1.6);
  \draw[lds,domain=2.2:4.1,samples=40,smooth] plot (\x,{1.6+0.5*(\x-2.2)-0.08*(\x-2.2)^2});
  % non-concave but non-decreasing alternative (long-range interactions)
  \draw[alt,domain=1:2.2,samples=40,smooth] plot (\x,{1+0.5*(\x-1)-0.17*(sin(deg(pi*(\x-1)/1.2)))^2});
  \draw[edge] (-0.2,0.4) -- (3.6,2.3);
  \node[orange!85!black,anchor=south,font=\scriptsize] at (3.5,2.3) {slope $1/T_1$};
  \draw[guide] (1,0) -- (1,1); \draw[guide] (2.2,0) -- (2.2,1.6);
  \draw[guide] (0,1) -- (1,1); \draw[guide] (0,1.6) -- (2.2,1.6);
  \node[below,font=\scriptsize] at (1,0) {$\varepsilon_-$};
  \node[below,font=\scriptsize] at (2.2,0) {$\varepsilon_+$};
  \draw[<->,thin] (0.15,1) -- (0.15,1.6) node[midway,right,font=\scriptsize]{$\Delta s$};
\end{tikzpicture}\hfill
\begin{tikzpicture}[ldsgraph]
  % panel (b): continuous transition, s = s_c + du/T_c - B|du|^3 (sigma = 3, T_c = 2)
  \node[font=\normalsize] at (-0.4,2.4) {(b)};
  \draw[->] (-0.3,0) -- (4.4,0) node[below left]{$\varepsilon$};
  \draw[->] (0,-0.2) -- (0,2.4) node[above]{$\lds$};
  % shaded: the singular part, i.e. the deviation of the logdos from its tangent at eps_c
  \fill[orange!25] plot[domain=0:4.1,samples=80,smooth] (\x,{1.3+0.5*(\x-2)-0.035*abs(\x-2)^3})
     -- (4.1,2.35) -- (0,0.3) -- cycle;
  \draw[lds,domain=0:4.1,samples=80,smooth] plot (\x,{1.3+0.5*(\x-2)-0.035*abs(\x-2)^3});
  \draw[edge] (0,0.3) -- (4.1,2.35);
  \node[orange!85!black,anchor=south east,font=\scriptsize] at (3.3,2.0) {slope $1/T_c$};
  \draw[guide] (2,0) -- (2,1.3);
  \node[below,font=\scriptsize] at (2,0) {$\varepsilon_c$};
  \fill[blue!60!black] (2,1.3) circle (1.5pt);
  \node[anchor=north west,font=\scriptsize] at (2.05,1.2) {$\lds_{\varepsilon\varepsilon}=0$};
  \node[anchor=north west,font=\scriptsize,orange!85!black] at (2.3,0.75) {$\lds_{\rm sing}\propto-|\delta\varepsilon|^{\sigma}$};
  \draw[thin,orange!85!black] (3.0,0.72) -- (3.75,2.08);
\end{tikzpicture}\hfill
\begin{tikzpicture}[ldsgraph]
  % panel (c): crossover, smooth change of form
  \node[font=\normalsize] at (-0.4,2.4) {(c)};
  \draw[->] (-0.3,0) -- (4.4,0) node[below left]{$\varepsilon$};
  \draw[->] (0,-0.2) -- (0,2.4) node[above]{$\lds$};
  \fill[gray!15] (0.7,0) rectangle (1.3,2.3);
  \draw[lds,domain=0:1,samples=40,smooth] plot (\x,{\x^0.75});
  \draw[lds,domain=1:4.1,samples=40,smooth] plot (\x,{1+0.75*(\x-1)-0.12*(\x-1)^2});
  % tangents at a temperature below and above the crossover temperature
  \draw[edge] (-0.2,0.088) -- (1.5,1.38);          % slope 0.76, tangent at x=0.95
  \draw[edge] (1.2,1.534) -- (4.3,2.37);           % slope 0.27, tangent at x=3
  \fill[orange!85!black] (0.95,0.962) circle (1.5pt);
  \fill[orange!85!black] (3,2.02) circle (1.5pt);
  \node[below,font=\scriptsize] at (1,0) {$\varepsilon_\times$};
  \node[font=\scriptsize,anchor=north,blue!60!black] at (3.4,1.8) {other form};
  \node[font=\scriptsize,orange!85!black,anchor=north west] at (1.6,1.35) {$T>T_\times$};
  \node[font=\scriptsize,orange!85!black,anchor=north west] at (0.5,0.42) {$T<T_\times$};
  \node[font=\scriptsize,anchor=south west,blue!60!black] at (0.03,0.8) {$\propto (\delta\varepsilon)^{\sigma}$};
\end{tikzpicture}

\vspace{1mm}

% ---------------- middle row: psi_T = eps - T s(eps) at three temperatures
\begin{tikzpicture}[psigraph]
  % (a): T = 1.5, 2 (= T_1), 2.6
  \draw[->] (-0.3,0) -- (4.4,0) node[above left]{$\varepsilon$};
  \draw[->] (0,-2.5) -- (0,1.1) node[above]{$\fen$};
  \draw[psiL,domain=0:1,samples=40,smooth] plot (\x,{\x-1.5*sqrt(\x)});
  \draw[psiL] (1,-0.5) -- (2.2,-0.2);
  \draw[psiL,domain=2.2:4.1,samples=30,smooth] plot (\x,{\x-1.5*(1.6+0.5*(\x-2.2)-0.08*(\x-2.2)^2)});
  \draw[psiC,domain=0:1,samples=40,smooth] plot (\x,{\x-2*sqrt(\x)});
  \draw[psiC] (1,-1) -- (2.2,-1);
  \draw[psiC,domain=2.2:4.1,samples=30,smooth] plot (\x,{-1+0.16*(\x-2.2)^2});
  \draw[alt,thick,domain=1:2.2,samples=40,smooth] plot (\x,{-1+0.34*(sin(deg(pi*(\x-1)/1.2)))^2});
  \draw[psiH,domain=0:1,samples=40,smooth] plot (\x,{\x-2.6*sqrt(\x)});
  \draw[psiH] (1,-1.6) -- (2.2,-1.96);
  \draw[psiH,domain=2.2:4.1,samples=30,smooth] plot (\x,{\x-2.6*(1.6+0.5*(\x-2.2)-0.08*(\x-2.2)^2)});
  \fill[blue!50] (0.5625,-0.5625) circle (1.5pt);
  \fill[red!70!black] (2.92,-2.07) circle (1.5pt);
  \fill[orange!85!black] (1,-1) circle (1.5pt);
  \fill[orange!85!black] (2.2,-1) circle (1.5pt);
  \node[font=\scriptsize,blue!50,anchor=south east] at (3.3,0.25) {$T<T_1$};
  \node[font=\scriptsize,orange!85!black,anchor=south east] at (3.7,-0.7) {$T=T_1$};
  \node[font=\scriptsize,red!70!black,anchor=south] at (3.6,-1.8) {$T>T_1$};
\end{tikzpicture}\hfill
\begin{tikzpicture}[psigraph]
  % (b): T = 1.6, 2 (= T_c), 2.6
  \draw[->] (-0.3,0) -- (4.4,0) node[above left]{$\varepsilon$};
  \draw[->] (0,-2.5) -- (0,1.1) node[above]{$\fen$};
  \draw[psiL,domain=0:4.1,samples=80,smooth] plot (\x,{\x-1.6*(1.3+0.5*(\x-2)-0.035*abs(\x-2)^3)});
  \draw[psiC,domain=0:4.1,samples=80,smooth] plot (\x,{\x-2*(1.3+0.5*(\x-2)-0.035*abs(\x-2)^3)});
  \draw[psiH,domain=0:4.1,samples=80,smooth] plot (\x,{\x-2.6*(1.3+0.5*(\x-2)-0.035*abs(\x-2)^3)});
  \fill[blue!50] (0.909,-0.225) circle (1.5pt);
  \fill[orange!85!black] (2,-0.6) circle (1.5pt);
  \fill[red!70!black] (3.048,-1.59) circle (1.5pt);
  \draw[guide] (2,0) -- (2,-0.6);
  \node[font=\scriptsize,blue!50,anchor=south east] at (3.2,0.3) {$T<T_c$};
  \node[font=\scriptsize,orange!85!black,anchor=north] at (3.1,-0.68) {$T=T_c$: $\propto|\delta\varepsilon|^{\sigma}$};
  \node[font=\scriptsize,red!70!black,anchor=north] at (3.0,-1.65) {$T>T_c$};
\end{tikzpicture}\hfill
\begin{tikzpicture}[psigraph]
  % (c): T = 1.1, 1.6, 2.5  (T_x = 4/3)
  \draw[->] (-0.3,0) -- (4.4,0) node[above left]{$\varepsilon$};
  \draw[->] (0,-2.5) -- (0,1.1) node[above]{$\fen$};
  \fill[gray!15] (0.7,-2.4) rectangle (1.3,1.0);
  \draw[->] (-0.3,0) -- (4.4,0);
  \draw[psiL,domain=0:1,samples=40,smooth] plot (\x,{\x-1.1*\x^0.75});
  \draw[psiL,domain=1:3.25,samples=40,smooth] plot (\x,{\x-1.1*(1+0.75*(\x-1)-0.12*(\x-1)^2)});
  \draw[psiC,domain=0:1,samples=40,smooth] plot (\x,{\x-1.6*\x^0.75});
  \draw[psiC,domain=1:4.1,samples=40,smooth] plot (\x,{\x-1.6*(1+0.75*(\x-1)-0.12*(\x-1)^2)});
  \draw[psiH,domain=0:1,samples=40,smooth] plot (\x,{\x-2.5*\x^0.75});
  \draw[psiH,domain=1:4.1,samples=40,smooth] plot (\x,{\x-2.5*(1+0.75*(\x-1)-0.12*(\x-1)^2)});
  \fill[blue!50] (0.463,-0.155) circle (1.5pt);
  \fill[orange!85!black] (1.52,-0.65) circle (1.5pt);
  \fill[red!70!black] (2.458,-2.14) circle (1.5pt);
  \node[font=\scriptsize,blue!50,anchor=south east] at (2.6,0.35) {$T<T_\times$};
  \node[font=\scriptsize,red!70!black,anchor=south] at (3.7,-1.55) {$T>T_\times$};
\end{tikzpicture}

\vspace{3mm}

% ---------------- lower row: quantum transition
\begin{tikzpicture}[ldsgraph]
  % panel (d): the lower edge of the spectrum as a function of a parameter;
  % two branches of different slopes, 1.1 + 0.25 x and 1.1 - 0.5 x
  \node[font=\normalsize] at (-2.4,2.4) {(d)};
  \draw[->] (-2.2,1.0) -- (2.3,1.0) node[above left]{$\hfield$};
  \draw[->] (0,-0.2) -- (0,2.4) node[right]{$\varepsilon_0(\hfield)$};
  \draw[alt] (-2,0.6) -- (2,1.6);
  \draw[alt] (-2,2.1) -- (2,0.1);
  \draw[lds] (-2,0.6) -- (0,1.1) -- (2,0.1);
  % finite size: avoided crossing of the two lowest eigenvalues
  \draw[gray!60!black,thick,domain=-2:2,samples=60,smooth] plot (\x,{1.1-0.125*\x-sqrt(0.1406*\x*\x+0.02)});
  \node[anchor=south west,font=\scriptsize] at (0.03,1.25) {$\hfield_c$};
  \draw[guide] (0,1.0) -- (0,1.1);
  \node[font=\scriptsize,anchor=north,gray!40!black] at (-1.0,0.45) {finite $V$: no crossing};
  \node[font=\scriptsize,anchor=east,blue!60!black] at (2.3,2.2) {$V\to\infty$: kink};
\end{tikzpicture}\hspace{1.5cm}
\begin{tikzpicture}[ldsgraph]
  % panel (e): the logdos above the lower edge for three values of the parameter
  \node[font=\normalsize] at (-0.4,2.4) {(e)};
  \draw[->] (-0.3,0) -- (4.4,0) node[below left]{$\varepsilon$};
  \draw[->] (0,-0.2) -- (0,2.4) node[above]{$\lds$};
  \draw[lds,variable=\t,domain=0:1,samples=60,smooth] plot ({0.3+3.8*\t*\t},{0.75*sqrt(3.8)*\t});
  \draw[very thick,teal!80!black,domain=0.95:4.1,samples=60,smooth] plot (\x,{0.55*(\x-0.95)^0.75});
  \draw[psiC,variable=\t,domain=0:1,samples=60,smooth] plot ({1.7+2.4*\t^2.5},{0.95*(2.4)^0.4*\t});
  \node[below,font=\scriptsize] at (1.7,0) {$\varepsilon_0(\hfield_c)$};
  \draw[lds] (2.6,0.75) -- (3.0,0.75) node[right,font=\scriptsize,black]{$\hfield<\hfield_c$};
  \draw[psiC] (2.6,0.5) -- (3.0,0.5) node[right,font=\scriptsize,black]{$\hfield=\hfield_c$};
  \draw[very thick,teal!80!black] (2.6,0.25) -- (3.0,0.25) node[right,font=\scriptsize,black]{$\hfield>\hfield_c$};
\end{tikzpicture}
\caption{The \lndos{} and the transitions of Section~\ref{sec:transitions}.
Upper row: $\lds(\varepsilon)$ (blue) and the tangent of slope $1/T$
(orange), Eq.~\eqref{eq:filt_tangent}. Middle row:
$\fen(\varepsilon)=\varepsilon-T\lds(\varepsilon)$ at three temperatures; the
dot marks its minimum, at $\varepsilon_*(T)$ with the value $\omega(T)$.
(a)~Discontinuous transition: $\lds$ is linear between $\varepsilon_-$ and
$\varepsilon_+$ (short-range interactions) or non-concave there (long-range
interactions; red, dashed); the minimum jumps at $T_1$, with the entropy
jump $\Delta s$.
(b)~Continuous transition: the curvature of $\lds$ vanishes at the single
energy $\varepsilon_c$; the shaded deviation from the tangent is
$\lds_{\rm sing}=-B_\pm|\delta\varepsilon|^{\sigma}$ of
Assumption~\ref{ass:crit} (drawn for $\sigma=3$).
(c)~Crossover: $\lds$ changes its form smoothly about $\varepsilon_\times$
(grey), and $\omega(T)$ is analytic.
(d)~Discontinuous quantum transition: two branches (red, dashed) of
different slopes exchange at $\hfield_c$, and the ground-state energy
$\varepsilon_0(\hfield)$ (blue) has a kink for $V\to\infty$; at finite size
the two lowest eigenvalues may avoid each other (grey).
(e)~Continuous quantum transition: the \lndos{} above the lower edge for
$\hfield<\hfield_c$, $\hfield=\hfield_c$ and $\hfield>\hfield_c$.}
\label{fig:transitions}
\end{figure}

\paragraph{Discontinuous transition.} Let $\fen$ have two local minima, at
$\varepsilon_-(T)<\varepsilon_+(T)$, with the values
$\omega_\mp(T):=\fen(\varepsilon_\mp)$; then $\omega=\min(\omega_-,\omega_+)$.
Since $\partial\omega_\mp/\partial T=-\lds(\varepsilon_\mp)$ and
$\lds(\varepsilon_+)>\lds(\varepsilon_-)$, $\omega_+$ decreases faster with
$T$ than $\omega_-$. With $T_1$ defined by $\omega_-(T_1)=\omega_+(T_1)$, the
thermal energy is $\varepsilon_-$ for $T<T_1$ and $\varepsilon_+$, where the
levels are more numerous, for $T>T_1$. At $T_1$ the potential $\omega$ has a
kink\footnote{A \emph{kink} of a function is a point at which it is
continuous and its first derivative jumps.}, the entropy jumps by
$\Delta s=\lds(\varepsilon_+)-\lds(\varepsilon_-)$, and the latent heat per
unit size is $\varepsilon_+-\varepsilon_-=T_1\,\Delta s$. The singularity
lies in the selection of the energy: $\lds$ is analytic at
$\varepsilon_\mp$, and no exponent is associated with the transition. The
minimum that is not selected persists beyond $T_1$ as a local minimum (a
metastable phase), which $\omega$ alone does not show. Between
$\varepsilon_-$ and $\varepsilon_+$ two forms of $\lds$ occur.
(i)~For long-range interactions $\lds$ may be non-concave there. Then
$\fen(\varepsilon)>\omega(T)$ on this interval at every temperature: no
temperature selects these levels.
(ii)~For short-range interactions $\lds$ is concave, and it is linear on
$[\varepsilon_-,\varepsilon_+]$ with the slope $1/T_1$. The levels in this
interval are the phase-separated states: a fraction $x$ of the sample has
the energy density $\varepsilon_+$ and the fraction $1-x$ has
$\varepsilon_-$, the interface between them costing an energy that is not
extensive. The energies of the two parts add,
$\varepsilon=(1-x)\varepsilon_-+x\,\varepsilon_+$, and their numbers of
levels multiply, $\lds=(1-x)\,\lds(\varepsilon_-)+x\,\lds(\varepsilon_+)$,
which is linear in $\varepsilon$. At $T_1$, and only at $T_1$, every energy
of the interval is a minimum of $\fen$.

A transition that is discontinuous in a field $\hfield$ at $\hfield=0$
(spontaneous order) has the same structure. Let the levels at the thermal
energy consist of two groups with the conjugate densities $a=\pm m$,
degenerate at $\hfield=0$. By Eq.~\eqref{eq:filt_af} a field changes the
potential of each group by $\mp m\,\hfield$ to first order, so that
$\omega(T,\hfield)=\omega(T,0)-m|\hfield|+O(\hfield^{2})$: the conjugate density
$-\partial\omega/\partial \hfield$ jumps from $-m$ to $m$.

\paragraph{Continuous transition.} The minimiser $\varepsilon_*(T)$ is
continuous at $T_c$, so that there is no latent heat, but
$c_\mu=d\varepsilon_*/dT$ is singular. By Eq.~\eqref{eq:filt_curvmain} a
divergent $c_\mu$ requires the curvature of the \lndos{} to vanish at the
energy $\varepsilon_c=\varepsilon_*(T_c)$; it vanishes at a single energy,
not on an interval as in case~(ii) above. The singularity is a property of
$\lds$ at $\varepsilon_c$, and it carries the exponents and the scaling
functions of Section~\ref{sec:filt_critical}.

\paragraph{Crossover.} If $\lds$ has two different forms below and above an
energy $\varepsilon_\times$ (two power laws, or an activated form below a
gap and a power law above), the observables change their behaviour when
$\varepsilon_*(T)$ passes $\varepsilon_\times$, at the temperature $T_\times$
defined by
$\partial \lds/\partial\varepsilon|_{\varepsilon_\times}=1/T_\times$. If $\lds$
is analytic at $\varepsilon_\times$ with a finite curvature, $\omega(T)$ is
analytic at $T_\times$, and the change extends over the range of
temperatures in which $\varepsilon_*(T)$ traverses the range of energies
over which the form of $\lds$ changes. Examples are the Debye crossover of
the heat capacity of a crystal from $c\propto T^{3}$ to the Dulong--Petit
value, the maximum of the heat capacity of a narrow group of levels
(Section~\ref{sec:frustration}), and the rounding of a transition in a
finite system (Section~\ref{sec:filt_size}).

A crossover does not correspond to a non-analytic point of $\lds$. For
$V\to\infty$ the distribution of the energy per unit size has the width
$(k_BT^{2}c_\mu/V)^{1/2}\to0$ (Section~\ref{sec:large}): the Boltzmann factor
selects the single energy $\varepsilon_*(T)$, and a non-analytic point of
$\lds$ at $\varepsilon_\times$ makes $\omega$ non-analytic at $T_\times$,
which is a transition, case~(b). The sharp change of the spectrum that
underlies a crossover lies in one of two other places.
\begin{enumerate}[label=(\roman*),itemsep=0pt]
\item \emph{In the spectrum of the single excitations}, i.e.\ of the
  excitations of which the many-body levels are sums, as in
  Eq.~\eqref{eq:lif_levels}; their energies are of order $k_BT$, not of
  order $V$. Examples of a sharp point of this spectrum, at an energy
  $E_\times$, are a gap, a band edge and the upper edge of the phonon
  spectrum. The thermal occupation of an
  excitation varies with its energy on the scale $k_BT$
  (Section~\ref{sec:filt_dist}), so that the point enters $\omega(T)$
  through an integral over an interval of width $k_BT$, as in
  Eq.~\eqref{eq:lif_omega}. The integral is analytic in $T$, and so is $\lds$;
  the observables change over the range $k_BT\sim E_\times$. A measurement
  with a resolution finer than $k_BT$ locates the point: the spectroscopy of
  the single excitations, or the thermodynamics for $T\to0$ as a function of
  a parameter $\hfield$ on which $E_\times$ depends. If $E_\times(\hfield)$
  vanishes at some $\hfield_c$, the crossover line
  $k_BT_\times\sim E_\times(\hfield)$ ends at $T=0$ in a quantum transition,
  discontinuous or continuous as described below; the Lifshitz transition of
  Section~\ref{sec:filt_scaling0} is an example.
\item \emph{In the spectrum at a neighbouring value of a parameter.} A
  transition that exists for $V=\infty$ or at $\hfield=0$ appears at a finite
  size or in a non-zero field as a crossover, which extends over
  $|t|\lesssim V^{-1/\nu_*}$ (Eq.~\eqref{eq:filt_fss}) or
  $|t|\lesssim|\hfield|^{1/\Delta_\hfield}$ (Eq.~\eqref{eq:filt_crit}), with
  $t$ the relative distance of the temperature from the transition. In the
  limit $V\to\infty$ or $\hfield\to0$ the range shrinks to a point, and the
  crossover becomes the discontinuous or the continuous transition.
\end{enumerate}

\paragraph{Quantum transition.} At $T=0$ the thermal energy is the
ground-state energy per unit size $\varepsilon_0(\hfield)$, the lower edge
of the spectrum, and a transition as a function of a parameter $\hfield$ is
a non-analyticity of $\varepsilon_0(\hfield)$ at some $\hfield_c$. At a finite
size the lowest eigenvalue $\enrg_0(\hfield)$ is analytic unless two levels
cross, which requires that the two eigenstates differ in a conserved
quantum number\footnote{Two eigenvalues of $\hat{\mathcal H}(\hfield)$ that
approach each other as $\hfield$ varies cross only if the corresponding
eigenstates differ in a conserved quantum number. Otherwise they repel
(an \emph{avoided crossing}): they approach to a minimum distance, the gap,
and the two eigenstates exchange their character.}.
(i)~\emph{Discontinuous.} Two branches of eigenvalues whose slopes
$\partial\enrg/\partial\hfield$ differ by an amount of order $V$ exchange at
$\hfield_c$, by a crossing or by an avoided crossing whose gap vanishes for
$V\to\infty$. Then $\varepsilon_0(\hfield)$ has a kink, and the conjugate
density $a=-\partial\varepsilon_0/\partial\hfield$ jumps
(Fig.~\ref{fig:transitions}d).
(ii)~\emph{Continuous.} The density $a$ is continuous at $\hfield_c$ and a
higher derivative of $\varepsilon_0$ is singular. The gap above the ground
state closes at $\hfield_c$ for $V\to\infty$, and the form of the \lndos{}
above the lower edge changes there (Fig.~\ref{fig:transitions}e). If this
form is scale-free, the thermodynamics at $T>0$ near $\hfield_c$ is that of
Section~\ref{sec:filt_scaling0}.

\section{Power laws and exponents from the spectrum}
\label{sec:filt_lowenergy}

Eq.~\eqref{eq:filt_legendre} turns an assumption on the form of
$\lds$, i.e.\ on the many-body level count on the
logarithmic scale, into a statement on the observables. In each case the
path is: \lndos, thermodynamic potential, then entropy, heat capacity and
responses by differentiation (see Table~\ref{tab:tools}).

\subsection{Low-temperature power laws}
\label{sec:filt_power_mb}

\begin{assumption}[Power law above the ground state]\label{ass:power}
With $\delta\varepsilon:=\varepsilon-\varepsilon_0$ the excitation energy per
unit size, for $\delta\varepsilon\to0^{+}$,
\begin{equation}\label{eq:filt_stretched}
  \lds(\varepsilon)=\cA\,(\delta\varepsilon)^{\sigma} ,
  \quad \cA>0 ,\quad 0<\sigma<1 ,
  \qquad\text{i.e.}\qquad
  \cN(\enrg)\asymp\exp\left[\frac{V\cA}{k_B}\Big(\frac{\enrg-\enrg_0}{V}\Big)^{\sigma}\right]
\end{equation}
to leading order in $\delta\varepsilon$: the level count grows above the ground state
as a stretched exponential of the excitation energy.\footnote{Precisely,
$\lds(\varepsilon)=\cA\,(\delta\varepsilon)^{\sigma}\,[1+o(1)]$. Then for
every $\eta>0$,
$\cA(1-\eta)(\delta\varepsilon)^{\sigma}\le\lds\le\cA(1+\eta)(\delta\varepsilon)^{\sigma}$
at sufficiently small $\delta\varepsilon$. Since
$\min_\varepsilon[\varepsilon-T\lds(\varepsilon)]$ decreases when $\lds$ is
increased, $\omega$ lies between the values obtained with the amplitudes
$\cA(1\pm\eta)$, and the results of Proposition~\ref{prop:power} hold to
leading order for $T\to0$.}
\end{assumption}
The amplitude and the exponent are positive because $\cN$ is non-decreasing,
and $\lds\to0$ for $\delta\varepsilon\to0$ excludes a residual entropy. The
condition $\sigma<1$ is required for the thermal excitation energy
$\delta\varepsilon_*:=\varepsilon_*-\varepsilon_0$ to be positive at every
$T>0$ and to vanish for $T\to0$: for $\sigma\ge1$ the minimum of
$\fen=\varepsilon_0+\delta\varepsilon-T\cA\,(\delta\varepsilon)^{\sigma}$ at
small $\delta\varepsilon$ and small $T$ stays at $\delta\varepsilon=0$. For
short-range interactions $\sigma>1$ is excluded by the concavity of $\lds$.

\begin{proposition}[Low-temperature power laws]\label{prop:power}
Under Assumption~\ref{ass:power}, to leading order for $T\to0$,
\begin{equation}\label{eq:filt_lowT}
\begin{gathered}
  \delta\varepsilon_*(T) = (\sigma\cA\,T)^{1/(1-\sigma)} ,\qquad
  \omega-\varepsilon_0 = -\frac{1-\sigma}{\sigma}\,\delta\varepsilon_* ,\qquad
  u-\mu n-\varepsilon_0=\delta\varepsilon_* ,
  \\
  \frac{S}{V} = \cA\,(\delta\varepsilon_*)^{\sigma}=\frac{\delta\varepsilon_*}{\sigma T} ,\qquad
  c_\mu = \frac{\delta\varepsilon_*}{(1-\sigma)\,T} .
\end{gathered}
\end{equation}
All are powers of $T$:
\begin{equation}\label{eq:filt_powerT}
  c_\mu \propto T^{\kappa} , \qquad
  \kappa = \frac{\sigma}{1-\sigma}
  \;\Longleftrightarrow\;
  \sigma = \frac{\kappa}{\kappa+1} ,
  \qquad
  \frac{S}{C_\mu} = \frac{1}{\kappa} .
\end{equation}
\end{proposition}
\begin{proof}
The function
$\fen(\varepsilon)=\varepsilon_0+\delta\varepsilon-T\cA\,(\delta\varepsilon)^{\sigma}$
is that of Lemma~\ref{lem:balance}, Eq.~\eqref{eq:tools_edge}: it has its
minimum where $\sigma\cA\,(\delta\varepsilon_*)^{\sigma-1}=1/T$, which is
Eq.~\eqref{eq:filt_stat}; there
$T\cA\,(\delta\varepsilon_*)^{\sigma}=\delta\varepsilon_*/\sigma$, which gives
$\omega=\fen(\varepsilon_*)$ and $S/V=\lds(\varepsilon_*)$. Since
$\delta\varepsilon_*\to0$ for $T\to0$, this minimum lies in the range
$\delta\varepsilon<\varepsilon_1$ in which the power law holds. It is the
global minimum: for $\delta\varepsilon\ge\varepsilon_1$ the
bound~\eqref{eq:filt_linbound} gives
$\fen-\varepsilon_0\ge\varepsilon_1(1-T/\Tmax)-(c_0+\varepsilon_0/\Tmax)\,T$,
which is positive at small $T$, whereas
$\fen(\varepsilon_*)-\varepsilon_0<0$. Finally $u-\mu n$ and $c_\mu$ follow
from Eq.~\eqref{eq:filt_SI}, with
$d(\delta\varepsilon_*)/dT=\delta\varepsilon_*/(1-\sigma)T$.
\end{proof}
A stretched-exponential level count with exponent $\sigma$ is thus equivalent
to a power-law heat capacity with exponent $\kappa=\sigma/(1-\sigma)$, for
any Hamiltonian that satisfies the assumption, interacting or not.

\paragraph{Inverse relation.} Conversely, let the heat capacity be
$ c_\mu = \mathcal C_0\,T^{\kappa} $ at low temperature, with $S(T=0)=0$.
Then
$S/V = \int_0^{T}c_\mu\,dT'/T' = \mathcal C_0T^{\kappa}/\kappa$ and
$u-\mu n-\varepsilon_0 = \int_0^{T}c_\mu\,dT' = \mathcal C_0T^{\kappa+1}/(\kappa+1)$.
By Eq.~\eqref{eq:filt_SI} these are $\lds(\varepsilon_*)$ and
$\delta\varepsilon_*$; eliminating $T$,
\begin{equation}\label{eq:filt_sinv}
  \lds(\varepsilon) = \frac{\mathcal C_0}{\kappa}
  \left(\frac{(\kappa+1)\,\delta\varepsilon}{\mathcal C_0}\right)^{\kappa/(\kappa+1)} ,
\end{equation}
which is Eq.~\eqref{eq:filt_stretched}. The parametric relation
$\lds(u-\mu n)=S/V$ holds at every temperature and returns the concave hull of
$\lds$: the \lndos{} is measured by calorimetry. In the mathematical
literature the equivalence of $\ln\cN\simeq B(\enrg-\enrg_0)^{\sigma}$ and
$\ln\Xi+\beta\enrg_0\propto T^{\kappa}$ is Kohlbecker's Tauberian
theorem\footnote{A \emph{Tauberian theorem} infers the asymptotic form of a
function from that of its Laplace transform.}~\cite[Sec.~4.12]{BGT}.

\paragraph{Test of the assumption.} By Eq.~\eqref{eq:filt_lowT} the ratios
\[
  \frac{S}{C_\mu}=\frac1\kappa ,\qquad
  \frac{u-\mu n-\varepsilon_0}{T\,c_\mu}=\frac{1}{\kappa+1} ,\qquad
  \frac{\varepsilon_0-\omega}{T\,c_\mu}=\frac{1}{\kappa(\kappa+1)}
\]
do not depend on the amplitude $\cA$ and give $\kappa$ at each temperature
without a fit. Their constancy in $T$ and their mutual consistency test
Assumption~\ref{ass:power}; a drift signals a correction to the power law,
e.g.\ a logarithm (Section~\ref{sec:filt_nonpower}).

\paragraph{Spectra that violate the assumption.}
\emph{(i)~Power-law level count.} If the level count itself is a power,
$\cN(\enrg)\propto(\enrg-\enrg_0)^{p}$, Eq.~\eqref{eq:filt_laplace} gives
$e^{\beta\enrg_0}\Xi\propto\Gamma(p+1)\,(k_BT)^{p}$ and $C_\mu=p\,k_B$,
independent of $T$: equipartition (the classical ideal gas of $N$ particles
has $\cN\propto E^{3N/2}$). A heat capacity that is extensive and vanishes
as a power of $T$ requires a level count that grows faster than any power,
as in Eq.~\eqref{eq:filt_stretched}.
\emph{(ii)~Gap.} Let the lowest level, of degeneracy $d_0$, be separated by
a gap $\Delta$ from the next one, of degeneracy $d_1$. For $k_BT$ small
compared with the distance to the following level,
$e^{\beta\enrg_0}\Xi=d_0(1+q)$ with $q:=(d_1/d_0)\,e^{-\beta\Delta}$, hence
$\Var(\hat{\mathcal H})=\Delta^{2}q/(1+q)^{2}$ and, to leading order in $q$,
\begin{equation}\label{eq:filt_gap}
  C_\mu \;\simeq\; k_B\left(\frac{\Delta}{k_BT}\right)^{2}\frac{d_1}{d_0}\,e^{-\Delta/k_BT} :
\end{equation}
the heat capacity is activated, with the many-body gap of
$ \hat{\mathcal{H}} $. For $n_1V$ independent local excitations of energy
$\Delta$, $d_1/d_0=n_1V$ and Eq.~\eqref{eq:filt_gap} is extensive. The
corresponding \lndos{} is
$\lds-\lds_0=k_Bn_1[-x\ln x-(1-x)\ln(1-x)]$ with
$x=\delta\varepsilon/n_1\Delta$, which rises from $\delta\varepsilon=0$ with a
logarithmically divergent slope (Fig.~\ref{fig:residual}a).
\emph{(iii)~Residual entropy.} If the number of ground states is exponential
in the size, $d_0=e^{V\lds_0/k_B}$, then $\lds\to\lds_0>0$ for
$\delta\varepsilon\to0$ and $S\to V\lds_0$ for $T\to0$
(Section~\ref{sec:frustration}).
\emph{(iv)~Non-extensive spectra.} Extensivity is not essential for the
exponent. If $\ln\cN(\enrg)=B\,(\enrg-\enrg_0)^{\sigma}[1+o(1)]$ in a range of
$\enrg$, with any dependence of $B$ on the size, the steepest-descent
evaluation of Eq.~\eqref{eq:filt_laplace} gives directly
\begin{equation}\label{eq:filt_total}
  \ln\Xi+\beta\enrg_0=\frac{1-\sigma}{\sigma}\,(\sigma B)^{1/(1-\sigma)}\,(k_BT)^{\kappa} ,
  \qquad
  C_\mu=k_B\,(\kappa+1)\,(\sigma B)^{1/(1-\sigma)}(k_BT)^{\kappa} ,
\end{equation}
for the temperatures at which
$\enrg_*-\enrg_0=(\sigma Bk_BT)^{1/(1-\sigma)}$ lies in that range and
$\ln\Xi+\beta\enrg_0\gg1$. The extensive case is $B=\cA V^{1-\sigma}/k_B$.

\paragraph{Examples: ideal quantum gases.} For non-interacting particles
both sides are known independently (Table~\ref{tab:gases}). With
$\varrho(E)$ the single-particle density of states per unit volume and
$x:=(E-\mu)/2k_BT$, the heat capacities of independent fermions and bosons
are
\begin{equation}\label{eq:filt_cF}
  c_\mu=k_B\int dE\;\varrho(E)\,\frac{x^{2}}{\cosh^{2}x}
  \quad\text{(fermions)},
  \qquad
  c_\mu=k_B\int dE\;\varrho(E)\,\frac{x^{2}}{\sinh^{2}x}
  \quad\text{(bosons)},
\end{equation}
and for a power law $\varrho=B|E-\mu|^{r}$ they give
$c_\mu=2k_BB\,\Gamma(r+3)\,\etaD(r+2)\,(k_BT)^{r+1}$ for fermions and, with
$\varrho=B(E-\mu)^{r}$ for $E>\mu$,
$c_\mu=k_BB\,\Gamma(r+3)\,\zeta(r+2)\,(k_BT)^{r+1}$ for bosons, where
$\etaD(s)=(1-2^{1-s})\zeta(s)$. In each case $\lds$ follows from
Eq.~\eqref{eq:filt_sinv}, and $\kappa-1=r$ is the exponent of the
single-particle density of states.
\begin{table}[htbp]
\centering
\small
\caption{Low-temperature heat capacity $c_\mu\propto T^{\kappa}$ of ideal
quantum gases and the exponent $\sigma=\kappa/(\kappa+1)$ of the \lndos{}
$\lds\propto(\delta\varepsilon)^{\sigma}$, Eq.~\eqref{eq:filt_powerT};
$\varrho$ is the single-particle density of states per unit volume and $v$
the velocity of the mode (of light, for photons).}
\label{tab:gases}
\begin{tabular}{llccl}
\toprule
System & heat capacity (low $T$) & $\kappa$ & $\sigma$ & Ref.~\cite{Schwabl} \\
\midrule
degenerate Fermi gas & $c=\tfrac{\pi^{2}}{3}k_B^{2}\varrho(\mu)\,T$ & $1$ & $\tfrac12$ & Eq.~(4.3.24a) \\[2pt]
fermions, $\varrho=B|E-\mu|^{r}$ & $c=2k_BB\,\Gamma(r+3)\,\etaD(r+2)\,(k_BT)^{r+1}$ & $r+1$ & $\tfrac{r+1}{r+2}$ & --- \\[2pt]
Bose gas, $T<T_c$, $d=3$ & $c=\tfrac{15}{4}\tfrac{\zeta(5/2)}{\zeta(3/2)}\,nk_B\,(T/T_c)^{3/2}$ & $\tfrac32$ & $\tfrac35$ & Eq.~(4.4.11) \\[2pt]
photons & $c=\tfrac{4\pi^{2}}{15}\,k_B\,(k_BT/\hbar v)^{3}$ & $3$ & $\tfrac34$ & Eq.~(4.5.7) \\[2pt]
acoustic phonons & $c=\tfrac{2\pi^{2}}{5}\,k_B\,(k_BT/\hbar v)^{3}$ & $3$ & $\tfrac34$ & Eq.~(4.6.13) \\
\bottomrule
\end{tabular}
\end{table}
\begin{itemize}\setlength{\itemsep}{0pt}
\item \emph{Degenerate Fermi gas.} The Sommerfeld law
  \begin{equation}\label{eq:filt_sommerfeld}
    c_\mu=\frac{\pi^{2}}{3}\,k_B^{2}\,\varrho(\mu)\,T
  \end{equation}
  holds whenever $\varrho$ is smooth on the scale $k_BT$ about $\mu$. With
  $\kappa=1$, Eq.~\eqref{eq:filt_sinv} is
  $\lds = \sqrt{2\,\mathcal C_0\,\delta\varepsilon}
  = 2k_B\sqrt{\pi^{2}\varrho(\mu)\,\delta\varepsilon/6}$,
  which is Bethe's formula for the level density of a Fermi
  gas~\cite{Bethe1936}, $ \rho \propto e^{2\sqrt{\tilde a(\enrg-\enrg_0)}} $
  with $ \tilde a = \pi^{2}\varrho(\mu)V/6 $, obtained there by counting the
  levels. The same form, $\ln\rho\propto\sqrt{\enrg-\enrg_0}$, holds for
  every one-dimensional system whose low-energy spectrum is conformally
  invariant (Cardy's formula~\cite{Cardy1986}).
\item \emph{Fermions at a zero of $ \varrho $.} At a Dirac point,
  $E=\pm\hbar v|\mathbf k|$ in $d$ dimensions, $\varrho\propto|E|^{d-1}$, so
  that $r=d-1$ and $c\sim T^{d}$.
\item \emph{Ideal Bose gas below the condensation temperature.} Here
  $ \mu=0 $ and $ \varrho(E)=BE^{1/2} $; the density
  $n=B\,\Gamma(\tfrac32)\zeta(\tfrac32)(k_BT_c)^{3/2}$ eliminates $B$. In $d$
  dimensions $\kappa=d/2$, for $d>2$ only: for $d\le2$ there is no
  condensation at a fixed density~\cite[Sec.~6.4]{Honerkamp}.
\item \emph{Photons and phonons.} $\varrho(E)\propto E^{2}$ and $\mu=0$,
  with two polarisations for photons and three branches for phonons
  (Debye's law).
\end{itemize}

\subsection{Scaling above the ground state and quantum criticality}
\label{sec:filt_scaling0}

Let the Hamiltonian depend on a parameter $ \hfield $, measured from a
reference value $\hfield=0$, and let
$\delta\varepsilon:=\varepsilon-\varepsilon_0-e_{\rm reg}(\hfield)$, where
$\varepsilon_0$ is the ground-state energy at $\hfield=0$ and
$e_{\rm reg}(\hfield)$, with $e_{\rm reg}(0)=0$, is an analytic displacement
of the lower edge of the spectrum by the parameter ($e_{\rm reg}=0$ if
there is none).

\begin{assumption}[Homogeneity above the ground state]\label{ass:hom0}
For small $\delta\varepsilon$ and $\hfield$ the \lndos{} has no intrinsic
scale: for all $\lambda>0$,
\begin{equation}\label{eq:filt_shom0}
  \lds(\lambda\,\delta\varepsilon,\lambda^{b}\hfield)=\lambda^{\sigma}\lds(\delta\varepsilon,\hfield) ,
  \qquad 0<\sigma<1 ,\quad b>0 ,
\end{equation}
i.e.\ $\ln\cN$ is a generalised homogeneous function\footnote{A function
$F(x_1,\dots,x_n)$ is \emph{generalised homogeneous} if
$F(\lambda^{a_1}x_1,\dots,\lambda^{a_n}x_n)=\lambda^{a_F}F(x_1,\dots,x_n)$ for
all $\lambda>0$, with fixed exponents $a_i$ and $a_F$. The choice
$\lambda=|x_1|^{-1/a_1}$ shows that $F$ is a power of $|x_1|$ times a
function of the ratios $x_i/|x_1|^{a_i/a_1}$.} of the excitation energy and
the parameter.
\end{assumption}
At $ \hfield=0 $ this is Assumption~\ref{ass:power} with $ \cA = \lds(1,0) $.
With $\lambda=|\hfield|^{-1/b}$, the parameter enters $\lds$ only through the
ratio $\delta\varepsilon/|\hfield|^{1/b}$: it changes the level count in the
range of energies $\delta\varepsilon\lesssim|\hfield|^{1/b}$ above the ground
state.

\begin{proposition}[Scaling form]\label{prop:scaling0}
Under Assumption~\ref{ass:hom0}, with
$\delta\omega:=\omega-\varepsilon_0-e_{\rm reg}(\hfield)$,
\begin{equation}\label{eq:filt_scalingform}
\begin{gathered}
  \delta\omega(\lambda^{1-\sigma}T,\lambda^{b}\hfield)=\lambda\,\delta\omega(T,\hfield)
  \quad\text{for all }\lambda>0 ,\qquad\text{hence}
  \\
  \delta\omega(T,\hfield)=T^{1/(1-\sigma)}\,\Phi\left(\frac{\hfield}{T^{\,b/(1-\sigma)}}\right),
\end{gathered}
\end{equation}
with $\Phi(w):=\delta\omega(1,w)$.
\end{proposition}
\begin{proof}
In $\delta\omega(\lambda^{1-\sigma}T,\lambda^{b}\hfield)
=\min_{\delta\varepsilon}[\delta\varepsilon-\lambda^{1-\sigma}T\,\lds(\delta\varepsilon,\lambda^{b}\hfield)]$
substitute $\delta\varepsilon=\lambda\,\delta\varepsilon'$ and use
Eq.~\eqref{eq:filt_shom0} (Lemma~\ref{lem:rescale}): the minimised function
becomes $\lambda[\delta\varepsilon'-T\,\lds(\delta\varepsilon',\hfield)]$, whose
minimum is $\lambda\,\delta\omega(T,\hfield)$. The choice
$\lambda=T^{-1/(1-\sigma)}$ makes the first argument equal to $1$.
\end{proof}
By differentiation, with $ w:=\hfield/T^{\,b/(1-\sigma)} $ and
$\kappa=\sigma/(1-\sigma)$,
\[
  \frac{S}{V}=T^{\kappa}\,\Phi_S(w) ,\qquad
  c_\mu=T^{\kappa}\,\Phi_c(w) ,\qquad
  a+\frac{de_{\rm reg}}{d\hfield}=-T^{(1-b)/(1-\sigma)}\,\Phi'(w) ,
\]
where $ \Phi_S $ and $ \Phi_c $ are combinations of $ \Phi $ and its
derivatives. No property of the Hamiltonian is used beyond
Eq.~\eqref{eq:filt_shom0}: the scaling form of the thermodynamic potential
is the homogeneity of the level count, carried through the Laplace
transform.

\paragraph{Quantum critical point.} At a continuous quantum transition
(Section~\ref{sec:transitions}), let Assumption~\ref{ass:hom0} hold about
the ground state at $\hfield_c$, with $\hfield$ measured from $\hfield_c$.
The observables depend on $\hfield$ only through $w$. For $|w|\ll1$ they have
the temperature dependence of the critical point,
$c_\mu\propto T^{\sigma/(1-\sigma)}$; for $|w|\gg1$ they have that of the
phase at $\hfield\gtrless0$. The region $|\hfield|\lesssim T^{\,b/(1-\sigma)}$
is the \emph{quantum critical fan}. Unless $\Phi$ is singular at a finite
$w$, its boundaries are crossovers, and the transition is confined to $T=0$.
A finite $\delta\omega$ for $T\to0$ at $\hfield\neq0$ requires
$\Phi(w)\propto|w|^{1/b}$ for $|w|\to\infty$ (Eq.~\eqref{eq:tools_match}):
the singular part of the ground-state energy is $\propto|\hfield|^{1/b}$. The
standard scaling form $f=T^{(d+z)/z}\,\Phi(\hfield/T^{1/\nu z})$ of a quantum
critical point in $d$ dimensions, with the dynamical exponent $z$ (the
energy of the critical modes of wave vector $q$ scales as $q^{z}$) and the
correlation-length exponent $\nu$~\cite{Lohneysen2007}, is
Proposition~\ref{prop:scaling0} with
\begin{equation}\label{eq:filt_qcp}
  \sigma=\frac{d}{d+z} ,\qquad b=\frac{1}{\nu\,(d+z)} .
\end{equation}

\paragraph{Example: the Lifshitz transition.}
The simplest quantum transition without an order parameter is the passage
of the chemical potential through a band edge of independent fermions, the
\emph{Lifshitz transition}~\cite{Lifshitz1960}: no symmetry distinguishes
the two sides. Here Assumption~\ref{ass:hom0} can be verified on the
spectrum itself, and the scaling function is known in closed form.

\emph{Spectrum.} Let the single-particle density of states per unit volume
be $\varrho(E)=B\,(E-E_0)^{r}$ above the band edge $E_0$ and zero below it,
with $r>-1$ ($r=d/2-1$ for a parabolic band in $d$ dimensions;
Fig.~\ref{fig:lifshitz}a), and let
\[
  \hfield:=\mu-E_0
\]
be the chemical potential measured from the band edge. It is a parameter of
$\hat{\mathcal H}=\hat H-\mu\hat N$ in the sense of
Section~\ref{sec:filt_dist}, with $\hat A=\hat N$ and the conjugate density
$a=n$. With $x_k:=E_k-E_0\ge0$ the single-particle energies measured from
the edge and $n_k\in\{0,1\}$ their occupations, the many-body levels are
\begin{equation}\label{eq:lif_levels}
  \enrg=\sum_k(x_k-\hfield)\,n_k .
\end{equation}
For $\hfield<0$ every term is positive: the ground state is the empty band,
$\enrg_0=0$, and the lowest excitation, one particle at the band edge, costs
the gap $|\hfield|$. For $\hfield>0$ the ground state has all states with
$x_k<\hfield$ occupied. The ground-state energy per unit size is therefore
\begin{equation}\label{eq:lif_ground}
  \varepsilon_0(\hfield)=-B\int_0^{\hfield}dx\,x^{r}(\hfield-x)
  =-\frac{B\,\hfield^{\,r+2}}{(r+1)(r+2)}\quad(\hfield>0) ,
  \qquad
  \varepsilon_0(\hfield)=0\quad(\hfield<0) .
\end{equation}
The lower edge of the spectrum is non-analytic at $\hfield=0$, and the
density $n=-\partial\varepsilon_0/\partial\hfield=B\,\hfield^{\,r+1}/(r+1)$
is continuous there: a continuous quantum transition in the sense of
Section~\ref{sec:transitions}.

\emph{Homogeneity of the level count.} Consider the system of size $V$ at
the parameter $\Lambda\hfield$, with $\Lambda>0$, and write its
single-particle energies as $x_k=\Lambda x'_k$. The number of the $x'_k$ in
an interval $dx'$ is $VB\,(\Lambda x')^{r}\,\Lambda\,dx'
=\Lambda^{r+1}V\cdot B\,x'^{\,r}dx'$: the $x'_k$ are the single-particle
energies of a system of size $\Lambda^{r+1}V$. By
Eq.~\eqref{eq:lif_levels}, $\enrg=\Lambda\sum_k(x'_k-\hfield)\,n_k$. The
many-body levels of the system of size $V$ at the parameter $\Lambda\hfield$
are thus $\Lambda$ times those of the system of size $\Lambda^{r+1}V$ at the
parameter $\hfield$, and the level counts are equal,
\[
  \cN_V(\Lambda\enrg)\Big|_{\Lambda\hfield}
  =\cN_{\Lambda^{r+1}V}(\enrg)\Big|_{\hfield} .
\]
Taking $(k_B/V)\ln$ of both sides for $V\to\infty$, with
$\varepsilon=\enrg/\Lambda^{r+1}V$ the energy per unit size of the larger
system, so that $\Lambda\enrg/V=\Lambda^{r+2}\varepsilon$,
\begin{equation}\label{eq:lif_hom}
  \lds(\Lambda^{r+2}\varepsilon,\Lambda\hfield)=\Lambda^{r+1}\,\lds(\varepsilon,\hfield) .
\end{equation}
With $\lambda:=\Lambda^{r+2}$ this is Eq.~\eqref{eq:filt_shom0}, with
$\delta\varepsilon=\varepsilon$, $e_{\rm reg}=0$ and
\begin{equation}\label{eq:lif_exp}
  \sigma=\frac{r+1}{r+2} ,\qquad b=\frac{1}{r+2} .
\end{equation}
Assumption~\ref{ass:hom0} holds exactly in the range of energies in which
$\varrho$ is a power law. It follows from the power law of the
single-particle density of states alone, without reference to a
thermodynamic potential. For $r=d/2-1$, $\sigma=d/(d+2)$ and $b=2/(d+2)$,
which is Eq.~\eqref{eq:filt_qcp} with $z=2$ and $\nu=\tfrac12$.

\emph{Scaling function.} Proposition~\ref{prop:scaling0} with
Eq.~\eqref{eq:lif_exp} gives
$\omega(T,\hfield)=T^{\,r+2}\,\Phi(\hfield/T)$. This is confirmed by the
direct evaluation. For independent fermions~\cite[Ch.~4]{Schwabl}
\[
  \omega=-k_BT\int dE\,\varrho(E)\,\ln\big(1+e^{-(E-\mu)/k_BT}\big) ,
\]
and the substitution $E-E_0=k_BT\,x$ followed by a partial integration gives
\begin{equation}\label{eq:lif_omega}
  \omega(T,\hfield)=B\,(k_BT)^{r+2}\,\Phi_r\Big(\frac{\hfield}{k_BT}\Big) ,
  \qquad
  \Phi_r(y):=-\frac{1}{r+1}\int_0^{\infty}dx\,\frac{x^{r+1}}{e^{x-y}+1} .
\end{equation}
The three limits of the scaling function are the three regimes of
Fig.~\ref{fig:lifshitz}b.
\begin{itemize}\setlength{\itemsep}{0pt}
\item \emph{Empty band}, $\hfield<0$ and $k_BT\ll|\hfield|$ ($y\to-\infty$):
  $\Phi_r\simeq-\Gamma(r+1)\,e^{y}$, so that
  $\omega\simeq-B\,\Gamma(r+1)\,(k_BT)^{r+2}e^{-|\hfield|/k_BT}$. All
  quantities are activated, with the gap $|\hfield|$.
\item \emph{Critical regime}, $|\hfield|\ll k_BT$ ($y\to0$):
  $\Phi_r(0)=-\Gamma(r+1)\,\etaD(r+2)$, and
  $c_\mu=-T\,\partial^{2}\omega/\partial T^{2}
  =k_BB\,\Gamma(r+3)\,\etaD(r+2)\,(k_BT)^{r+1}$: a power law with
  $\kappa=r+1=\sigma/(1-\sigma)$, as in Proposition~\ref{prop:power} (one
  half of the value of Table~\ref{tab:gases}, where $\varrho$ extends to
  both sides of $\mu$).
\item \emph{Fermi liquid}, $\hfield>0$ and $k_BT\ll\hfield$ ($y\to+\infty$):
  by the Sommerfeld expansion\footnote{For a smooth function $g$ and
  $y\gg1$,
  $\int_0^{\infty}dx\,g(x)/(e^{x-y}+1)=\int_0^{y}dx\,g(x)+\tfrac{\pi^{2}}{6}\,g'(y)+\dots$~\cite[Sec.~4.3]{Schwabl}.},
  $\Phi_r\simeq-y^{r+2}/(r+1)(r+2)-\tfrac{\pi^{2}}{6}\,y^{r}$, so that
  $\omega\simeq\varepsilon_0(\hfield)-\tfrac{\pi^{2}}{6}B\,\hfield^{\,r}(k_BT)^{2}$.
  The first term is the ground-state energy~\eqref{eq:lif_ground},
  $\propto\hfield^{1/b}$ as required by Eq.~\eqref{eq:tools_match}; the
  second gives the Sommerfeld law~\eqref{eq:filt_sommerfeld},
  $c_\mu=\tfrac{\pi^{2}}{3}k_B^{2}\varrho(\mu)\,T$ with
  $\varrho(\mu)=B\,\hfield^{\,r}$.
\end{itemize}
Since $\Phi_r(y)$ is analytic at every finite $y$, the boundaries
$k_BT\sim|\hfield|$ of the critical regime are crossovers, and the
transition is confined to $T=0$. The density is
$n=-\partial\omega/\partial\hfield=B\,(k_BT)^{r+1}\int_0^{\infty}dx\,x^{r}/(e^{x-y}+1)$,
with the prefactor $T^{(1-b)/(1-\sigma)}=T^{\,r+1}$ of
Proposition~\ref{prop:scaling0}. No symmetry and no order parameter has been
used: the exponents~\eqref{eq:lif_exp} and the three regimes follow from one
property of the spectrum, the power law of the single-particle density of
states at the band edge, which makes the many-body level count homogeneous.

\begin{figure}[htbp]
\centering
\tikzset{
  ldsgraph/.style={>=Stealth,font=\small,x=1.4cm,y=1.6cm},
  lds/.style={very thick,blue!60!black},
  guide/.style={thin,dotted,gray!80!black},
}
\begin{tikzpicture}[ldsgraph]
  % panel (a): single-particle density of states at the band edge
  \node[font=\normalsize] at (-1.45,1.55) {(a)};
  \fill[orange!25] (0,0) -- plot[domain=0:1.2,samples=40,smooth] (\x,{0.9*sqrt(\x)}) -- (1.2,0) -- cycle;
  \draw[->] (-1.1,0) -- (2.5,0) node[below left]{$E$};
  \draw[->] (-1.0,-0.1) -- (-1.0,1.5) node[above]{$\varrho$};
  \draw[lds] (-1.0,0) -- (0,0);
  \draw[lds,domain=0:2.3,samples=60,smooth] plot (\x,{0.9*sqrt(\x)});
  \draw[guide] (-0.5,0) -- (-0.5,1.15);
  \draw[guide] (1.2,0) -- (1.2,1.15);
  \node[below,font=\scriptsize] at (0,0) {$E_0$};
  \node[below,font=\scriptsize] at (-0.5,0) {$\mu$};
  \node[below,font=\scriptsize] at (1.2,0) {$\mu$};
  \node[above,font=\scriptsize] at (-0.5,1.15) {$\hfield<0$};
  \node[above,font=\scriptsize] at (1.2,1.15) {$\hfield>0$};
  \node[font=\scriptsize,blue!60!black,anchor=east] at (2.45,0.85) {$B\,(E-E_0)^{r}$};
\end{tikzpicture}\hfill
\begin{tikzpicture}[ldsgraph]
  % panel (b): regimes in the plane of the parameter and the temperature
  \node[font=\normalsize] at (-2.55,1.55) {(b)};
  \fill[orange!20] (0,0) -- (1.35,1.35) -- (-1.35,1.35) -- cycle;
  \draw[->] (-2.3,0) -- (2.4,0) node[below left]{$\hfield$};
  \draw[->] (0,-0.1) -- (0,1.5) node[above]{$k_BT$};
  \draw[dashed,thick,gray!60!black] (0,0) -- (1.35,1.35);
  \draw[dashed,thick,gray!60!black] (0,0) -- (-1.35,1.35);
  \fill (0,0) circle (1.6pt);
  \node[below,font=\scriptsize] at (0,-0.02) {$\hfield=0$};
  \node[font=\scriptsize,align=center] at (-1.55,0.45) {empty band:\\ activated, gap $|\hfield|$};
  \node[font=\scriptsize,align=center] at (1.6,0.45) {Fermi liquid:\\ $c_\mu\propto\varrho(\mu)\,T$};
  \node[font=\scriptsize,align=center,fill=orange!20,inner sep=1pt] at (0,1.0) {critical:\\ $c_\mu\propto T^{\,r+1}$};
  \node[font=\scriptsize,gray!40!black,anchor=west] at (1.38,1.3) {$k_BT=\hfield$};
\end{tikzpicture}
\caption{The Lifshitz transition.
(a)~Single-particle density of states at a band edge $E_0$; the chemical
potential lies below the edge for $\hfield<0$ (empty band) and inside the
band for $\hfield>0$ (occupied states shaded).
(b)~The regimes of the scaling function $\Phi_r(\hfield/k_BT)$,
Eq.~\eqref{eq:lif_omega}: the critical regime $|\hfield|\lesssim k_BT$
(shaded) is bounded by crossovers (dashed), and the transition (dot) is at
$T=0$.}
\label{fig:lifshitz}
\end{figure}

\subsection{Critical point at a finite temperature}
\label{sec:filt_critical}

\paragraph{Definitions.} Near a critical point at $T_c$, with the reduced
temperature $t:=(T-T_c)/T_c$ and the field $\hfield$ conjugate to the order
parameter of density $m=-\partial\omega/\partial\hfield$, the \emph{critical
exponents} are defined by
\begin{equation}\label{eq:exponents}
\begin{gathered}
  c_\mu\sim|t|^{-\alpha},\qquad
  \chi\sim|t|^{-\gamma}\quad(\hfield=0),
  \\
  m\sim(-t)^{\tilde\beta}\quad(\hfield\to0^{+},\ t<0),\qquad
  m\sim|\hfield|^{1/\delta}\quad(t=0),
\end{gathered}
\end{equation}
with $\chi=\partial m/\partial \hfield$ (the order-parameter exponent is
denoted by $\tilde\beta$, since $\beta$ is the inverse temperature),
together with the correlation length $\ell_c\sim|t|^{-\nu}$ and the decay
$G(r)\sim r^{-(d-2+\eta)}$ of the order-parameter correlation function at
$t=0$ in $d$ dimensions. The \emph{scaling hypothesis}~\cite{Widom1965}
states that the singular part of the thermodynamic potential is a
generalised homogeneous function,
\begin{equation}\label{eq:scaling}
  \omega_{\rm sing}(t,\hfield)=|t|^{2-\alpha}\,\Phi_\pm\left(\frac{\hfield}{|t|^{\Delta_\hfield}}\right) ,
\end{equation}
with the gap exponent $\Delta_\hfield$ and two scaling functions $\Phi_\pm$,
for $t>0$ and $t<0$. Differentiation gives
\begin{equation}\label{eq:rushbrooke}
  \Delta_\hfield=\tilde\beta\delta=\tilde\beta+\gamma ,\qquad \alpha+2\tilde\beta+\gamma=2 ,
\end{equation}
the second being the Rushbrooke relation~\cite[Ch.~7]{Schwabl}. Here
Eq.~\eqref{eq:scaling} is derived from an assumption on the spectrum about
an energy $\varepsilon_c$ in its interior.

\begin{assumption}[Homogeneity about $\varepsilon_c$]\label{ass:crit}
With $\delta\varepsilon:=\varepsilon-\varepsilon_c$,
\begin{equation}\label{eq:filt_gcrit}
  \lds(\varepsilon,\hfield)=\lds_c+\frac{\delta\varepsilon}{T_c}
  +\lds_{\rm sing}(\delta\varepsilon,\hfield) ,
  \qquad
  \frac{1}{T_c}:=\frac{\partial \lds}{\partial\varepsilon}\bigg|_{\varepsilon_c,\,\hfield=0} ,
\end{equation}
where the singular part has no intrinsic scale: for all $\lambda>0$ and small
arguments,
\begin{equation}\label{eq:filt_shom}
  \lds_{\rm sing}(\lambda\,\delta\varepsilon,\lambda^{b}\hfield)=\lambda^{\sigma}\lds_{\rm sing}(\delta\varepsilon,\hfield) ,
  \qquad \sigma>1 ,\quad b>0 .
\end{equation}
\end{assumption}
The first two terms are the value and the slope of $\lds$ at $\varepsilon_c$.
The slope fixes the critical temperature: by Eq.~\eqref{eq:filt_stat} the
thermal energy reaches $\varepsilon_c$ at $T=T_c$. The condition $\sigma>1$
makes $\lds_{\rm sing}$ subleading to the linear term. The simplest instance
is $\lds_{\rm sing}(\delta\varepsilon,0)=-B_\pm|\delta\varepsilon|^{\sigma}$
with amplitudes $B_\pm>0$ for $\delta\varepsilon\gtrless0$. In terms of the
level count,
\[
  \cN(\enrg)\asymp e^{V\lds_c/k_B}\;e^{(\enrg-\enrg_c)/k_BT_c}\;
  \exp\Big[\frac{V}{k_B}\,\lds_{\rm sing}\Big(\frac{\enrg-\enrg_c}{V},\hfield\Big)\Big] :
\]
an exponential growth with the energy, $\partial\ln\cN/\partial\enrg=1/k_BT_c$,
times a deviation whose logarithm is extensive. For $\sigma>2$ the deviation
has no quadratic term: the curvature of the \lndos{} vanishes at
$\varepsilon_c$, whereas at a non-critical energy it is finite and negative
(Eq.~\eqref{eq:filt_curvmain}). A range of energies about $\varepsilon_c$
then has almost the same slope $1/T_c$, and a small change of $T$ moves
$\varepsilon_*$ by a large amount. Correspondingly, the distribution of the
energy at $T_c$, in which the linear terms cancel, is
$P_{T_c}(V\varepsilon)\propto\exp[-VB_\pm|\delta\varepsilon|^{\sigma}/k_B]$.
Its width in the total energy is $V^{1-1/\sigma}$, which for $\sigma>2$
exceeds the width $V^{1/2}$ of a non-critical point.

The exponent $b$ fixes the range of energies in which a field changes the
level count. With $\lambda=|\delta\varepsilon|^{-1}$,
Eq.~\eqref{eq:filt_shom} gives
$\lds_{\rm sing}(\delta\varepsilon,\hfield)
=|\delta\varepsilon|^{\sigma}\,\lds_{\rm sing}(\pm1,\hfield/|\delta\varepsilon|^{b})$:
for $|\delta\varepsilon|\gg|\hfield|^{1/b}$ the singular part has its form at
$\hfield=0$, and for $|\delta\varepsilon|\lesssim|\hfield|^{1/b}$ it is changed
by the field to order one.

\begin{proposition}[Scaling at a critical point]\label{prop:crit}
Under Assumption~\ref{ass:crit}, if $\lds_{\rm sing}$ is the leading
deviation of $\lds$ from its tangent at $\varepsilon_c$, the thermodynamic
potential is $\omega=\varepsilon_c-T\lds_c+\omega_{\rm sing}$ with
\begin{equation}\label{eq:filt_crit}
  \omega_{\rm sing}(t,\hfield)=|t|^{\sigma/(\sigma-1)}\,
  \Phi_\pm\left(\frac{\hfield}{|t|^{\,b/(\sigma-1)}}\right),
  \qquad
  2-\alpha=\frac{\sigma}{\sigma-1} ,\quad \Delta_\hfield=\frac{b}{\sigma-1} ,
\end{equation}
to leading order in $t$, and the exponents of Eq.~\eqref{eq:exponents} are
\begin{equation}\label{eq:filt_critexp}
  \alpha=\frac{\sigma-2}{\sigma-1} ,\qquad
  \tilde\beta=\frac{\sigma-b}{\sigma-1} ,\qquad
  \gamma=\frac{2b-\sigma}{\sigma-1} ,\qquad
  \delta=\frac{b}{\sigma-b} .
\end{equation}
They satisfy the scaling relations~\eqref{eq:rushbrooke} identically.
\end{proposition}
\begin{proof}
Insert Eq.~\eqref{eq:filt_gcrit} in Eq.~\eqref{eq:filt_legendre}. Since
$1-T/T_c=-t$,
\[
  \varepsilon-T\,\lds(\varepsilon,\hfield)
  =\varepsilon_c-T\,\lds_c-t\,\delta\varepsilon-T\,\lds_{\rm sing}(\delta\varepsilon,\hfield) .
\]
The first two terms are analytic in $T$. To leading order in $t$ (with $T$
replaced by $T_c$ in the factor of $\lds_{\rm sing}$) the other two define
\begin{equation}\label{eq:filt_omsing}
  \omega_{\rm sing}(t,\hfield):=\min_{\delta\varepsilon}
  \big[-t\,\delta\varepsilon-T_c\,\lds_{\rm sing}(\delta\varepsilon,\hfield)\big] :
\end{equation}
the reduced temperature is the variable conjugate to the energy.
Substituting $\delta\varepsilon=\lambda\,\delta\varepsilon'$ and using
Eq.~\eqref{eq:filt_shom} (Lemma~\ref{lem:rescale}),
\[
  \omega_{\rm sing}(\lambda^{\sigma-1}t,\lambda^{b}\hfield)
  =\min_{\delta\varepsilon'}\lambda^{\sigma}
  \big[-t\,\delta\varepsilon'-T_c\,\lds_{\rm sing}(\delta\varepsilon',\hfield)\big]
  =\lambda^{\sigma}\,\omega_{\rm sing}(t,\hfield) ,
\]
and the choice $\lambda=|t|^{-1/(\sigma-1)}$ gives Eq.~\eqref{eq:filt_crit}
with $\Phi_\pm(z):=\omega_{\rm sing}(\pm1,z)$. The singular part of the heat
capacity is
$-T_c^{-1}\partial^{2}\omega_{\rm sing}/\partial t^{2}\propto|t|^{-\alpha}$,
and $m=-\partial\omega_{\rm sing}/\partial \hfield\propto|t|^{2-\alpha-\Delta_\hfield}$,
$\chi=\partial m/\partial \hfield\propto|t|^{2-\alpha-2\Delta_\hfield}$
(Eq.~\eqref{eq:tools_expo}), which gives Eq.~\eqref{eq:filt_critexp};
$\delta=\Delta_\hfield/\tilde\beta$ follows from the requirement that $m$ be
finite at $t=0$, $\hfield\neq0$ (Eq.~\eqref{eq:tools_match}).
\end{proof}

For $\lds_{\rm sing}=-B_\pm|\delta\varepsilon|^{\sigma}$ at $\hfield=0$ the
minimisation is explicit (Lemma~\ref{lem:balance},
Eq.~\eqref{eq:tools_crit}): the thermal energy lies on the side of the sign
of $t$,
\begin{equation}\label{eq:filt_omexpl}
  \delta\varepsilon_*=\pm\left(\frac{|t|}{\sigma T_cB_\pm}\right)^{1/(\sigma-1)} ,
  \qquad
  \omega_{\rm sing}=-\frac{\sigma-1}{\sigma}\,|t|\,|\delta\varepsilon_*|
  =-\frac{\sigma-1}{\sigma}\,\frac{|t|^{\sigma/(\sigma-1)}}{(\sigma T_cB_\pm)^{1/(\sigma-1)}} ,
\end{equation}
and the amplitude ratio of the heat capacity above and below $T_c$ is
$(B_-/B_+)^{1/(\sigma-1)}$, the ratio of the amplitudes of the spectrum on
the two sides of $\varepsilon_c$. Since
$|\delta\varepsilon_*|\propto|t|^{1/(\sigma-1)}=|t|^{1-\alpha}$, the
observables as functions of the energy have the exponents of
Eq.~\eqref{eq:filt_critexp} divided by $1-\alpha$, e.g.\
$m\propto|\delta\varepsilon|^{\sigma-b}$ with
$\sigma-b=\tilde\beta/(1-\alpha)$: these are the exponents of the
microcanonical ensemble~\cite{Kastner2000}.

\paragraph{Leading and subleading singular terms.} Whether
Proposition~\ref{prop:crit} applies depends on the regular quadratic term of
$\lds$ at $\varepsilon_c$.
\begin{itemize}\setlength{\itemsep}{2pt}
\item \emph{Divergent heat capacity.} If $\sigma>2$ and $\lds$ has no regular
  quadratic term, $\lds_{\rm sing}$ is the leading deviation, the curvature
  of $\lds$ vanishes at $\varepsilon_c$, and $0<\alpha<1$. All static
  exponents are fixed by the two exponents $(\sigma,b)$ of the spectrum. For
  the three-dimensional Ising class ($\alpha\simeq0.110$,
  $\tilde\beta\delta\simeq1.564$), $\sigma\simeq2.12$ and $b\simeq1.76$.
\item \emph{Finite heat capacity with a cusp\footnote{A \emph{cusp} of a
  function is a point at which it is continuous and finite but its slope
  diverges, as for $|t|^{p}$ with $0<p<1$ at $t=0$.}.} If the heat capacity
  at $T_c$ is finite, $c_\mu(T_c)=c_{\rm reg}>0$, then by
  Eq.~\eqref{eq:filt_curvmain} $\lds$ has the regular term $-(\delta\varepsilon)^{2}/2T_c^{2}c_{\rm reg}$ in addition to the terms
  of Eq.~\eqref{eq:filt_gcrit}, and a singular term with $\sigma>2$ is
  subleading to it. The quadratic term then fixes the minimiser
  (Lemma~\ref{lem:sublead}), $\delta\varepsilon_*=c_{\rm reg}T_c\,t$ to
  leading order, and
  $\omega_{\rm sing}(t,\hfield)\simeq-T_c\,\lds_{\rm sing}(c_{\rm reg}T_c\,t,\hfield)$:
  the exponents of the spectrum are transmitted unchanged,
  \begin{equation}\label{eq:filt_critexp2}
    2-\alpha=\sigma ,\qquad \Delta_\hfield=b ,\qquad
    \tilde\beta=\sigma-b ,\qquad \gamma=2b-\sigma ,\qquad
    \delta=\frac{b}{\sigma-b} ,
  \end{equation}
  in place of Eqs.~\eqref{eq:filt_crit} and~\eqref{eq:filt_critexp}. Here
  $\alpha<0$, as in the three-dimensional XY class ($\alpha\simeq-0.015$,
  $\sigma\simeq2.015$) and the Heisenberg class.
\item \emph{Marginal case.} For $\sigma=2$ the term
  $-B_\pm(\delta\varepsilon)^{2}$ is a finite curvature of $\lds$, different
  on the two sides: $\alpha=0$, with a jump of the heat capacity at $T_c$ if
  $B_+\neq B_-$, or with a logarithmic divergence if the term carries a
  logarithmic factor (Section~\ref{sec:filt_nonpower}).
\item \emph{Vanishing heat capacity.} For $1<\sigma<2$ and a leading singular
  term the curvature of $\lds$ diverges at $\varepsilon_c$, and the whole
  heat capacity vanishes at $T_c$ as $|t|^{(2-\sigma)/(\sigma-1)}$. This
  requires that no degree of freedom contributes a regular part to $c_\mu$.
\end{itemize}
A continuous transition is thus a singular term of $\lds$ at $\varepsilon_c$
with $\sigma>2$ (or $\sigma=2$ with a logarithm). It is the leading deviation
from the tangent if the heat capacity diverges, and it is subleading to a
regular quadratic term if the heat capacity remains finite. The formulas of
Sections~\ref{sec:filt_length}--\ref{sec:filt_rg} are written for the
first case; in the second, $\sigma-1$ is replaced by $1$ in the denominators
of the exponents.

\paragraph{Critical point at the edge of the spectrum.}
Assumption~\ref{ass:crit} places $\varepsilon_c$ in the interior of the
spectrum. In the mean-field (Curie--Weiss) model,
$\hat H=-(J_0/2V)\big(\sum_ix_i\big)^{2}$ with $x_i=\pm1$ on $V$
sites~\cite[Sec.~4.5.1]{Honerkamp}, it lies at the upper edge. The energy is
$\varepsilon=-J_0y^{2}/2$ with $y$ the magnetisation per site, and the
number of configurations with magnetisation $y$ is, by Stirling's formula,
$\exp[V(\ln2-\mathcal S(y))]$ with
$\mathcal S(y)=\tfrac{1+y}{2}\ln(1+y)+\tfrac{1-y}{2}\ln(1-y)=\tfrac12y^{2}+\tfrac1{12}y^{4}+\dots$,
so that
\[
  \frac{\lds(\varepsilon)}{k_B}=\ln2-\mathcal S\Big(\sqrt{-2\varepsilon/J_0}\Big)
  =\ln2+\frac{\varepsilon}{J_0}-\frac{\varepsilon^{2}}{3J_0^{2}}+\dots ,
  \qquad -\tfrac12J_0\le\varepsilon\le0 ,
\]
and $\lds$ is constant above $\varepsilon=0$, where there are no levels. The
slope of $\lds$ at the upper edge is $1/T_c$ with $k_BT_c=J_0$, and it is
larger everywhere below the edge. For $T>T_c$ no energy satisfies
Eq.~\eqref{eq:filt_stat}: $\fen$ decreases up to the edge and increases
beyond it, so that $\varepsilon_*=0$, with $S/V=k_B\ln2$ and $c_\mu=0$. For
$T<T_c$ the minimum lies inside the spectrum, at
$\varepsilon_*=-\tfrac32J_0(T_c-T)/T$ to leading order, and
$c_\mu\to\tfrac32k_B$ at $T_c^{-}$: the mean-field jump, $\alpha=0$. The
singularity is the arrival of the thermal energy at the upper edge of the
spectrum; it is the thermal counterpart of the quantum transition, at which
the edge is the ground state. With short-range couplings the spectrum
continues above $\varepsilon_c$, and Assumption~\ref{ass:crit} applies.

\subsection{Wave vector and correlation length}
\label{sec:filt_length}\label{sec:filt_corrq}

A length enters the spectral formulation through the dependence of
$\chi(\mathbf q)$ on the wave vector, i.e.\ of the response of the \lndos{}
to a modulated field, Eq.~\eqref{eq:filt_chiq}. Wave vectors are measured in
units of the inverse lattice spacing and lengths in units of the lattice
spacing. For an isotropic system and small $q=|\mathbf q|$, the expansion of
$\chi(\mathbf q)=\sum_{\mathbf r}e^{-i\mathbf q\cdot\mathbf r}\chi(\mathbf r)$
to second order is
\begin{equation}\label{eq:filt_elldef}
  \chi(\mathbf q)=\chi\,\big[1-q^{2}\ell_c^{2}+O(q^{4})\big] ,
  \qquad
  \ell_c^{2}:=-\frac{1}{\chi}\,\frac{\partial\chi(\mathbf q)}{\partial q^{2}}\bigg|_{q=0}
  =\frac{\sum_{\mathbf r}r^{2}\,\chi(\mathbf r)}{2d\,\sum_{\mathbf r}\chi(\mathbf r)} ,
\end{equation}
which defines the \emph{correlation length} $\ell_c$. By the second form, and
$k_BT\chi(\mathbf r)=G(\mathbf r)$ (Section~\ref{sec:filt_corr}), $\ell_c$ is
the range of the correlations.

\begin{assumption}[Homogeneity in the wave vector]\label{ass:q}
Near $\varepsilon_c$ the singular part of the \lndos{} in the presence of a
modulated field of amplitude $\hfield_{\mathbf q}$ and wave vector
$\mathbf q$ is homogeneous in $q$ as well: for $q\to0$,
\begin{equation}\label{eq:filt_shomq}
  \lds_{\rm sing}(\lambda\,\delta\varepsilon,\lambda^{b}\hfield_{\mathbf q};\lambda^{b_q}q)
  =\lambda^{\sigma}\lds_{\rm sing}(\delta\varepsilon,\hfield_{\mathbf q};q) ,
  \qquad b_q>0 ,
\end{equation}
with the exponent $b$ of the uniform field, and the response
$\partial^{2}\lds/\partial\hfield_{\mathbf q}^{2}$ at $\varepsilon_c$ is finite
for $q\neq0$.
\end{assumption}
With $\lambda=|\delta\varepsilon|^{-1}$ the wave vector enters only through
the ratio $q/|\delta\varepsilon|^{b_q}$: for the levels at the distance
$\delta\varepsilon$ from $\varepsilon_c$, a field modulated on a length that
exceeds $\ell\propto|\delta\varepsilon|^{-b_q}$ acts as a uniform field. The
assumption states that this single length governs the dependence on the wave
vector.

\paragraph{Homogeneity and physical dimensions.} The exponent $b_q$ is a
further exponent of the spectrum. It is not fixed by $(\sigma,b)$, and it is
not fixed by the physical dimensions of energy and length. Let $a_0$ denote
the lattice spacing and $J$ a microscopic energy (a coupling of the
Hamiltonian). The arguments of $\lds_{\rm sing}$ are the pure numbers
$\delta\varepsilon\,a_0^{d}/J$ and $q\,a_0$. A change of the units of energy
and length changes $J$ and $a_0$ together with $\delta\varepsilon$ and $q$,
and leaves both numbers unchanged; it holds for every function and gives no
relation between exponents. Eq.~\eqref{eq:filt_shomq} is a different
statement. In it $J$ and $a_0$ are fixed, and the level counts of one
Hamiltonian are compared at different energies and at different wave vectors
of the applied field. The wave vector is a parameter of the Hamiltonian that
is chosen independently of the energy, as the amplitude $\hfield$ is, and its
exponent is as free as $b$. With the constants restored,
\[
  \ell\propto a_0\,\Big(\frac{|\delta\varepsilon|\,a_0^{d}}{J}\Big)^{-b_q} :
\]
the dimension of length is carried by $a_0$, and no value of $b_q$ is
excluded. Dimensional analysis would fix the exponent only if $a_0$ did not
appear in $\ell$; then $\ell\propto(|\delta\varepsilon|/J)^{-1/d}$, i.e.\
$b_q=1/d$, whereas the three-dimensional Ising class has $b_q\simeq0.71$.
The exponent ceases to be independent under two further properties of the
spectrum (Section~\ref{sec:filt_size}): if the size of the system enters
only through the ratio $L/\ell$, then $b_V=d\,b_q$, and if hyperscaling
holds, $b_V=\sigma$, so that $b_q=\sigma/d$. The first property fails above
the upper critical dimension.

By Lemma~\ref{lem:rescale}, with $q$ as a further parameter,
\[
  \omega_{\rm sing}(\lambda^{\sigma-1}t,\lambda^{b}\hfield_{\mathbf q};\lambda^{b_q}q)
  =\lambda^{\sigma}\omega_{\rm sing}(t,\hfield_{\mathbf q};q) .
\]
Two derivatives with respect to $\hfield_{\mathbf q}$ at
$\hfield_{\mathbf q}=0$, Eq.~\eqref{eq:filt_chiq}, and the choice
$\lambda=|t|^{-1/(\sigma-1)}$ give
\begin{equation}\label{eq:filt_chiqscal}
  \chi(\mathbf q,t)=|t|^{-\gamma}\,X_\pm\big(q\,|t|^{-\nu}\big) ,
  \qquad
  \nu:=\frac{b_q}{\sigma-1} ,
\end{equation}
with $\gamma=(2b-\sigma)/(\sigma-1)$ as in Eq.~\eqref{eq:filt_critexp}. By
the definition~\eqref{eq:filt_elldef}, $\ell_c\propto|t|^{-\nu}$: the
susceptibility depends on the wave vector through the product $q\,\ell_c$.
This is tested in scattering experiments: the ratio
$\chi(\mathbf q,t)/\chi(0,t)$, measured at different temperatures on one
side of the transition, is a single function of $q\,\ell_c$.
Since $\chi(\mathbf q,t)$ is finite for $t\to0$ at $q\neq0$, the factor
$|t|^{-\gamma}$ must be cancelled by $X_\pm$, which requires
$X_\pm(x)\propto x^{-\gamma/\nu}$ for $x\to\infty$
(Eq.~\eqref{eq:tools_match}). Hence
\begin{equation}\label{eq:filt_fisher}
  \chi(\mathbf q,0)\propto q^{-(2-\eta)} ,
  \qquad 2-\eta=\frac{\gamma}{\nu}=\frac{2b-\sigma}{b_q} ,
\end{equation}
which is the Fisher relation $\gamma=\nu(2-\eta)$. In real space the
correlation function is the Fourier transform
$G(\mathbf r)=k_BT\int d^{d}q\,e^{i\mathbf q\cdot\mathbf r}\chi(\mathbf q)/(2\pi)^{d}$,
and the Fourier transform of a power is a power: for $0<2-\eta<d$,
\[
  G(r)\propto\frac{1}{r^{\,d-2+\eta}}\quad(t=0) ,\qquad
  G(r,t)=\frac{1}{r^{\,d-2+\eta}}\;\mathcal G_\pm\left(\frac{r}{\ell_c}\right)\quad(t\neq0) .
\]
The dimension $d$ enters only through the Fourier transform. For the
three-dimensional Ising class $\nu\simeq0.630$, i.e.\ $b_q\simeq0.71$.

\subsection{Size dependence and hyperscaling}
\label{sec:filt_size}

Let $\lds(\varepsilon,\hfield,V^{-1}):=(k_B/V)\ln\cN_V(V\varepsilon(\hfield))$
denote the logarithm of the level count per unit size of the finite system
(Eq.~\eqref{eq:filt_sV}), with $V^{-1}$ as a third argument.

\begin{assumption}[Homogeneity in the size]\label{ass:V}
\begin{equation}\label{eq:filt_shomV}
  \lds_{\rm sing}(\lambda\,\delta\varepsilon,\lambda^{b}\hfield,\lambda^{b_V}V^{-1})
  =\lambda^{\sigma}\lds_{\rm sing}(\delta\varepsilon,\hfield,V^{-1}) ,\qquad b_V>0 .
\end{equation}
\end{assumption}
As in the proof of Proposition~\ref{prop:crit},
$\omega_{\rm sing}(\lambda^{\sigma-1}t,\lambda^{b}\hfield,\lambda^{b_V}V^{-1})
=\lambda^{\sigma}\,\omega_{\rm sing}(t,\hfield,V^{-1})$, and the choice
$\lambda=|t|^{-1/(\sigma-1)}$ gives the finite-size scaling form
\begin{equation}\label{eq:filt_fss}
  \omega_{\rm sing}(t,\hfield,V^{-1})=|t|^{2-\alpha}\,
  \Phi_\pm\Big(\frac{\hfield}{|t|^{\Delta_\hfield}},\frac{V_c}{V}\Big) ,
  \qquad V_c:=|t|^{-\nu_*} ,\quad \nu_*:=\frac{b_V}{\sigma-1} ,
\end{equation}
with $2-\alpha$ and $\Delta_\hfield$ as in Eq.~\eqref{eq:filt_crit}, up to
the corrections $O(\ln V/V)$ to Eq.~\eqref{eq:filt_legendre}. The size
enters through the ratio $V_c/V$: $V_c$ is the number of components above
which the system has the properties of the infinite one. For
$V\lesssim V_c$, i.e.\ $|t|\lesssim V^{-1/\nu_*}$, the singularity is rounded,
as it must be in a finite system (Section~\ref{sec:filt_dist}). In
particular the susceptibility $\chi=|t|^{-\gamma}X_\pm(V_c/V)$ is finite at
$t=\hfield=0$, which requires $X_\pm(x)\propto x^{-\gamma/\nu_*}$ for
$x\to\infty$, so that $\chi(t=0,\hfield=0)\propto V^{\gamma/\nu_*}$.

\paragraph{Hyperscaling.} Under the rescaling of Eq.~\eqref{eq:filt_shomV},
$\delta\varepsilon\to\lambda\,\delta\varepsilon$, $\hfield\to\lambda^{b}\hfield$,
$V\to\lambda^{-b_V}V$, the singular part of the logarithm of the \emph{total}
level count, $V \lds_{\rm sing}/k_B$, is multiplied by $\lambda^{\sigma-b_V}$.
\emph{Hyperscaling} is the further assumption that it is invariant,
$b_V=\sigma$, i.e.
\begin{equation}\label{eq:filt_hyper}
  2-\alpha=\nu_* .
\end{equation}
Equivalently, $V_c\,\omega_{\rm sing}\propto|t|^{2-\alpha-\nu_*}$ does not
depend on $t$: the singular part of the thermodynamic potential is of order
$k_BT_c$ per group of $V_c$ components. With $b_V=\sigma$ and
$\lambda=V^{1/\sigma}$, Eq.~\eqref{eq:filt_shomV} gives
\[
  V\,\lds_{\rm sing}(\delta\varepsilon,\hfield,V^{-1})
  =\lds_{\rm sing}\big(\delta\varepsilon\,V^{1/\sigma},\,\hfield\,V^{b/\sigma},\,1\big) :
\]
the singular part of $\ln\cN_V$ depends on $V$ only through the combinations
$\delta\varepsilon\,V^{1/\sigma}$ and $\hfield\,V^{b/\sigma}$. This is a
statement about the level count that can be tested directly on spectra of
finite systems. No length has been used, and the statement holds also where
none exists: for the infinite-range Ising model the scaling variable is
$tV^{1/2}$, so $\nu_*=2$, in agreement with $\alpha=0$.

\paragraph{Relation to the correlation length.} If the size of a system of
linear dimension $L=V^{1/d}$ enters the singular part only through the ratio
$L/\ell_c$, then $V_c=\ell_c^{\,d}$, i.e.\ $b_V=d\,b_q$ and $\nu_*=\nu\,d$,
and Eq.~\eqref{eq:filt_hyper} becomes the hyperscaling relation
$2-\alpha=\nu\,d$. The susceptibility of the finite system at $t=\hfield=0$
then grows as $V^{\gamma/\nu_*}=L^{\gamma/\nu}=L^{2-\eta}$.\footnote{Above the
upper critical dimension the premise fails, because the uniform mode is not
controlled by the modes of small non-zero $\mathbf q$: the Ising model in
$d>4$ with periodic boundary conditions has $\nu=\tfrac12$ but
$\nu_*=2$~\cite{Brezin1982}, as for infinite range, so that
$2-\alpha=\nu_*$ holds and $2-\alpha=\nu\,d$ does not.}

\subsection{Logarithms and essential singularities}
\label{sec:filt_nonpower}

Eq.~\eqref{eq:filt_legendre} holds for any form of $\lds$, and homogeneity is
a special case. Other forms are treated by the rules of
Appendix~\ref{sec:tools_logs}.
\begin{itemize}\setlength{\itemsep}{0pt}
\item \emph{Logarithmic divergence.} For the two-dimensional Ising model
  ($\alpha=0$),
  $\lds_{\rm sing}\simeq-B\,(\delta\varepsilon)^{2}/\ln(1/|\delta\varepsilon|)$.
  Stationarity in Eq.~\eqref{eq:filt_omsing} gives
  $t\simeq2T_cB\,\delta\varepsilon_*/\ln(1/|\delta\varepsilon_*|)$, i.e.\
  $\delta\varepsilon_*\propto t\,\ln(1/|t|)$, and
  $c_\mu=T_c^{-1}\partial\delta\varepsilon_*/\partial t\propto\ln(1/|t|)$.
\item \emph{Essential singularity.} At a Kosterlitz--Thouless
  transition~\cite{Kosterlitz1974} the singular part of the thermodynamic
  potential is $\omega_{\rm sing}\propto\exp(-{\rm const}/\sqrt{t})$ for
  $t>0$, and all its derivatives vanish at $t=0$. Correspondingly $\lds$ has
  a finite curvature and an essential singularity at $\varepsilon_c$, and no
  pair $(\sigma,b)$ describes it.
\item \emph{Logarithmic correction at low temperature.} A heat capacity
  $c_\mu=\mathcal C_0\,T\ln(T_0/T)$ (a marginal Fermi liquid) gives, by the
  integrals of Section~\ref{sec:filt_power_mb} to leading logarithmic order,
  $u-\mu n-\varepsilon_0=\tfrac12\mathcal C_0T^{2}\ln(T_0/T)$ and
  $S/V=\mathcal C_0T\ln(T_0/T)$, hence
  $\lds\simeq\sqrt{\mathcal C_0\,\delta\varepsilon\,\ln(\varepsilon_1/\delta\varepsilon)}$
  with a constant $\varepsilon_1$: Bethe's formula with a logarithm.
\end{itemize}

\subsection{Relation to the scaling hypothesis and to the renormalisation group approach}
\label{sec:filt_rg}

\paragraph{Scaling hypothesis.} For a concave $\lds$ the Legendre
transform~\eqref{eq:filt_legendre} is invertible,
$\lds(\varepsilon)=\min_T[(\varepsilon-\omega(T))/T]$, so that the
homogeneity of $\lds_{\rm sing}$ and the homogeneity of $\omega_{\rm sing}$
imply one another. Assumption~\ref{ass:crit} is therefore equivalent to the
scaling hypothesis~\eqref{eq:scaling}, and the relations among the exponents
are the known ones. The spectral formulation does not add a relation; it
identifies what the exponents describe: $(\sigma,b,b_q,b_V)$ are properties
of the many-body level count about $\varepsilon_c$, defined for a given
Hamiltonian without reference to a free energy. The assumptions on the
spectrum also apply where no critical point is involved
(Assumptions~\ref{ass:power} and~\ref{ass:hom0}, and
Section~\ref{sec:beyond}).

\paragraph{Renormalisation group.} In the renormalisation group a rescaling
of lengths by the factor $\ell$ maps $t\to\ell^{\,y_t}t$ and
$\hfield\to\ell^{\,y_\hfield}\hfield$, and the singular free energy per unit
volume obeys
$f_{\rm sing}(t,\hfield)=\ell^{-d}f_{\rm sing}(\ell^{\,y_t}t,\ell^{\,y_\hfield}\hfield)$~\cite[Ch.~7]{Plischke}.
The choice $\ell=|t|^{-1/y_t}$ gives
$f_{\rm sing}=|t|^{d/y_t}\,\Phi_\pm(\hfield/|t|^{y_\hfield/y_t})$ and
$\ell_c\propto|t|^{-1/y_t}$. Comparison with Eqs.~\eqref{eq:filt_crit}
and~\eqref{eq:filt_chiqscal} identifies
\[
  y_t=\frac1\nu=\frac{\sigma-1}{b_q} ,\qquad
  y_\hfield=\frac{b}{b_q} ,\qquad
  d=\frac{\sigma}{b_q} ,
\]
the last being hyperscaling, $b_V=\sigma=d\,b_q$. The
renormalisation-group eigenvalues are ratios of the exponents of the
many-body spectrum, and the thermal energy scales as
$|\delta\varepsilon_*|\propto\ell_c^{-1/b_q}=\ell_c^{-(d-y_t)}$, with $d-y_t$
the scaling dimension of the energy density. A fixed point corresponds to a
Hamiltonian for which the homogeneity~\eqref{eq:filt_shom} holds exactly,
not only asymptotically for small $\delta\varepsilon$ and $\hfield$.

The two formulations differ in what is assumed. The renormalisation group
constructs the rescaling as a map on Hamiltonians and obtains the
homogeneity from the existence of a fixed point, which also explains why
different Hamiltonians share exponents and provides the means of computing
them. The spectral formulation takes the homogeneity as a property of the
spectrum of the given Hamiltonian; it does not compute the exponents, and it
applies unchanged to power laws that are not universal.

\section{Beyond the symmetry-breaking paradigm}
\label{sec:beyond}

In Landau's description~\cite{LandauLifshitz5} a phase is characterised by
the symmetry it breaks and by the corresponding order parameter; a
continuous transition is a change of symmetry, and the gapless excitations
of an ordered phase are the Goldstone modes of the broken symmetry.
Propositions~\ref{prop:power}--\ref{prop:crit} use none of these notions.
Their input is the \lndos{} and its dependence on the parameters of the
Hamiltonian. A symmetry, where there is one, is a reason for a
particular form of the \lndos{} and for its robustness; the thermodynamics
depends on the form. The formulation therefore applies without change to
transitions without a change of symmetry (the Lifshitz transition,
Section~\ref{sec:filt_scaling0}) or without a local order parameter (the
Kosterlitz--Thouless transition, Section~\ref{sec:filt_nonpower}), to
gapless phases whose power laws are not due to a broken symmetry (Fermi and
non-Fermi liquids, glasses), and to frustrated systems, in which several
orders compete. The ordered phase with Goldstone modes is treated first, as
the case in which both descriptions apply.

\subsection{Goldstone modes and degenerate manifolds}
\label{sec:goldstone}

A phase with a broken continuous symmetry has a spectrum that is scale-free
at low energy without any tuning~\cite{Goldstone1961,GSW1962}: a slow twist
of the order parameter, of wave vector $\mathbf q$, costs an energy that
vanishes with $q$, so that for short-range interactions there are gapless
excitations of energy $\nu(\mathbf q)\propto q^{z}$, the \emph{Goldstone
modes}\footnote{$z=1$ for antiferromagnets, superfluids and the acoustic
phonons of a crystal; $z=2$ for ferromagnets. In this section
$\nu(\mathbf q)$ denotes a dispersion; the exponent $\nu$ of
Section~\ref{sec:filt_length} does not occur.}. Their contribution to $\lds$
follows from a general statement on independent modes.

\begin{lemma}[Homogeneous dispersion]\label{lem:disp}
Let the modes have a dispersion $\nu(\mathbf q)\ge0$, the energy of the mode
of wave vector $\mathbf q$ ($\hbar=1$), that is homogeneous near its zeros,
$\nu(\lambda^{a_1}q_1,\dots,\lambda^{a_d}q_d)=\lambda\,\nu(\mathbf q)$
for all $\lambda>0$, with exponents $a_i\ge0$ and $\mathfrak a:=\sum_ia_i>0$. Then
the density of modes per unit volume is a power law of the mode energy $E$,
\[
  \varrho(E):=\int\frac{d^{d}q}{(2\pi)^{d}}\,\delta\big(E-\nu(\mathbf q)\big)
  =\varrho(1)\,E^{\mathfrak a-1} ,
\]
and independent bosonic modes of this density give
\begin{equation}\label{eq:filt_lemma}
  c_\mu=k_B\,\varrho(1)\,\Gamma(\mathfrak a+2)\,\zeta(\mathfrak a+1)\,(k_BT)^{\mathfrak a} ,
  \qquad
  \lds\propto (\delta\varepsilon)^{\mathfrak a/(\mathfrak a+1)} .
\end{equation}
\end{lemma}
\begin{proof}
Substitute $q_i=E^{a_i}x_i$: then $d^{d}q=E^{\mathfrak a}d^{d}x$,
$\nu(\mathbf q)=E\,\nu(\mathbf x)$ by homogeneity with $\lambda=E$, and
$\delta(E-E\,\nu(\mathbf x))=E^{-1}\delta(1-\nu(\mathbf x))$.
Components with $a_i=0$ are not rescaled and are integrated over their
finite range. The heat capacity is the bosonic one of
Section~\ref{sec:filt_power_mb} with $\mu=0$ and $r=\mathfrak a-1$, and
$\lds(\varepsilon)$ follows from Eq.~\eqref{eq:filt_sinv} with
$\kappa=\mathfrak a$.
\end{proof}
For Goldstone modes $a_i=1/z$ in all $d$ directions, $\mathfrak a=d/z$:
$c_{\mu}\sim T^{d/z}$ and $\lds\propto(\delta\varepsilon)^{\,d/(d+z)}$, i.e.\
$(\delta\varepsilon)^{3/4}$ for $d=3$, $z=1$ (the Debye result) and
$(\delta\varepsilon)^{3/5}$ for a three-dimensional ferromagnet.

\paragraph{Degenerate manifold of wave vectors.} If several ordering wave
vectors are exactly degenerate, as for the competing orders of a frustrated
system, the gapless modes have $\nu(\mathbf q)=0$ on a manifold of dimension
$d_m$ in $\mathbf q$ space, and $\nu=v\,|\mathbf q_\perp|^{z}$ depends only on
the $d-d_m$ components $\mathbf q_\perp$ of $\mathbf q$ normal to the
manifold. In Lemma~\ref{lem:disp}, $a_i=1/z$ for these components and
$a_i=0$ along the manifold, so that $\mathfrak a=(d-d_m)/z$:
\begin{equation}\label{eq:filt_manifold}
  \varrho(E)\propto E^{(d-d_m)/z-1} ,\qquad
  c_\mu\sim T^{(d-d_m)/z} ,\qquad
  \lds\propto(\delta\varepsilon)^{(d-d_m)/(d-d_m+z)} .
\end{equation}
The exponent is fixed by the codimension $d-d_m$ of the manifold, i.e.\ by
the number of components of $\mathbf q$ on which the mode energy depends,
and not by the dimension of space (Fig.~\ref{fig:manifold}). For a surface
of degenerate wave vectors, $d_m=d-1$, $c_\mu\sim T^{1/z}$ in any dimension.

\begin{figure}[htbp]
\centering
\tikzset{
  ldsgraph/.style={>=Stealth,font=\small,x=1.6cm,y=1.6cm},
  lds/.style={very thick,blue!60!black},
  edge/.style={very thick,orange!85!black},
  guide/.style={thin,dotted,gray!80!black},
  alt/.style={very thick,red!75!black,dashed},
}
\begin{tikzpicture}[ldsgraph]
  \draw[->] (-0.15,0) -- (2.5,0) node[below left]{$\delta\varepsilon$};
  \draw[->] (0,-0.1) -- (0,1.5) node[above]{$\lds$};
  \draw[lds,domain=0:2.3,samples=60,smooth] plot (\x,{0.8*\x^0.75});
  \draw[teal!80!black,very thick,domain=0:2.3,samples=60,smooth] plot (\x,{sqrt(\x)});
  \draw[alt,domain=0.001:2.3,samples=80,smooth] plot (\x,{sqrt(\x)*(\x/(\x+0.3))^(1/6)});
  \draw[guide] (0.3,0) -- (0.3,0.5);
  \node[below,font=\scriptsize] at (0.3,0) {$\sim\delta$};
  \node[font=\scriptsize,blue!60!black,anchor=north west] at (1.5,0.95) {$d_m=0$: $\sigma=\frac{d}{d+z}$};
  \node[font=\scriptsize,teal!80!black,anchor=south west] at (0.5,1.22) {$d_m>0$: $\sigma=\frac{d-d_m}{d-d_m+z}$};
  \node[font=\scriptsize,red!75!black,anchor=north west] at (0.35,0.35) {lifted by $\delta$};
\end{tikzpicture}
\caption{The \lndos{} above the ground state,
$\lds\propto(\delta\varepsilon)^{\sigma}$, for gapless modes on a degenerate
manifold of wave vectors of dimension $d_m$; $d_m=0$ is the Goldstone case.
When the degeneracy is lifted by an energy $\delta$ (red, dashed), the law
of the degenerate manifold holds for $\delta\varepsilon$ above a scale set
by $\delta$, and that of the selected order below,
Eq.~\eqref{eq:filt_neardeg}.}
\label{fig:manifold}
\end{figure}

\paragraph{Lifted degeneracy.} Let the degeneracy on the manifold be lifted
by a small energy $\delta$:
$\nu(\mathbf q)=\delta\,u(\mathbf q_\parallel)+v\,|\mathbf q_\perp|^{z}$,
with $0\le u\le1$ vanishing quadratically at isolated points (the selected
wave vectors). With $n:=d-d_m$, the integral over $\mathbf q_\perp$ gives
$\varrho(E)\propto\int d^{d_m}q_\parallel\,
[E-\delta\,u(\mathbf q_\parallel)]_+^{\,n/z-1}$, with $[x]_+:=\max(x,0)$.
For $E\gg\delta$ the integrand is $E^{n/z-1}$ on the whole manifold and
Eq.~\eqref{eq:filt_manifold} holds. For $E\ll\delta$ only the
neighbourhoods $u<E/\delta$ of the minima contribute, of volume
$\propto(E/\delta)^{d_m/2}$. Hence
\begin{equation}\label{eq:filt_neardeg}
  c_\mu\sim
  \begin{cases}
    T^{(d-d_m)/z} , & \delta\ll k_BT ,\\[2pt]
    \delta^{-d_m/2}\,T^{(d-d_m)/z+d_m/2} , & k_BT\ll\delta ,
  \end{cases}
\end{equation}
the second line being Lemma~\ref{lem:disp} with $a_i=\tfrac12$ along the
manifold. The thermodynamics is that of degenerate orders down to
$k_BT\sim\delta$, and the selected order appears only below. A manifold at
a finite energy gives the activated law~\eqref{eq:filt_gap} with a prefactor
fixed by the same geometry: the rotons of superfluid helium,
$\nu=\Delta+(q-q_0)^{2}/2m^{*}$ on a sphere ($d_m=2$), have
$\varrho(E)\propto(E-\Delta)^{-1/2}$ and
$c\propto\Delta^{2}T^{-3/2}e^{-\Delta/k_BT}$~\cite[Sec.~6.9.2]{Honerkamp}.

The broken symmetry is not an ingredient of Lemma~\ref{lem:disp} or of
Proposition~\ref{prop:power}: any Hamiltonian with the same low-energy level
count has the same power laws. The symmetry supplies the reason for the
absence of a gap, and the exponent $z$, for every perturbation that respects
it. Conversely, $\rho(\enrg)$ alone does not identify the modes (acoustic
phonons and antiferromagnetic magnons both give $c\sim T^{3}$); that
requires the dependence of $\rho$ on the field that couples to them.

\subsection{Frustration}
\label{sec:frustration}

In a frustrated system\footnote{Couplings that cannot all be satisfied at
once: antiferromagnetic bonds on a triangle, or first- and second-neighbour
couplings that favour different ordering wave vectors.} several ordering
tendencies compete, and the candidate ordered states are nearly degenerate.
In the spectrum this is a definite statement: an unusually large number of
many-body levels lies just above the ground state, and the thermodynamics
is decided by the form of $\ln\cN$ in this group of levels. Three cases
occur: a degenerate manifold of ordering wave vectors
(Section~\ref{sec:goldstone}), a finite number of competing states, and a
number of ground states that is exponential in the size.

\begin{figure}[htbp]
  % s_0 in blue, s_1 in teal
\centering
\tikzset{
  ldsgraph/.style={>=Stealth,font=\small,x=1.6cm,y=1.6cm},
  lds/.style={very thick,blue!60!black},
  edge/.style={very thick,orange!85!black},
  guide/.style={thin,dotted,gray!80!black},
  alt/.style={very thick,red!75!black,dashed},
}
\begin{tikzpicture}[ldsgraph]
  % two competing families and the common tangent of their logdos
  \draw[->] (-0.15,0) -- (2.5,0) node[below left]{$\delta\varepsilon$};
  \draw[->] (0,-0.1) -- (0,1.5) node[above]{$\lds$};
  \draw[lds,domain=0:2.3,samples=50,smooth] plot (\x,{0.55*sqrt(\x)});
  \draw[teal!80!black,very thick,domain=0.5:2.3,samples=50,smooth] plot (\x,{0.95*sqrt(\x-0.5)});
  % common tangent: touches 0.55 sqrt(x) at x1 and 0.95 sqrt(x-0.5) at x2 with equal slope
  \draw[edge] (-0.1,0.083) -- (2.4,1.452);
  \fill[orange!85!black] (0.252,0.276) circle (1.4pt);
  \fill[orange!85!black] (1.252,0.824) circle (1.4pt);
  \draw[guide] (0.5,0) -- (0.5,0.1);
  \node[below,font=\scriptsize] at (0.5,0) {$\delta_1$};
  \node[font=\scriptsize,blue!60!black,anchor=north west] at (1.3,0.55) {$\lds_0$};
  \node[font=\scriptsize,teal!80!black,anchor=north west] at (2.2,1.0) {$\lds_1$};
  \node[font=\scriptsize,orange!85!black,anchor=south east] at (1.1,0.85) {slope $1/T_\times$};
\end{tikzpicture}
\caption{Two competing families of levels above the ground state: $\lds_0$
built on the ground state, and $\lds_1$, of larger amplitude, built on a
state that lies higher by $\delta_1$. The \lndos{} is the larger of the two.
The common tangent of slope $1/T_\times$ is a discontinuous transition
between the two orders, Eq.~\eqref{eq:filt_compete}.}
\label{fig:frustration}
\end{figure}

\paragraph{Finite number of competing states.} Let the low-energy spectrum
consist of $K$ families of levels, $k=0,\dots,K-1$. The family $k$ is the
set of eigenstates built on a candidate ordered state of energy
$\enrg_0+V\delta_k$ ($\delta_0=0\le\delta_k$), with the level count
$\cN_k(\enrg)=\exp[V\lds_k(\delta\varepsilon-\delta_k)/k_B]$,
$\delta\varepsilon=(\enrg-\enrg_0)/V$, and with an order parameter of density
$m_k$ conjugate to a field $\hfield_k$ (Fig.~\ref{fig:frustration}). The
families are disjoint sets of eigenstates of one system, not independent
subsystems. The level counts therefore add, $\cN=\sum_k\cN_k$, and since
$\Xi$ is linear in $\rho$, $\Xi=\sum_k\Xi_k$ (a product holds for
independent subsystems, Section~\ref{sec:filt_dist}). Since
$\max_k\Xi_k\le\Xi\le K\max_k\Xi_k$, the potential is the smallest of the
potentials of the families, up to a term $O(k_BT\ln K/V)$. To first order
in the fields, which shift the levels of the family $k$ by $-Vm_k\hfield_k$,
\begin{equation}\label{eq:filt_compete}
\begin{gathered}
  \omega(T,\{\hfield_k\})-\varepsilon_0=\min_k\,g_k ,\qquad
  g_k:=\delta_k-m_k\hfield_k+\omega_k(T) ,
  \\
  \omega_k(T):=\min_{\delta\varepsilon}\big[\delta\varepsilon-T\,\lds_k(\delta\varepsilon)\big] .
\end{gathered}
\end{equation}
For $\lds_k=\mathcal{A}_k\,(\delta\varepsilon)^{\sigma_k}$,
Eq.~\eqref{eq:filt_lowT} gives
$\omega_k=-\tfrac{1-\sigma_k}{\sigma_k}(\sigma_k \mathcal{A}_k T)^{1/(1-\sigma_k)}$.
\begin{itemize}\setlength{\itemsep}{0pt}
\item \emph{Selection by the level count.} For exactly degenerate states,
  $\delta_k=0$, and without fields, the family realised at $T>0$ is the one
  with the lowest $\omega_k(T)$: for $T\to0$ the family with the smallest
  $\sigma_k$ and, among equal $\sigma_k$, the largest $\mathcal{A}_k$, i.e.\
  the family with the largest number of low-lying levels. An order is
  selected at every $T>0$ although no state is preferred at $T=0$
  (\emph{order by disorder}~\cite{Villain1980,Henley1989}).
\item \emph{Small temperature scale.} Let the family with more low-lying
  levels lie higher: $\delta_1>0$, with $\sigma_1=\sigma_0=\sigma$ and
  $\mathcal{A}_1>\mathcal{A}_0$. Then $g_0<g_1$ at low temperature, and
  $g_0=g_1$ at the temperature $T_\times$ given by
  \[
    \delta_1=\frac{1-\sigma}{\sigma}\,(\sigma T_\times)^{1/(1-\sigma)}
    \Big[\mathcal{A}_1^{1/(1-\sigma)}-\mathcal{A}_0^{1/(1-\sigma)}\Big] ,
    \qquad\text{i.e.}\qquad T_\times\propto\delta_1^{\,1-\sigma} .
  \]
  At $T_\times$ the system passes from the family $0$ to the family $1$ in a
  discontinuous transition. The temperature $T_\times$ is set by the
  near-degeneracy $\delta_1$, not by the couplings, and vanishes with it.
\item \emph{Coexistence at zero field.} Let $g_k$ at zero fields have the
  same value $g(T)$ for the families of a set $\mathcal K$ (exact degeneracy
  with equal spectra, or the point $T_\times$). For small fields,
  $\omega-\varepsilon_0=g(T)-\max_{k\in\mathcal K}m_k\hfield_k$. As a function
  of one field $\hfield_k$, the others being zero, $\omega$ has a kink at
  $\hfield_k=0$, and the order parameter $-\partial\omega/\partial \hfield_k$
  jumps from $0$ to $m_k$, for every $k\in\mathcal K$: an infinitesimal field
  selects the corresponding order. At zero fields the equilibrium state is
  the mixture of the families of $\mathcal K$ with equal weights, and no
  single order parameter characterises it.
\end{itemize}
In the response, competing orders are several maxima of
$\chi(\mathbf q)=2T\,\partial^{2}\lds/\partial\hfield_{\mathbf q}^{2}$,
Eq.~\eqref{eq:filt_chiq}, as a function of $\mathbf q$, and order sets in at
the wave vector at which $\chi(\mathbf q)$ first diverges. When the maxima
are degenerate on a manifold of wave vectors, the fluctuations of the
gapless modes on the manifold (Section~\ref{sec:goldstone}) suppress the
divergence, and the transition is driven to low temperature or becomes
discontinuous~\cite{Brazovskii1975}.

\paragraph{Exponentially many ground states.} If the number of ground
states is exponential in the size, $d_0=e^{V\lds_0/k_B}$, the Boltzmann
weights tend to $1/d_0$ on the ground states for $T\to0$, and
$S\to k_B\ln d_0=V\lds_0$. Assumption~\ref{ass:power} fails, and the entropy
sum rule $\int_0^\infty C_\mu\,dT/T=S(\infty)-S(0)$ is reduced by $V\lds_0$.
For the triangular Ising antiferromagnet $\lds_0=0.323\,k_B$ per
spin~\cite{Wannier1950}, for spin ice
$\lds_0\simeq\tfrac12k_B\ln\tfrac32$ per
spin~\cite{Pauling1935,Ramirez1999}. If the first excitations are separated
by a gap $\Delta$, the heat capacity is activated, Eq.~\eqref{eq:filt_gap}
(Fig.~\ref{fig:residual}a).

Let a weak perturbation spread the $d_0$ degenerate levels over energies
$\enrg_j$ within a local scale $W$ above $\enrg_0$, small compared with the
gap $\Delta$ to the other levels (Fig.~\ref{fig:residual}b). For
$k_BT\ll\Delta$,
$e^{\beta\enrg_0}\Xi=d_0\,\langle e^{-\beta(\enrg-\enrg_0)}\rangle_{\rm b}$, with
$\langle\cdot\rangle_{\rm b}$ the average over the group with equal weights.
The expansion of its logarithm in the cumulants of $\enrg$ within the group,
with the extensive variance $Vv_2$ ($v_2\sim W^{2}$ per component), gives
\[
  \frac{S}{V}=\lds_0-\frac{v_2}{2k_BT^{2}}+\dots ,\qquad
  c_\mu=\frac{v_2}{k_BT^{2}}+\dots \qquad (W\ll k_BT\ll\Delta) .
\]
The entropy $\lds_0$ is released below $k_BT\sim W$: $c_\mu(T)$ has a maximum
at $k_BT\sim W$, with a $T^{-2}$ tail above it, in addition to the maximum
at the scale of the couplings. Power laws, if any, appear only below $W$ and
are those of the level count within the group. Exact spectra of small
clusters of the spin-$\tfrac12$ kagome antiferromagnet show such a group: a
number of singlet levels below the lowest triplet that grows exponentially
with the size~\cite{Waldtmann1998}. Experimentally, the ordering temperature
of frustrated magnets lies far below the scale of the
couplings~\cite{Ramirez1994}.

\begin{figure}[htbp]
\centering
\tikzset{
  ldsgraph/.style={>=Stealth,font=\small,x=1.6cm,y=1.6cm},
  lds/.style={very thick,blue!60!black},
  edge/.style={very thick,orange!85!black},
  guide/.style={thin,dotted,gray!80!black},
  alt/.style={very thick,red!75!black,dashed},
}
\begin{tikzpicture}[ldsgraph]
  % panel (a): exponentially many ground states, gap above them
  \node[font=\normalsize] at (-0.45,1.55) {(a)};
  \draw[->] (-0.15,0) -- (2.5,0) node[below left]{$\delta\varepsilon$};
  \draw[->] (0,-0.1) -- (0,1.5) node[above]{$\lds$};
  % s - s_0 = A x ln(1 + B/x): -A x ln x near x = 0 (vertical tangent), saturating at large x;
  % sampled at x = 2.3 t^4 so that the steep start is resolved
  \draw[edge] (0,0.5) -- (0,1.05);
  \draw[lds,variable=\t,domain=0.1:1,samples=150,smooth]
    plot ({2.3*\t*\t*\t*\t},{0.5+2.84*(2.3*\t*\t*\t*\t)*ln(1+0.3/(2.3*\t*\t*\t*\t))});
  \draw[lds] (0,0.5) -- ({2.3*0.0001},{0.5+2.84*0.00023*ln(1+0.3/0.00023)});
  \fill[blue!60!black] (0,0.5) circle (1.4pt);
  \node[left,font=\scriptsize] at (0,0.5) {$\lds_0$};
  \node[font=\scriptsize,orange!85!black,anchor=west,align=left] at (0.75,0.55) {vertical tangent at $\lds_0$:\\ activated, gap $\Delta$};
  \draw[thin,orange!85!black] (0.75,0.55) -- (0.03,0.62);
\end{tikzpicture}\hfill
\begin{tikzpicture}[ldsgraph]
  % panel (b): the degenerate levels spread over a narrow group of width W
  \node[font=\normalsize] at (-0.45,1.55) {(b)};
  \draw[->] (-0.15,0) -- (2.5,0) node[below left]{$\delta\varepsilon$};
  \draw[->] (0,-0.1) -- (0,1.5) node[above]{$\lds$};
  \fill[gray!15] (0,0) rectangle (0.15,1.4);
  \draw[lds,domain=0:0.15,samples=30,smooth] plot (\x,{0.5*sqrt(\x/0.15)});
  \draw[lds,domain=0.152:2.3,samples=120,smooth] plot (\x,{0.5+0.8*((\x-0.15)/2.15)*(1-ln((\x-0.15)/2.15))});
  \draw[guide] (0,0.5) -- (0.15,0.5);
  \node[left,font=\scriptsize] at (0,0.5) {$\lds_0$};
  \node[font=\scriptsize,anchor=west] at (0.2,1.3) {$\sim W$};
  \node[font=\scriptsize,blue!60!black,anchor=north west,align=left] at (0.9,0.9) {entropy $\lds_0$ released\\ below $k_BT\sim W$};
\end{tikzpicture}
\caption{The \lndos{} above the ground state.
(a)~Exponentially many ground states and a gap $\Delta$: $\lds$ starts at the
residual entropy $\lds_0$ and rises with a logarithmically divergent slope
(orange; drawn with an enlarged prefactor), $\lds-\lds_0\propto-x\ln x$ with
$x\propto\delta\varepsilon/\Delta$, which gives the activated
law~\eqref{eq:filt_gap}.
(b)~The same levels spread over a narrow group of width $W$ (grey): $\lds$
rises from $0$ to $\lds_0$ within $\delta\varepsilon\sim W$, and the residual
entropy is released in a maximum of $c_\mu$ at $k_BT\sim W$.}
\label{fig:residual}
\end{figure}

\subsection{Disorder}
\label{sec:disorder}

Quenched disorder\footnote{Random terms in the Hamiltonian (impurities,
random bonds or fields) that are fixed in each sample and do not
equilibrate.} does not create levels; it redistributes them. Each random
term is a parameter of the Hamiltonian in the sense of
Section~\ref{sec:filt_dist}, and as it is varied the levels move along the
energy axis without changing their number, Eq.~\eqref{eq:filt_gA}:
degenerate levels are split into bands whose width is set by the disorder,
levels are moved into gaps and beyond band edges, and the exponents of
$\lds$ may change. For short-range disorder $\ln\cN$ is extensive and
self-averaging\footnote{A quantity is \emph{self-averaging} if its relative
fluctuations from sample to sample vanish for $V\to\infty$.}, so that
$\lds(\varepsilon)$ is the same for almost every large sample. The
consequences follow from its form (Fig.~\ref{fig:disorder}).

\begin{figure}[htbp]
\centering
\tikzset{
  ldsgraph/.style={>=Stealth,font=\small,x=1.4cm,y=1.6cm},
  lds/.style={very thick,blue!60!black},
  clean/.style={very thick,red!75!black},
  guide/.style={thin,dotted,gray!80!black},
}
\begin{tikzpicture}[ldsgraph]
  % panel (a): single-particle density of states at a band edge, clean and with a Lifshitz tail
  \node[font=\normalsize] at (-1.45,1.55) {(a)};
  \draw[->] (-1.1,0) -- (2.2,0) node[below left]{$E$};
  \draw[->] (-1.0,-0.1) -- (-1.0,1.5) node[above]{$\varrho$};
  \draw[clean] (-1.0,0) -- (0,0);
  \draw[clean,domain=0:2,samples=50,smooth] plot (\x,{0.9*sqrt(\x)});
  \draw[lds,domain=-0.59:0.2,samples=80,smooth] plot (\x,{0.9*0.68*exp(-0.3/((\x+0.6)^1.5))});
  \draw[lds,domain=0.2:2,samples=50,smooth] plot (\x,{0.9*sqrt(\x+0.0)});
  \draw[guide] (0,0) -- (0,1.0);
  \node[below,font=\scriptsize] at (0,0) {$E_0$};
  \node[below,font=\scriptsize] at (-0.6,0) {$E_0'$};
  \node[font=\scriptsize,red!75!black,anchor=west,align=left] at (0.9,0.5) {clean: $B|E-E_0|^{r}$};
  \node[font=\scriptsize,blue!60!black,anchor=south east,align=right] at (0.05,0.35) {Lifshitz\\ tail};
\end{tikzpicture}\hfill
\begin{tikzpicture}[ldsgraph]
  % panel (b): a gap of the clean spectrum filled by the disorder
  \node[font=\normalsize] at (-0.45,1.55) {(b)};
  \draw[->] (-0.1,0) -- (3.2,0) node[below left]{$E$};
  \draw[->] (0,-0.1) -- (0,1.5) node[above]{$\varrho$};
  \draw[clean,domain=0:0.8,samples=40,smooth] plot (\x,{0.9*sqrt(0.8-\x)});
  \draw[clean] (0.8,0) -- (1.8,0);
  \draw[clean,domain=1.8:3.0,samples=40,smooth] plot (\x,{0.9*sqrt(\x-1.8)});
  \draw[lds,domain=0:3.0,samples=120,smooth] plot (\x,{0.9*sqrt(max(0.8-\x,0)+0.0)*0+0.9*sqrt(0.5*(0.8-\x+sqrt((0.8-\x)^2+0.03)))*(\x<1.3)+0.9*sqrt(0.5*(\x-1.8+sqrt((\x-1.8)^2+0.03)))*(\x>=1.3)+0.07});
  \draw[guide] (1.3,0) -- (1.3,0.6);
  \node[below,font=\scriptsize] at (1.3,0) {$\mu$};
  \node[font=\scriptsize,blue!60!black,anchor=south] at (1.3,0.62) {$\varrho(\mu)>0$: $c_\mu\propto T$};
  \node[font=\scriptsize,red!75!black,anchor=north] at (1.3,-0.3) {clean gap: activated};
\end{tikzpicture}

\vspace{3mm}

\begin{tikzpicture}[ldsgraph]
  % panel (c): distributions of local splittings
  \node[font=\normalsize] at (-0.45,1.55) {(c)};
  \draw[->] (-0.1,0) -- (2.6,0) node[below left]{$\Delta$};
  \draw[->] (0,-0.1) -- (0,1.5) node[above]{$P(\Delta)$};
  \fill[gray!15] (0,0) rectangle (0.6,1.4);
  \node[font=\scriptsize,anchor=south] at (0.3,1.1) {$\Delta<k_BT$};
  \draw[lds] (0,0.6) -- (2.4,0.6);
  \draw[teal!80!black,domain=0.1:2.4,samples=60,smooth] plot (\x,{0.45/sqrt(\x)});
  \draw[red!75!black,very thick,dashed,domain=0:2.4,samples=60,smooth] plot (\x,{0.6*sqrt(\x/1.5)});
  \node[font=\scriptsize,blue!60!black,anchor=south west] at (0.65,0.65) {$\theta=1$: $P_0$ (glass)};
  \node[font=\scriptsize,teal!80!black,anchor=south west] at (0.75,1.0) {$\theta<1$: $\chi$ diverges};
  \node[font=\scriptsize,red!75!black,anchor=south west] at (1.9,0.8) {$\theta>1$};
\end{tikzpicture}\hfill
\begin{tikzpicture}[ldsgraph]
  % panel (d): the resulting logdos above the ground state
  \node[font=\normalsize] at (-0.45,1.55) {(d)};
  \draw[->] (-0.1,0) -- (2.6,0) node[below left]{$\delta\varepsilon$};
  \draw[->] (0,-0.1) -- (0,1.5) node[above]{$\lds$};
  \draw[clean,domain=0.002:2.4,samples=120,smooth] plot (\x,{0.55*(\x/2.4)*(1-ln(\x/2.4))});
  \draw[lds,domain=0:2.4,samples=60,smooth] plot (\x,{0.6*sqrt(\x)});
  \draw[teal!80!black,domain=0:2.4,samples=60,smooth] plot (\x,{0.85*\x^(1/3)});
  \node[font=\scriptsize,red!75!black,anchor=north west] at (1.3,0.42) {clean gap: $-x\ln x$};
  \node[font=\scriptsize,blue!60!black,anchor=north west] at (1.75,0.74) {$(\delta\varepsilon)^{1/2}$};
  \node[font=\scriptsize,teal!80!black,anchor=south west] at (1.5,1.2) {$(\delta\varepsilon)^{\theta/(\theta+1)}$, $\theta<1$};
\end{tikzpicture}
\caption{Disorder: clean spectra in red, disordered ones in blue and green.
(a)~Single-particle density of states at a band edge: the clean power law
$B|E-E_0|^{r}$ and the Lifshitz tail of the disordered system, which
vanishes with all its derivatives at the new edge $E_0'$.
(b)~A gap of the clean spectrum filled by the disorder: a non-zero
$\varrho(\mu)$ replaces the activated law~\eqref{eq:filt_gap} by the linear
law~\eqref{eq:filt_sommerfeld}.
(c)~Distributions $P(\Delta)\propto\Delta^{\theta-1}$ of the splittings of
local two-level systems; the groups with $\Delta<k_BT$ (grey) are thermally
excited.
(d)~The \lndos{} in these cases: $-x\ln x$ for the clean gap,
$(\delta\varepsilon)^{1/2}$ for a filled gap or a flat $P(\Delta)$, and
$(\delta\varepsilon)^{\theta/(\theta+1)}$ for a power-law $P(\Delta)$.}
\label{fig:disorder}
\end{figure}
\begin{itemize}\setlength{\itemsep}{2pt}
\item \emph{Band edges and gaps.} At a band edge $E_0$ of the clean system,
  with $\varrho(E)=B|E-E_0|^{r}$ on the side of the band and zero on the
  other, the ground-state energy has the singular part
  $\propto|\mu-E_0|^{r+2}$ on the side of the band,
  Eq.~\eqref{eq:lif_ground}: the Lifshitz transition of
  Section~\ref{sec:filt_scaling0} (for a three-dimensional band edge,
  $r=\tfrac12$, the third derivative with respect to $\mu$ diverges: a
  transition of order $\tfrac52$). In the disordered system the band edge moves to $E_0'$, the
  lowest energy that any configuration of the disorder allows. A state of
  energy $E$ just above $E_0'$ must be confined to a region throughout which
  the random potential is close to its lowest value. The kinetic energy of
  confinement, $\propto R^{-2}$ for a region of linear size $R$, requires
  $R\gtrsim(E-E_0')^{-1/2}$, and if each site is favourable independently
  with probability $p<1$, such a region occurs with probability $p^{R^{d}}$.
  The density of states near $E_0'$ is therefore exponentially small (a
  Lifshitz tail~\cite{Lifshitz1964}),
  $\varrho(E)\propto\exp[-{\rm const}\,(E-E_0')^{-d/2}]$, and the band-edge
  singularity is replaced by an essential singularity: the transition is
  rounded. Likewise, if the disorder leaves a non-zero $\varrho(\mu)$ in a
  gap of the clean system, the activated law~\eqref{eq:filt_gap} is replaced
  by the linear law~\eqref{eq:filt_sommerfeld}.
\item \emph{Flat distribution of local splittings.} In a glass, small groups
  of atoms have two nearly degenerate configurations. The energy difference
  $\Delta$ between the two lowest levels of such a group, its
  \emph{splitting}, is distributed with a density $P(\Delta)$ per unit
  volume that tends to a constant $P_0$ for
  $\Delta\to0$~\cite{Anderson1972,Phillips1972}. For independent two-level
  systems a many-body level is specified by the set of excited groups. With
  $n(\Delta)$ the fraction of excited groups of splitting $\Delta$, the
  logarithm of the number of such sets at a given energy is dominated by its
  largest term,
  \begin{gather*}
    \lds(\varepsilon)=\max_{n}\;k_B\int_0^{\infty}d\Delta\,P(\Delta)\,
    \big[-n\ln n-(1-n)\ln(1-n)\big]
    \\
    \text{at fixed}\quad
    \delta\varepsilon=\int_0^{\infty}d\Delta\,P(\Delta)\,\Delta\,n(\Delta) .
  \end{gather*}
  With the Lagrange multiplier $1/T=\partial\lds/\partial\varepsilon$,
  Eq.~\eqref{eq:filt_stat}, the maximum is at $n=1/(e^{\Delta/k_BT}+1)$. For
  $P=P_0$, with $\int_0^{\infty}dx\,x/(e^{x}+1)=\pi^{2}/12$,
  \begin{equation}\label{eq:filt_tls}
    \delta\varepsilon_*=\frac{\pi^{2}}{12}\,P_0\,(k_BT)^{2} ,\qquad
    \lds(\varepsilon)=k_B\sqrt{\tfrac{\pi^{2}}{3}\,P_0\,\delta\varepsilon} ,
    \qquad
    c_\mu=\frac{\pi^{2}}{6}\,k_B^{2}P_0\,T .
  \end{equation}
  This is Assumption~\ref{ass:power} with $\sigma=\tfrac12$: the \lndos{} of
  a glass has the form of that of a Fermi liquid, in an insulator and in any
  dimension.
\item \emph{Power-law distribution of local splittings.} For
  $P(\Delta)=P_\theta\Delta^{\theta-1}$ at small $\Delta$ ($\theta>0$), the
  same maximisation gives
  \[
    \delta\varepsilon_*=P_\theta\,\Gamma(\theta+1)\,\etaD(\theta+1)\,(k_BT)^{\theta+1} ,\qquad
    \lds\propto(\delta\varepsilon)^{\theta/(\theta+1)} ,\qquad
    c_\mu\propto T^{\theta} .
  \]
  If the two-level systems are pairs of spins $\tfrac12$ (moment $m_0$ each)
  with a singlet ground state and a triplet at $\Delta$, the susceptibility
  of one pair is, by Eq.~\eqref{eq:filt_mb_response},
  $\Var(\hat M)/k_BT=8m_0^{2}/k_BT\,(e^{\Delta/k_BT}+3)$: a Curie law for
  $\Delta\ll k_BT$, and exponentially small for $\Delta\gg k_BT$. Hence
  \[
    \chi=\int_0^{\infty}d\Delta\,P(\Delta)\,\frac{8m_0^{2}}{k_BT\,(e^{\Delta/k_BT}+3)}
    \;\propto\;T^{\theta-1} :
  \]
  the number of pairs with $\Delta<k_BT$, which respond as free moments, is
  $\propto T^{\theta}$, and each contributes a Curie term $\propto1/T$. The
  susceptibility diverges for $\theta<1$. The exponent $\theta$ is not
  universal and varies continuously with the parameters of the Hamiltonian
  (random-singlet and Griffiths phases\footnote{Phases of disordered
  magnets in which rare, weakly coupled spins or regions dominate the
  response at low temperature.}~\cite{BhattLee1982,Griffiths1969}): power
  laws without a critical point and without universality.
\item \emph{Stability of a critical point.} Let the disorder shift the local
  value of $t$ by independent random amounts of variance $w^{2}$. Over the
  $V_c\sim|t|^{-\nu_*}$ components of Section~\ref{sec:filt_size} the mean
  shift is of order $w\,V_c^{-1/2}=w\,|t|^{\nu_*/2}$. It is negligible
  compared with $|t|$ for $t\to0$ if and only if $\nu_*>2$, i.e.\ $\alpha<0$
  by Eq.~\eqref{eq:filt_hyper}. This is the Harris
  criterion~\cite{Harris1974} ($\nu\,d>2$ with $\nu_*=\nu\,d$). In terms of
  the spectrum: weak disorder that couples to the local energy density
  changes the exponents precisely when the curvature of $\lds$ vanishes at
  $\varepsilon_c$ (the case $0<\alpha<1$ of Section~\ref{sec:filt_critical}).
\end{itemize}
A limit of the formulation appears here. The density of states does not
register whether the eigenstates are extended or localised: at an Anderson
transition of non-interacting electrons $\rho$ is smooth, and the
thermodynamics with it. That distinction lies in the eigenvectors.

\subsection{Fermi liquids and non-Fermi liquids}
\label{sec:nfl}
A Fermi liquid has $c_\mu\propto T$, i.e.\ $\kappa=1$ and $\sigma=\tfrac12$ in
Eq.~\eqref{eq:filt_powerT}: its low-energy many-body spectrum is that of
independent quasiparticles, and
$\lds=\sqrt{2\,\mathcal C_0\,\delta\varepsilon}$ (Bethe's formula,
Section~\ref{sec:filt_power_mb}). In the present terms, a non-Fermi liquid
is a metal whose level count above the ground state has another form.
Two-dimensional fermions coupled to a gapless boson (a gauge field, or the
order-parameter fluctuations at a quantum critical point with dynamical
exponent $z=3$) have $c\sim T^{2/3}$~\cite{HLR1993,Lohneysen2007}, i.e.\
$\kappa=\tfrac23$ and
\[
  \lds\propto(\delta\varepsilon)^{2/5}
  \qquad\text{against}\qquad
  \lds\propto(\delta\varepsilon)^{1/2}\ \ \text{for a Fermi liquid} ,
\]
in agreement with $\sigma=d/(d+z)$, Eq.~\eqref{eq:filt_qcp}: there are more
low-lying levels than in any Fermi liquid. The marginal Fermi liquid is the
logarithmic case of Section~\ref{sec:filt_nonpower}. In these systems the
description by independent quasiparticles fails, and the interaction
mediated by the gapless boson is long-ranged in space and time. Neither
affects Propositions~\ref{prop:power} and~\ref{prop:scaling0}, which require
only the stated form of the level count: the concavity of $\lds$, which
rests on a division of the system into parts, is not guaranteed there, and
the steepest-descent method does not need it.

\section{Conclusion}
\label{sec:principle}\label{sec:summary}

The thermodynamics of a Hamiltonian is encoded in the \lndos{}
$\lds(\varepsilon(\hfield))$ on the domain $\varepsilon\ge\varepsilon_0(\hfield)$
and in the lower edge $\varepsilon_0(\hfield)$ of this domain. Both are
properties of the spectrum and do not depend on $T$. In the limit of large
size a state $(T,\hfield)$ selects the single energy $\varepsilon_*(T,\hfield)$
at which the slope of $\lds$ is $1/T$, and every thermodynamic quantity is a
function of $\lds$ and of its derivatives at that energy
(Table~\ref{tab:relations}). The thermodynamic potential is non-analytic
only where $\lds$ is non-analytic or has zero curvature at $\varepsilon_*$,
where $\varepsilon_*$ jumps, or where $\varepsilon_0(\hfield)$ is non-analytic
(Section~\ref{sec:transitions}). The exponents and the scaling functions of
a transition or of a low-temperature regime are fixed by the form of $\lds$
about the energy involved (Table~\ref{tab:assumptions}), without reference
to a symmetry, an order parameter or a fixed point; a symmetry, where there
is one, is a reason for the form.

\begin{table}[htbp]
\centering
\small
\caption{The basic relations. Rows~1--4 are exact for a finite system; the
others hold in the limit $V\to\infty$. $\varepsilon_*$ is the thermal energy,
and subscripts of $\lds$ denote partial derivatives.}
\label{tab:relations}
\renewcommand{\arraystretch}{1.5}
\begin{tabular}{p{0.26\linewidth}p{0.50\linewidth}l}
\toprule
Quantity & Relation & Eq. \\
\midrule
density of states, level count &
$\rho(\enrg)=\sum_a\delta(\enrg-\enrg_a)$,\quad $\cN(\enrg)=\int_{-\infty}^{\enrg}d\enrg'\,\rho(\enrg')$ & \eqref{eq:filt_gN} \\
partition function &
$\Xi=\int d\enrg\,\rho(\enrg)\,e^{-\beta\enrg}=\beta\int d\enrg\,\cN(\enrg)\,e^{-\beta\enrg}$ & \eqref{eq:filt_laplace} \\
entropy, heat capacity &
$S=(\langle\hat{\mathcal H}\rangle-\Omega)/T$,\quad $C_\mu=\Var(\hat{\mathcal H})/k_BT^{2}$ & \eqref{eq:filt_mb} \\
response to a field &
$\langle\hat A\rangle=k_BT\,\Xi^{(1)}/\Xi$,\quad $\partial_\hfield\rho=\partial_\enrg\rho_A$ & \eqref{eq:filt_param}, \eqref{eq:filt_gA} \\
\midrule
\lndos &
$\lds(\varepsilon)=\lim_{V\to\infty}(k_B/V)\ln\cN_V(V\varepsilon)$ & \eqref{eq:filt_sdef} \\
thermodynamic potential &
$\omega(T)=\min_\varepsilon[\varepsilon-T\lds(\varepsilon)]=\varepsilon_*-T\lds(\varepsilon_*)$ & \eqref{eq:filt_legendre} \\
temperature &
$\lds_\varepsilon(\varepsilon_*)=1/T$ & \eqref{eq:filt_stat} \\
energy, entropy &
$u-\mu n=\varepsilon_*$,\quad $S/V=\lds(\varepsilon_*)$ & \eqref{eq:filt_SI} \\
heat capacity &
$c_\mu=d\varepsilon_*/dT=-1/T^{2}\lds_{\varepsilon\varepsilon}$ & \eqref{eq:filt_curvmain} \\
conjugate density, susceptibility &
$a=T\lds_\hfield$,\quad $\chi=T\,[\lds_{\hfield\hfield}-\lds_{\varepsilon\hfield}^{2}/\lds_{\varepsilon\varepsilon}]$ & \eqref{eq:filt_af}, \eqref{eq:filt_chis} \\
susceptibility at wave vector $\mathbf q$ &
$\chi(\mathbf q)=2T\,\partial^{2}\lds/\partial \hfield_{\mathbf q}^{2}$ & \eqref{eq:filt_chiq} \\
\bottomrule
\end{tabular}
\end{table}

\begin{table}[htbp]
\centering
\small
\caption{Assumptions on the spectrum and their consequences.}
\label{tab:assumptions}
\begin{tabular}{p{0.34\linewidth}p{0.58\linewidth}}
\toprule
Assumption & Result \\
\midrule
none (any $\hat H$, finite size) &
$\Xi=\int d\enrg\,\rho\,e^{-\beta\enrg}$; $S$, $C_\mu$ and the responses from
$\Xi$ and $\partial_\hfield^{n}\rho$,
Eqs.~\eqref{eq:filt_mb}, \eqref{eq:filt_param} \\[2pt]
extensive $\ln\cN$ (Section~\ref{sec:large}) &
$\omega=\min_\varepsilon[\varepsilon-T\lds]$, $S/V=\lds(\varepsilon_*)$,
Eqs.~\eqref{eq:filt_legendre}, \eqref{eq:filt_SI} \\[2pt]
\ref{ass:power}: $\lds=\cA\,(\delta\varepsilon)^{\sigma}$ above the ground state &
$c_\mu\propto T^{\sigma/(1-\sigma)}$, $S/C_\mu=(1-\sigma)/\sigma$,
Eq.~\eqref{eq:filt_powerT} \\[2pt]
\ref{ass:hom0}: homogeneous in $(\delta\varepsilon,\hfield)$ above the ground state &
scaling function, quantum critical fan, Eq.~\eqref{eq:filt_scalingform} \\[2pt]
\ref{ass:crit}: homogeneous in $(\delta\varepsilon,\hfield)$ about $\varepsilon_c$ &
$\alpha$, $\tilde\beta$, $\gamma$, $\delta$ from $(\sigma,b)$,
Eqs.~\eqref{eq:filt_critexp}, \eqref{eq:filt_critexp2} \\[2pt]
\ref{ass:q}: homogeneous in the wave vector &
$\nu=b_q/(\sigma-1)$, $\eta$ and the Fisher relation,
Eqs.~\eqref{eq:filt_chiqscal}, \eqref{eq:filt_fisher} \\[2pt]
\ref{ass:V}: homogeneous in the size &
finite-size scaling, Eq.~\eqref{eq:filt_fss}; hyperscaling $2-\alpha=\nu_*$
if $b_V=\sigma$, Eq.~\eqref{eq:filt_hyper} \\
\bottomrule
\end{tabular}
\end{table}

\paragraph{Limits.}\label{sec:scope}
First, at a critical point the formulation is equivalent to the scaling
hypothesis and adds no relation among the exponents
(Section~\ref{sec:filt_rg}). Second, it relates the exponents to the
spectrum but does not compute the spectrum: $(\sigma,b,b_q,b_V)$ must be
obtained from the Hamiltonian, exactly or numerically, and the fact that
different Hamiltonians share them is explained by the renormalisation
group. Third, properties of the eigenvectors that leave the level count
unchanged are invisible to it: localisation, a ground-state degeneracy that
is not exponential in the size (topological order), and all dynamical
quantities. Fourth, the thermodynamic potential determines only the concave hull of
$\lds$; where $\lds$ is not concave (long-range interactions) the
microcanonical and the canonical descriptions differ~\cite{Campa2009}.

\paragraph{Outlook.} The assumptions are statements about a microscopic
object that can be computed directly, by exact diagonalisation or by
sampling methods that estimate the density of states~\cite{WangLandau2001}.
Three tests on spectra of finite lattice models follow from the text: the
vanishing of the curvature of $\lds_V$ at $\varepsilon_c$ with increasing
size; the dependence of the singular part of $\ln\cN_V$ on the combinations
$\delta\varepsilon\,V^{1/\sigma}$ and $\hfield\,V^{b/\sigma}$ alone
(Section~\ref{sec:filt_size}), which tests hyperscaling without a
correlation length; and the determination of $b_q$ from the level count in
a modulated field. On the experimental side the \lndos{} is obtained from
calorimetry through the parametric relation $\lds(u-\mu n)=S/V$, and the
amplitude-free ratios of Section~\ref{sec:filt_power_mb} test a power law
without a fit; in frustrated magnets the width $W$ of the low-lying group of
levels follows from the $T^{-2}$ tail of the heat capacity, and the laws
\eqref{eq:filt_manifold} and~\eqref{eq:filt_neardeg} relate the measured
exponent to the dimension of the degenerate manifold. Three extensions are
open: the joint density of states in the energy and a conserved density,
which gives the ensemble of fixed particle number and the distribution of
the order parameter; the derivation of the form of $\lds$ from the
Hamiltonian where the spectrum is known (integrable models, conformal field
theories); and the inclusion of the eigenvectors, through the dependence of
the diagonal matrix elements $A_{aa}$ on the energy that enters
Eq.~\eqref{eq:filt_af}, which is needed for the dynamics.

\clearpage
\appendix

\section{The minimum of a function near a singular point}
\label{sec:tools}

Every result of Sections~\ref{sec:transitions} and~\ref{sec:filt_lowenergy}
is obtained by one operation: the function
$\fen(\varepsilon)=\varepsilon-T\,\lds(\varepsilon)$ is minimised over the
energy, and the minimum value $\omega$ and the minimiser $\varepsilon_*$ are
differentiated with respect to $T$ and to the parameters of the
Hamiltonian. This appendix collects, in a form that does not refer to
thermodynamics, the elementary facts about such minima that the main text
uses. Throughout, $x$ is the distance from the singular point, $\tau$ the
small parameter that drives the minimiser towards it, $g(x;\tau)$ the
function to be minimised, $x_*(\tau)$ its minimiser and
\begin{equation}\label{eq:tools_m}
  m(\tau):=\min_{x}g(x;\tau)=g\big(x_*(\tau);\tau\big)
\end{equation}
its minimum value. Subscripts denote partial derivatives. The correspondence
with the quantities of the main text is given in Table~\ref{tab:tools}.

\subsection{Power laws and analyticity}
\label{sec:tools_power}

For two functions of $x\to0$, $f\simeq g$ means $f/g\to1$ (\emph{equality to
leading order}). A statement ``$\lds=\cA\,x^{\sigma}$ for $x\to0^{+}$'' is
meant to leading order, $\lds=\cA\,x^{\sigma}\,[1+o(1)]$; two functions that
agree to leading order give observables that agree to leading order
(footnote to Assumption~\ref{ass:power}).

The function $f(x)=x^{\sigma}$, $x>0$, is up to a constant factor the only
one that satisfies
\begin{equation}\label{eq:tools_hom1}
  f(\lambda x)=\lambda^{\sigma}f(x)\qquad\text{for all }\lambda>0
\end{equation}
(set $x=1$): a power law is the statement that the function has no intrinsic
scale. For $x$ of either sign the same holds for $B_\pm|x|^{\sigma}$, with
independent amplitudes for $x>0$ and $x<0$. The $k$th derivative of
$x^{\sigma}$ is $\propto x^{\sigma-k}$; for $x\to0^{+}$ it vanishes if
$k<\sigma$ and diverges if $k>\sigma$, unless $\sigma$ is an integer. Hence
$x^{\sigma}$ is analytic at $x=0$ only if $\sigma$ is a non-negative
integer, and $|x|^{\sigma}$ only if $\sigma$ is an even integer. Near a point
$x=0$ at which it is not analytic, every function considered in this note
has the local form
\begin{equation}\label{eq:tools_local}
  f(x)=\underbrace{f_0+f_1x+\tfrac12f_2x^{2}+\dots}_{\text{regular part}}
  +\underbrace{B_\pm|x|^{\sigma}\,[1+o(1)]}_{\text{singular part}} ,
\end{equation}
with $\sigma$ not an even integer (or $\sigma=2$ with $B_+\neq B_-$). The
singular term is \emph{subleading} to the regular term $f_kx^{k}$ if
$\sigma>k$ and $f_k\neq0$, and \emph{leading} otherwise.

If $g$ is analytic in $(x,\tau)$ near $(x_*,\tau)$ and
$g_{xx}(x_*;\tau)\neq0$, the stationarity condition $g_x(x_*;\tau)=0$ can be
solved for an analytic $x_*(\tau)$ by the implicit-function theorem, and
$m(\tau)$ is analytic. A non-analytic dependence of the minimum value on the
parameter therefore requires that $g$ be non-analytic at $x_*$, that
$g_{xx}=0$ at $x_*$, or that the minimiser be not unique. These are the
cases of Section~\ref{sec:transitions}.

\subsection{The minimum of two competing powers}
\label{sec:tools_min}

\begin{lemma}[Balance of two powers]\label{lem:balance}
Let $\lambda,\Lambda>0$ and $0<p<q$. The function
\begin{equation}\label{eq:tools_g2}
  g(x)=-\lambda\,x^{p}+\Lambda\,x^{q} , \qquad x\ge0 ,
\end{equation}
has a unique minimum, at
\begin{equation}\label{eq:tools_xstar}
  x_*=\left(\frac{p\,\lambda}{q\,\Lambda}\right)^{1/(q-p)} ,
  \qquad
  m=g(x_*)=-\frac{q-p}{q}\,\lambda\,x_*^{\,p}
  \propto-\lambda^{\,q/(q-p)} .
\end{equation}
\end{lemma}
\begin{proof}
$g'(x)=x^{p-1}\big(-p\lambda+q\Lambda x^{q-p}\big)$ vanishes for $x>0$ only at
$x_*$, and is negative below and positive above. There
$\Lambda x_*^{q}=(p/q)\lambda x_*^{p}$, which gives $m$.
\end{proof}
As $\lambda\to0$ the minimiser moves to the singular point as
$\lambda^{1/(q-p)}$, and the minimum value vanishes as $\lambda^{q/(q-p)}$;
both exponents are fixed by the difference of the two exponents of $g$. Two
instances cover the applications.
\begin{itemize}\setlength{\itemsep}{2pt}
\item \emph{Power against a linear term} ($p=\sigma<1$, $q=1$, $\Lambda=1$):
  $g=x-\lambda x^{\sigma}$, with
  \begin{equation}\label{eq:tools_edge}
    x_*=(\sigma\lambda)^{1/(1-\sigma)} ,\qquad
    m=-\frac{1-\sigma}{\sigma}\,x_*\propto-\lambda^{1/(1-\sigma)} .
  \end{equation}
  This is the low-temperature case, $\lambda=T\cA$. For $\sigma\ge1$ the
  minimum stays at $x=0$ for small $\lambda$.
\item \emph{Linear term against a power} ($p=1$, $q=\sigma>1$): for $x$ of
  either sign, $g=-\lambda x+\Lambda_\pm|x|^{\sigma}$ has its minimum on the
  side of the sign of $\lambda$, with the amplitude of that side,
  \begin{equation}\label{eq:tools_crit}
    x_*=\pm\left(\frac{|\lambda|}{\sigma\Lambda_\pm}\right)^{1/(\sigma-1)} ,\qquad
    m=-\frac{\sigma-1}{\sigma}\,|\lambda|\,|x_*|\propto-|\lambda|^{\sigma/(\sigma-1)} .
  \end{equation}
  This is the critical case, $\lambda=t$. The exponents $\sigma$ of $g$ and
  $\sigma'=\sigma/(\sigma-1)$ of $m$ are conjugate,
  $1/\sigma+1/\sigma'=1$: the Legendre transform of $|x|^{\sigma}/\sigma$ is
  $|\lambda|^{\sigma'}/\sigma'$.
\end{itemize}
The lemma locates the minimum of the local form~\eqref{eq:tools_local},
which is valid for small $x$ only. It is the global minimum if, as
$\lambda\to0$, the full function stays positive for $|x|\ge x_1$, while the
local minimum value is negative; this is how the proof of
Proposition~\ref{prop:power} uses the bound~\eqref{eq:filt_linbound}. When
the condition fails, two separate minima can exchange their roles: a
discontinuous transition.

\subsection{Derivatives of a minimum value}
\label{sec:tools_env}

\begin{lemma}[Envelope theorem]\label{lem:envelope}
Let $g(x;\tau)$ be twice differentiable, and let $x_*(\tau)$ be a minimiser
in the interior of the range of $x$, with $g_{xx}(x_*;\tau)>0$. Then
\begin{equation}\label{eq:tools_env}
  \frac{dm}{d\tau}=g_\tau\Big|_{x_*} ,\qquad
  \frac{dx_*}{d\tau}=-\frac{g_{x\tau}}{g_{xx}}\Big|_{x_*} ,\qquad
  \frac{d^{2}m}{d\tau^{2}}=\left[g_{\tau\tau}-\frac{g_{x\tau}^{2}}{g_{xx}}\right]_{x_*} ,
\end{equation}
and likewise for a second parameter.
\end{lemma}
\begin{proof}
$dm/d\tau=g_x\,dx_*/d\tau+g_\tau$, and $g_x=0$ at $x_*$. Differentiating
$g_x(x_*(\tau);\tau)=0$ gives $g_{xx}\,dx_*/d\tau+g_{x\tau}=0$, and
differentiating $dm/d\tau=g_\tau(x_*(\tau);\tau)$ once more gives
$d^{2}m/d\tau^{2}=g_{\tau x}\,dx_*/d\tau+g_{\tau\tau}$.
\end{proof}
A parameter enters the first derivative of the minimum value only through
its explicit presence in $g$; the minimiser responds to the parameter in
inverse proportion to the curvature of $g$ at the minimum. With $g=\fen$ the
lemma gives Eqs.~\eqref{eq:filt_SI}, \eqref{eq:filt_curvmain},
\eqref{eq:filt_af} and~\eqref{eq:filt_chis}. Applied to a small term,
$g=g_0+\eta\,g_1$ with $x_0$ the minimiser of $g_0$, it gives
\begin{equation}\label{eq:tools_pert}
  m=m_0+\eta\,g_1(x_0)+O(\eta^{2}) :
\end{equation}
to first order the minimum value changes by the perturbation evaluated at
the unperturbed minimiser.

\begin{lemma}[Subleading singular term]\label{lem:sublead}
Let $\Lambda_2>0$, $\sigma>2$ and
\begin{equation}\label{eq:tools_gsub}
  g(x;\tau)=-\tau\,x+\Lambda_2\,x^{2}+\Lambda_\pm|x|^{\sigma} .
\end{equation}
Then, for $\tau\to0$,
\begin{equation}\label{eq:tools_msub}
  x_*=\frac{\tau}{2\Lambda_2}\,\big[1+O(|\tau|^{\sigma-2})\big] ,\qquad
  m(\tau)=-\frac{\tau^{2}}{4\Lambda_2}
  +\Lambda_\pm\Big|\frac{\tau}{2\Lambda_2}\Big|^{\sigma}\big[1+O(|\tau|^{\sigma-2})\big] .
\end{equation}
\end{lemma}
\begin{proof}
Eq.~\eqref{eq:tools_pert} with $g_0=-\tau x+\Lambda_2x^{2}$,
$x_0=\tau/2\Lambda_2$, $m_0=-\tau^{2}/4\Lambda_2$ and
$\eta\,g_1=\Lambda_\pm|x|^{\sigma}$; the term of second order in
$\Lambda_\pm$ is smaller than the first-order one by a factor
$O(|x_0|^{\sigma-2})$.
\end{proof}
When the singular term is leading, Eq.~\eqref{eq:tools_crit}, the minimiser
scales as $|\tau|^{1/(\sigma-1)}$ and the minimum value as
$|\tau|^{\sigma/(\sigma-1)}$: the exponent is transformed. When it is
subleading to a quadratic term, the minimiser is fixed by the quadratic
term and the singular part of the minimum value is the singular term itself
at $x_*$: the exponent is transmitted unchanged. Both cases occur at
critical points (Section~\ref{sec:filt_critical}). The regimes are collected
in Table~\ref{tab:tools_regimes}.
\begin{table}[htbp]
\centering
\small
\caption{The minimum of $g(x;\tau)$ near a singular term
$\Lambda|x|^{\sigma}$ of $g$ at $x=0$, for $\tau\to0$. The last column is the
second derivative of the minimum value, which gives the heat capacity in the
applications.}
\label{tab:tools_regimes}
\begin{tabular}{lllll}
\toprule
exponent & $g(x;\tau)$ & $x_*$ & singular part of $m$ & $d^{2}m/d\tau^{2}$ \\
\midrule
$0<\sigma<1$ & $x-\tau\Lambda x^{\sigma}$ & $\propto\tau^{1/(1-\sigma)}$ & $\propto-\tau^{1/(1-\sigma)}$ & $\propto\tau^{(2\sigma-1)/(1-\sigma)}$ \\[2pt]
$1<\sigma<2$ & $-\tau x+\Lambda|x|^{\sigma}$ & $\propto|\tau|^{1/(\sigma-1)}$ & $\propto-|\tau|^{\sigma/(\sigma-1)}$ & vanishes as $|\tau|^{(2-\sigma)/(\sigma-1)}$ \\[2pt]
$\sigma>2$ & $-\tau x+\Lambda|x|^{\sigma}$ & $\propto|\tau|^{1/(\sigma-1)}$ & $\propto-|\tau|^{\sigma/(\sigma-1)}$ & diverges as $|\tau|^{-(\sigma-2)/(\sigma-1)}$ \\[2pt]
$\sigma>2$ & $-\tau x+\Lambda_2x^{2}+\Lambda|x|^{\sigma}$ & $\simeq\tau/2\Lambda_2$ & $\simeq\Lambda|\tau/2\Lambda_2|^{\sigma}$ & finite; singular part $\propto|\tau|^{\sigma-2}$ \\
\bottomrule
\end{tabular}
\end{table}

\subsection{Scaling functions}
\label{sec:tools_scaling}

\begin{lemma}[Minimisation and rescaling]\label{lem:rescale}
Let $g$ depend on a second parameter $h$ and satisfy, for all $\lambda>0$
and fixed exponents $a$, $b$, $c$,
\begin{equation}\label{eq:tools_ghom}
  g(\lambda x;\lambda^{a}\tau,\lambda^{b}h)=\lambda^{c}\,g(x;\tau,h) .
\end{equation}
Then the minimum value and the minimiser satisfy
\begin{equation}\label{eq:tools_mhom}
  m(\lambda^{a}\tau,\lambda^{b}h)=\lambda^{c}\,m(\tau,h) ,\qquad
  x_*(\lambda^{a}\tau,\lambda^{b}h)=\lambda\,x_*(\tau,h) .
\end{equation}
\end{lemma}
\begin{proof}
Substitute $x=\lambda x'$ in the minimisation:
$m(\lambda^{a}\tau,\lambda^{b}h)=\min_{x'}g(\lambda x';\lambda^{a}\tau,\lambda^{b}h)
=\lambda^{c}\min_{x'}g(x';\tau,h)$, and the minimiser of the right-hand side
is $x'=x_*(\tau,h)$.
\end{proof}
The homogeneity of $g$ is inherited by the minimum value, with the same
degree $c$ and the same exponents of the parameters; further parameters are
treated in the same way. The function $g=x-\tau\,\lds(x,h)$ with $\lds$ as in
Eq.~\eqref{eq:filt_shom0} satisfies Eq.~\eqref{eq:tools_ghom} with
$a=1-\sigma$, $c=1$, and $g=-\tau x-T_c\lds_{\rm sing}(x,h)$ with
Eq.~\eqref{eq:filt_shom} satisfies it with $a=\sigma-1$, $c=\sigma$.

\paragraph{Scaling form.} The choice $\lambda=|\tau|^{-1/a}$ in
Eq.~\eqref{eq:tools_mhom} gives
\begin{equation}\label{eq:tools_scalingform}
  m(\tau,h)=|\tau|^{c/a}\,\Phi_\pm(z) ,\qquad
  z:=\frac{h}{|\tau|^{b/a}} ,\qquad
  \Phi_\pm(z):=m(\pm1,z) ,
\end{equation}
with one function for $\tau>0$ and one for $\tau<0$: the function of two
variables reduces to a power of $\tau$ times a function of one ratio.

\paragraph{Exponents by differentiation.} Each derivative of
Eq.~\eqref{eq:tools_mhom} with respect to $\tau$ or $h$ lowers the degree by
the exponent of that variable. Hence
\begin{equation}\label{eq:tools_expo}
  \frac{\partial^{\,j+k}m}{\partial\tau^{j}\,\partial h^{k}}\bigg|_{h=0}
  \propto|\tau|^{(c-ja-kb)/a} ,
  \qquad
  \frac{\partial^{k}m}{\partial h^{k}}\bigg|_{\tau=0}
  \propto|h|^{(c-kb)/b} ,
\end{equation}
provided that the corresponding derivative of the scaling function is finite
and non-zero. Every exponent of a derivative of $m$ is a linear function of
the two ratios $c/a$ and $b/a$. Four exponents therefore satisfy two linear
relations identically, whatever the values of $\sigma$ and $b$: the scaling
relations~\eqref{eq:rushbrooke} test the homogeneity, not the values of its
exponents.

\paragraph{Scaling function at large argument.} If $m(0,h)$ is finite and
non-zero for $h\neq0$, the factor $|\tau|^{c/a}=|h|^{c/b}\,|z|^{-c/b}$ in
Eq.~\eqref{eq:tools_scalingform} must be compensated:
\begin{equation}\label{eq:tools_match}
  \Phi_\pm(z)\simeq\Phi_\infty\,|z|^{c/b}\quad(|z|\to\infty) ,
  \qquad m(0,h)=\Phi_\infty|h|^{c/b} .
\end{equation}
This argument gives $\delta$ in Proposition~\ref{prop:crit}, the Fisher
relation in Section~\ref{sec:filt_length} (with $q$ in the role of $h$) and
the finite-size behaviour of the susceptibility in
Section~\ref{sec:filt_size} (with $V^{-1}$ in the role of $h$).

\subsection{Logarithmic corrections}
\label{sec:tools_logs}

A singular term $\Lambda|x|^{\sigma}L(x)$ with a \emph{slowly varying}
factor, $L(\lambda x)/L(x)\to1$ for $x\to0$ at every fixed $\lambda$ (a power
of $\ln(1/x)$, for instance), satisfies Eq.~\eqref{eq:tools_hom1} only in the
limit. Lemma~\ref{lem:balance} then holds with $\Lambda$ replaced by
$\Lambda L(x_*)$, and in $L(x_*)$ the argument may be replaced by its
power-law estimate. For $L=\ln^{\theta}(1/x)$ this gives
$|x_*|\propto\big[|\lambda|/\ln^{\theta}(1/|\lambda|)\big]^{1/(\sigma-1)}$ and
$m\propto-|\lambda|\,|x_*|$. The case $\sigma=2$, $L=1/\ln(1/x)$ gives
$x_*\propto\lambda\ln(1/|\lambda|)$ and a second derivative of $m$ that
diverges as $\ln(1/|\lambda|)$: the logarithmic heat capacity of
Section~\ref{sec:filt_nonpower}. An essential singularity,
$e^{-1/|x|^{\theta}}$, has no exponent, and the minimum value inherits an
essential singularity.

\begin{table}[htbp]
\centering
\small
\caption{The variables of Appendix~\ref{sec:tools} in the applications. In
each case $g$ is $\fen=\varepsilon-T\lds$ and $m$ is $\omega$, up to terms
independent of $\varepsilon$.}
\label{tab:tools}
\begin{tabular}{lllll}
\toprule
Application & $x$ & $\tau$ & $g(x;\tau,h)$ & Lemma \\
\midrule
low $T$ (Sec.~\ref{sec:filt_power_mb}) & $\varepsilon-\varepsilon_0$ & $T$ & $x-T\cA x^{\sigma}$ & \ref{lem:balance} \\[2pt]
scaling at low $T$ (Sec.~\ref{sec:filt_scaling0}) & $\varepsilon-\varepsilon_0$ & $T$ & $x-T\lds(x,\hfield)$ & \ref{lem:rescale} \\[2pt]
critical, $\alpha>0$ (Sec.~\ref{sec:filt_critical}) & $\varepsilon-\varepsilon_c$ & $t$ & $-tx-T_c\lds_{\rm sing}(x,\hfield)$ & \ref{lem:balance}, \ref{lem:rescale} \\[2pt]
critical, $\alpha<0$ (Sec.~\ref{sec:filt_critical}) & $\varepsilon-\varepsilon_c$ & $t$ & $-tx+\dfrac{x^{2}}{2T_cc_{\rm reg}}-T_c\lds_{\rm sing}$ & \ref{lem:sublead} \\[4pt]
wave vector, size (Secs.~\ref{sec:filt_length}, \ref{sec:filt_size}) & $\varepsilon-\varepsilon_c$ & $t$ & as above, $q$ or $V^{-1}$ a further parameter & \ref{lem:rescale} \\[2pt]
responses (Sec.~\ref{sec:large_laplace}) & $\varepsilon$ & $T$, $\hfield$ & $\varepsilon-T\lds(\varepsilon(\hfield))$ & \ref{lem:envelope} \\
\bottomrule
\end{tabular}
\end{table}

\section{Static response and the motion of the levels}
\label{app:lehmann}\label{app:continuity}

This appendix proves Eq.~\eqref{eq:filt_gA} and derives the static response
to a field $ \hfield $ coupled to a Hermitian operator $ \hat A $,
$ \hat{\mathcal{H}}(\hfield) = \hat{\mathcal{H}} - \hfield\hat A $, in the
notation of Section~\ref{sec:filt_resp}.

\paragraph{Derivative of a trace.} For a Fock space of finite dimension and
every continuously differentiable function $\varphi$,
\begin{equation}\label{eq:cont_trace}
  \frac{\partial}{\partial\hfield}\Tr\varphi\big(\hat{\mathcal H}(\hfield)\big)
  =-\Tr\big[\hat A\,\varphi'\big(\hat{\mathcal H}(\hfield)\big)\big] ,
\end{equation}
without any assumption on the commutator $[\hat A,\hat{\mathcal H}]$, on
degeneracies or on crossings of the levels. For $\varphi(x)=x^{n}$ the
product rule gives
$\partial_\hfield\hat{\mathcal H}^{n}=-\sum_{k=0}^{n-1}\hat{\mathcal H}^{k}\hat A\,\hat{\mathcal H}^{\,n-1-k}$,
and under the trace each term can be cycled to
$-\Tr[\hat A\,\hat{\mathcal H}^{\,n-1}]$. By linearity
Eq.~\eqref{eq:cont_trace} holds for polynomials, and a continuously
differentiable $\varphi$ is approximated, together with its derivative, by
polynomials uniformly on the bounded interval that contains the spectra. On
the right of Eq.~\eqref{eq:cont_trace},
$\varphi'(\hat{\mathcal H})=\sum_a\varphi'(\enrg_a)|a\rangle\langle a|$ is
diagonal in the eigenbasis, so that only the diagonal elements of $\hat A$
enter:
$\partial_\hfield\langle\rho,\varphi\rangle=-\sum_aA_{aa}\varphi'(\enrg_a)
=-\langle\rho_A,\varphi'\rangle=\langle\partial_\enrg\rho_A,\varphi\rangle$,
which is Eq.~\eqref{eq:filt_gA}. With $\varphi\equiv1$ the total number of
levels is independent of $\hfield$, and with $\varphi=e_\beta$ (for which the
identity holds also in an infinite-dimensional Fock space, when both traces
converge) $\partial_\hfield\Xi=\beta\,\Xi\,\langle\hat A\rangle$.

\paragraph{Slope and curvature of the levels.} Where a level is a
differentiable function of the field,
$ \partial\enrg_a/\partial \hfield = -A_{aa} $ (the Hellmann--Feynman
theorem\footnote{For a normalised eigenvector $ |a(\hfield)\rangle $ of
$ \hat{\mathcal{H}}(\hfield) $,
$ \partial\enrg_a/\partial \hfield = \langle a|\partial\hat{\mathcal{H}}/\partial \hfield|a\rangle $:
the terms containing the derivative of the eigenvector add up to
$ \enrg_a\,\partial\langle a|a\rangle/\partial \hfield = 0 $. Within a
degenerate level the basis is chosen to diagonalise $\hat A$.}). The
derivatives of the delta functions in $\rho$ give
$ \Xi^{(1)} = \beta\sum_aA_{aa}\,e^{-\beta\enrg_a} $ and
$ \Xi^{(2)} = \sum_a\big[\beta^{2}A_{aa}^{2} - \beta\,\partial^{2}\enrg_a/\partial \hfield^{2}\big]e^{-\beta\enrg_a} $,
so that Eq.~\eqref{eq:filt_param} becomes
\begin{equation}\label{eq:filt_vanvleck}
  \langle\hat A\rangle = \sum_ap_a\,A_{aa} , \qquad
  \chi_{AA} = \frac{1}{k_BT}\Big[\sum_ap_a\,A_{aa}^{2} - \Big(\sum_ap_a\,A_{aa}\Big)^{2}\Big]
  - \sum_ap_a\,\frac{\partial^{2}\enrg_a}{\partial \hfield^{2}} .
\end{equation}
The first term of $ \chi_{AA} $ is the redistribution of the populations
over levels of different slope (the Curie term); the second is the
curvature of the levels as functions of the field (the Van Vleck term),
with
$ -\partial^{2}\enrg_a/\partial \hfield^{2} = 2\sum_{b:\,\enrg_b\neq\enrg_a}|A_{ab}|^{2}/(\enrg_b - \enrg_a) $
from second-order perturbation theory and
$ A_{ab} := \langle a|\hat A|b\rangle $. Eq.~\eqref{eq:filt_corr} is the same
result for two fields.

\paragraph{Conserved quantities.} If $ \hat A $ commutes with $ \hat H $ and
$ \hat N $, the eigenbasis may be chosen to diagonalise it,
$ \hat A|a\rangle = A_a|a\rangle $; the field shifts the eigenvalues to
$ \enrg_a - \hfield\,A_a $ without changing the eigenvectors, and the Van
Vleck term vanishes. For any $ \hat B $ diagonal in the same basis, the
derivative of $\langle\hat B\rangle=\sum_aB_aw_a/\Xi$ with the weights
$w_a=e^{-\beta(\enrg_a-\hfield A_a)}$, for which
$\partial w_a/\partial\hfield=\beta A_aw_a$ and
$\partial w_a/\partial T=(\enrg_a/k_BT^{2})\,w_a$, is
\begin{equation}\label{eq:filt_mb_response}
  \frac{\partial\langle\hat B\rangle}{\partial \hfield}\bigg|_{T,\mu}
  = \frac{\Cov(\hat A, \hat B)}{k_BT} ,
  \qquad
  \frac{\partial\langle\hat B\rangle}{\partial T}\bigg|_{\mu}
  = \frac{\Cov(\hat{\mathcal{H}}, \hat B)}{k_BT^{2}} .
\end{equation}
Instances are $ (\partial N/\partial\mu)_T = \Var(\hat N)/k_BT $, the spin
susceptibility $ \Var(\hat M)/Vk_BT $ whenever the total moment $ \hat M $
is conserved, and $ C_\mu $ itself, Eq.~\eqref{eq:filt_CU}.

\paragraph{Non-conserved quantities.} If $ \hat A $ is not conserved, the
field changes the eigenvectors as well. To first order in $ \hfield $,
\[
  |a(\hfield)\rangle = |a\rangle - \hfield\sum_{b\,:\,\enrg_b\neq\enrg_a}|b\rangle\,\frac{A_{ba}}{\enrg_a - \enrg_b} .
\]
In $ \langle\hat B\rangle = \sum_ap_a(\hfield)\,\langle a(\hfield)|\hat B|a(\hfield)\rangle $
the populations contribute
$ \beta\sum_ap_aA_{aa}B_{aa} - \beta\langle\hat A\rangle\langle\hat B\rangle $
to the derivative at $ \hfield = 0 $, and the eigenvectors
$ -\sum_{a,b:\,\enrg_a\neq\enrg_b}p_a(A_{ab}B_{ba} + B_{ab}A_{ba})/(\enrg_a - \enrg_b) $.
Exchanging $ a $ and $ b $ in the second term of the latter gives the
\emph{Lehmann} (spectral) representation of the isothermal susceptibility,
\begin{equation}\label{eq:filt_mb_kubo}
  \chi_{AB} := \frac{\partial\langle\hat B\rangle}{\partial \hfield}\bigg|_{\hfield=0}
  = \sum_{a,b}A_{ab}B_{ba}\,w_{ab} - \beta\langle\hat A\rangle\langle\hat B\rangle ,
  \qquad
  w_{ab} := \frac{p_b - p_a}{\enrg_a - \enrg_b} ,
\end{equation}
with $ w_{ab} := \beta p_a $ when $ \enrg_a = \enrg_b $ (the limit of the
quotient). The weight $ w_{ab} > 0 $ acts on pairs of eigenstates connected
by the perturbation. With $ y = \beta(\enrg_a - \enrg_b) $,
$ w_{ab} = \beta p_b(1 - e^{-y})/y $, while
$ \tfrac12\beta(p_a + p_b) = \tfrac12\beta p_b(1 + e^{-y}) $, and
$ (1 - e^{-y})/y \leq \tfrac12(1 + e^{-y}) $ is the inequality
$ \tanh(y/2) \leq y/2 $. Hence, with
$ \delta\hat A = \hat A - \langle\hat A\rangle $,
\begin{equation}\label{eq:filt_chibound}
  0 \;\leq\; \chi_{AA} = \sum_{a,b}\big|\langle a|\delta\hat A|b\rangle\big|^{2}\,w_{ab}
  \;\leq\; \frac{\Var(\hat A)}{k_BT} ,
\end{equation}
with equality on the right if and only if $ \hat A $ connects only states
of equal energy, i.e.\ is conserved. A conserved quantity responds exactly
as much as it fluctuates; a non-conserved one responds less.

\section{The partition function as a generating function}
\label{app:generating}

\paragraph{Cumulants and large deviations.} For a Fock space of finite
dimension $D$, $\rho/D$ is the probability density of the energy when all
eigenstates have equal weight ($T=\infty$), and
$\Xi(\beta)/D=\langle e^{-\beta\enrg}\rangle_\infty$ is its moment-generating
function: $\ln\Xi$ generates the cumulants of the
energy~\cite[Sec.~4.3]{Honerkamp},
$\partial^{n}\ln\Xi/\partial\beta^{n}=(-1)^{n}\kappa_n$ with $\kappa_n$ the
$n$th cumulant under $P_T$ (Eq.~\eqref{eq:filt_moments} for $n=1,2$). A
change of temperature is an exponential reweighting of one and the same
distribution $\rho$. In these terms Section~\ref{sec:large} is an instance
of the theory of large deviations~\cite{Touchette2009}: up to a constant,
$-\lds/k_B$ is the rate function of the energy per unit size, $-\beta\omega$
is the scaled cumulant-generating function, and
Eq.~\eqref{eq:filt_legendre} is the Legendre transform that relates them.

\paragraph{Inversion and resolution.} $\Xi(\beta)$ is analytic in the
half-plane ${\rm Re}\,\beta>0$. On the imaginary axis, $\beta=it/\hbar$,
$\Xi=\Tr e^{-i\hat{\mathcal H}t/\hbar}$ is the trace of the evolution
operator and the Fourier transform of $\rho$, which is inverted stably: a
structure of width $\Delta$ in energy is resolved from times up to
$\hbar/\Delta$ (spectroscopy). Thermodynamics has access to $\Xi$ on the real
$\beta$ axis only. The inverse Laplace transform,
\[
  \rho(\enrg)=\frac{1}{2\pi i}\int_{\beta_0-i\infty}^{\beta_0+i\infty}
  d\beta\;\Xi(\beta)\,e^{\beta\enrg} ,
\]
runs parallel to the imaginary axis and requires the analytic continuation
of $\Xi$ from the real axis, which is unique but unstable against small
errors of $\Xi$: the Laplace transform is injective and yet of finite
resolution (Section~\ref{sec:filt_dist}). For $V\to\infty$ the inversion
integral is dominated by the saddle point on the real axis at which
$\enrg=-\partial\ln\Xi/\partial\beta$; there
$\ln\rho\asymp\beta\enrg+\ln\Xi$, i.e.\
$S=(\langle\hat{\mathcal H}\rangle-\Omega)/T$, the Legendre transform inverse
to Eq.~\eqref{eq:filt_legendre}.

\paragraph{Zeros.} For a finite Fock space $\Xi(\beta)$ is a finite sum of
exponentials. It is positive on the real axis, and its zeros lie at complex
$\beta$. A singularity of $\omega$ at a real temperature can arise for
$V\to\infty$ only where zeros approach the real axis~\cite{Fisher1965}.

%% ---------------------------------------------------------------------------
%% Appendix D is a self-contained suggestion. It is not referenced from any other
%% part of the text and can be deleted as a whole (down to the line "END OF
%% APPENDIX D"), together with the bibliography entries NejatiLenk2026 and
%% FerrenbergSwendsen1989, which are cited only here.
%% ---------------------------------------------------------------------------
\section{Suggested test: the Falicov--Kimball model on the triangular lattice}
\label{app:fkm}

This appendix outlines a test of Section~\ref{sec:frustration} on a model
whose spectrum consists of families of levels exactly. The test has not been
carried out. Eqs.~\eqref{eq:fkm_levels}--\eqref{eq:fkm_hist} are exact
properties of the spectrum; the three statements at the end are predictions.

\paragraph{Model.} Itinerant spinless fermions (operators $\hat c_i$) hop
between neighbouring sites of a triangular lattice of $V$ sites and repel
localised particles, whose occupation $f_i\in\{0,1\}$ of the site $i$ is
conserved. The two particle numbers are conserved separately, each with its
chemical potential:
\begin{equation}\label{eq:fkm_H}
  \hat{\mathcal H}=-t\sum_{\langle ij\rangle}\big(\hat c_i^{\dagger}\hat c_j+{\rm h.c.}\big)
  +U\sum_i f_i\,\hat c_i^{\dagger}\hat c_i
  -\mu_c\sum_i\hat c_i^{\dagger}\hat c_i-\mu_f\sum_i f_i .
\end{equation}
A Monte Carlo study~\cite{NejatiLenk2026} finds a charge-density wave at low
temperature and, between it and the weakly correlated regime, a regime in
which several ordering wave vectors compete and the criteria of universal
scaling fail; the order of the transition was not determined. No order
parameter characterises the intermediate regime.

\paragraph{Families of levels.} For a fixed configuration
$\mathbf f=(f_1,\dots,f_V)$ the itinerant fermions are independent. Let
$E_k[\mathbf f]$, $k=1,\dots,V$, be the single-particle energies of the
hopping term in the potential $Uf_i$, and $n_k\in\{0,1\}$ their occupations.
The many-body levels are
\begin{equation}\label{eq:fkm_levels}
  \enrg[\mathbf f,\{n_k\}]=-\mu_f\sum_if_i+\sum_k\big(E_k[\mathbf f]-\mu_c\big)\,n_k .
\end{equation}
Each configuration $\mathbf f$ defines a family of $2^{V}$ levels in the sense
of Section~\ref{sec:frustration}: the families are disjoint sets of
eigenstates, so that the level counts and the partition functions add,
\begin{equation}\label{eq:fkm_sum}
  \cN_V(\enrg)=\sum_{\mathbf f}\cN_{\mathbf f}(\enrg) ,\qquad
  \Xi=\sum_{\mathbf f}\Xi_{\mathbf f} ,\qquad
  \Xi_{\mathbf f}=e^{\beta\mu_f\sum_if_i}\prod_k\Big(1+e^{-\beta(E_k[\mathbf f]-\mu_c)}\Big) .
\end{equation}
The equilibrium probability of a family is
$P_T(\mathbf f)=\Xi_{\mathbf f}/\Xi$; this is the weight sampled by a Monte
Carlo simulation over the configurations. The number of families, $2^{V}$,
is exponential in the size. Let the families be grouped into classes
$\mathfrak c$ of equal intensive properties (density and correlations of the
$f_i$), with $e^{Vs_{\rm conf}(\mathfrak c)/k_B}$ members, and let the \lndos{}
$\lds_{\mathfrak c}(\varepsilon)$ of the levels~\eqref{eq:fkm_levels} be the
same for almost all members of a class. The sum~\eqref{eq:fkm_sum} is then
dominated by its largest term,
\begin{equation}\label{eq:fkm_s}
  \lds(\varepsilon)=\max_{\mathfrak c}\big[s_{\rm conf}(\mathfrak c)+\lds_{\mathfrak c}(\varepsilon)\big] ,
  \qquad
  \omega(T)=\min_{\mathfrak c}\Big\{\min_\varepsilon\big[\varepsilon-T\lds_{\mathfrak c}(\varepsilon)\big]-T\,s_{\rm conf}(\mathfrak c)\Big\} ,
\end{equation}
which is Eq.~\eqref{eq:filt_compete} with the entropy of the class added. The
\lndos{} ranges from $0$ to $2k_B\ln2$, the logarithm of the total number
$4^{V}$ of levels per unit size.

\paragraph{Heat capacity within and between the families.} Let
$\bar n_k=1/(e^{\beta(E_k[\mathbf f]-\mu_c)}+1)$ be the mean occupation of
the single-particle level $k$. The mean and the variance of the energy
within the family $\mathbf f$ are
\[
  \bar\enrg_{\mathbf f}=-\mu_f\sum_if_i+\sum_k\big(E_k[\mathbf f]-\mu_c\big)\,\bar n_k ,
  \qquad
  v_{\mathbf f}=\sum_k\big(E_k[\mathbf f]-\mu_c\big)^{2}\,\bar n_k(1-\bar n_k) ,
\]
since the occupations are independent, each with the variance
$\bar n_k(1-\bar n_k)$. The variance of the energy over all levels is the
mean of the variances within the families plus the variance of the means,
\[
  \Var(\hat{\mathcal H})
  =\sum_{\mathbf f}P_T(\mathbf f)\,v_{\mathbf f}
  +\Big[\sum_{\mathbf f}P_T(\mathbf f)\,\bar\enrg_{\mathbf f}^{2}
  -\Big(\sum_{\mathbf f}P_T(\mathbf f)\,\bar\enrg_{\mathbf f}\Big)^{2}\Big] ,
\]
so that Eq.~\eqref{eq:filt_mb} becomes, at fixed $\mu_c$ and $\mu_f$,
\begin{equation}\label{eq:fkm_C}
  C_\mu=C_{\rm in}+C_{\rm ex} ,\qquad
  C_{\rm in}=\frac{1}{k_BT^{2}}\sum_{\mathbf f}P_T(\mathbf f)\,v_{\mathbf f} ,\qquad
  C_{\rm ex}=\frac{1}{k_BT^{2}}\,\Var_{P_T}\big(\bar\enrg_{\mathbf f}\big) .
\end{equation}
$C_{\rm in}$ is the heat capacity of the levels within the families, and
$C_{\rm ex}$ that of the redistribution of the weight among families of
different energy. Both are averages over the sampled configurations. Since
$\bar n_k(1-\bar n_k)=1/4\cosh^{2}x$ with $x=(E_k-\mu_c)/2k_BT$,
$C_{\rm in}/V$ is the fermionic heat capacity of Eq.~\eqref{eq:filt_cF} with
the single-particle density of states averaged over the families,
$\bar\varrho_T(E):=V^{-1}\sum_{\mathbf f}P_T(\mathbf f)\sum_k\delta(E-E_k[\mathbf f])$.

\paragraph{The \lndos{} from the distribution of the energy.} By
Section~\ref{sec:filt_exact} the distribution of the energy is
$P_T(\enrg)=\rho(\enrg)\,e^{-\beta\enrg}/\Xi$. With
$\rho_V(V\varepsilon)\asymp e^{V\lds_V(\varepsilon)/k_B}$ and
$\ln\Xi=-V\omega/k_BT$,
\begin{equation}\label{eq:fkm_hist}
  \lds_V(\varepsilon)=\frac{\varepsilon-\omega(T)}{T}+\frac{k_B}{V}\ln P_T(V\varepsilon)
\end{equation}
to logarithmic accuracy, for the energies that occur at the temperature $T$.
In the simulation $P_T(\enrg)$ is the average over the sampled families of
the distribution within a family, a Gaussian of mean $\bar\enrg_{\mathbf f}$
and variance $v_{\mathbf f}$ for large $V$. Histograms at neighbouring
temperatures overlap and determine $\lds_V$ over the whole range of energies
up to one additive constant (the multiple-histogram
method~\cite{FerrenbergSwendsen1989}), which is fixed by the value
$2k_B\ln2$ at the upper end or by the exact enumeration of a small lattice.

\paragraph{Predictions.}
\begin{itemize}\setlength{\itemsep}{2pt}
\item \emph{Order of the transition.} By Section~\ref{sec:transitions}, at a
  continuous transition the curvature of $\lds$ vanishes at a single energy,
  and at a discontinuous one $\lds$ is linear on the interval
  $[\varepsilon_-,\varepsilon_+]$. At a finite size the latter appears as an
  interval of positive curvature of $\lds_V$, since the levels of the
  phase-separated states are reduced in number by the interface, and as two
  maxima of $P_T(\enrg)$. The criterion is the dependence on the size: the
  width of the interval tends to the latent heat at a discontinuous
  transition and to zero at a continuous one, where such an interval can
  also occur at finite size~\cite{Behringer2006}. For classes of families
  with different ordering wave vectors and nearly equal energies,
  Section~\ref{sec:frustration} gives discontinuous exchanges between the
  classes, with a latent heat set by the near-degeneracy. Whether this
  accounts for the absence of universal scaling in the intermediate regime
  is a conjecture that Eq.~\eqref{eq:fkm_hist} can decide.
\item \emph{Entropy of the frustrated configurations.} For $U\gg t$ and the
  densities $n_f=n_c=\tfrac12$, each site carries one particle at low
  energy, and to leading order in $t/U$ the configurations $\mathbf f$
  interact as an antiferromagnetic Ising model on the triangular lattice,
  with a coupling of order $t^{2}/U$; higher orders lift its
  degeneracy~\cite{NejatiLenk2026}. The ground states of that Ising model
  number $e^{V\lds_0/k_B}$ with
  $\lds_0=0.323\,k_B$~\cite{Wannier1950}. They form the narrow group of
  levels of Section~\ref{sec:frustration}, of a width $W\ll t^{2}/U$ per
  site: $\lds_V(\varepsilon)$ rises to $\lds_0$ within
  $\delta\varepsilon\sim W$, and the heat capacity has a maximum at
  $k_BT\sim W$ in addition to a broad one at $k_BT\sim t^{2}/U$. Since the
  itinerant fermions fill the sites without a localised particle and have a
  gap of order $U$, $C_{\rm in}$ is activated, and the test is free of
  parameters:
  \[
    \int\frac{dT}{T}\,\frac{C_{\rm ex}}{V}=\lds_0=0.323\,k_B
  \]
  over the maximum at $k_BT\sim W$. The heat capacity reported for these
  densities at large $U$ has a broad feature and a sharp peak at a much
  lower temperature~\cite{NejatiLenk2026}; the integral has not been
  evaluated.
\item \emph{Fermi-gas form without quasiparticles.} Where
  $\bar\varrho_T(E)$ is finite and smooth on the scale $k_BT$ about $\mu_c$,
  Eq.~\eqref{eq:fkm_C} gives the Sommerfeld law for the first part,
  $C_{\rm in}/V=\tfrac{\pi^{2}}{3}k_B^{2}\,\bar\varrho_T(\mu_c)\,T$, with the
  density of states at the Fermi level that the simulation computes. It
  holds whether the single-particle states are extended or localised, and
  although the itinerant fermions scatter from the static configuration and
  are not quasiparticles of a Fermi liquid. A deviation of the heat capacity
  from the linear law in the intermediate regime must therefore lie in
  $C_{\rm ex}$ or in the dependence of $\bar\varrho_T(\mu_c)$ on the
  temperature; the level count does not register the localisation
  (Section~\ref{sec:disorder}).
\end{itemize}
Eq.~\eqref{eq:fkm_C} holds at fixed chemical potentials. If the densities
are held fixed instead, the chemical potentials vary with the temperature,
and the measured heat capacity differs from $C_\mu$ by the term given in the
footnote of Section~\ref{sec:filt_mb}.
%% ----------------------------- END OF APPENDIX D ----------------------------

\bibliographystyle{plain}
\bibliography{thermo_dos}

\end{document}
